% Lesson 11 — Words and expressions related to anticipating one's second language needs at work — Anglais Tle A — Cours signé OnBuch+
% Programme officiel MINESEC. Compiler avec : tectonic EN11-vocabulary-words-and-expressions-related-to-anticipatin.tex
\newcommand{\DOCMATIERE}{Anglais}
\newcommand{\DOCNIVEAU}{Tle A}
\newcommand{\DOCMODULE}{Chapter 1 : Family and Social Life}
\newcommand{\DOCLECON}{EN11}
\newcommand{\DOCTITRE}{Lesson 11 — Words and expressions related to anticipating one's second language needs at work}
% OnBuch+ — préambule « cours signés » ENRICHI (compilé avec tectonic / XeLaTeX)
% Système visuel « Pop » d'OnBuch+ : fond crème (jamais blanc), bordures encre
% épaisses, ombres « dures » décalées, titres Archivo Black en capitales.
% Version MATHÉMATIQUES : graphes (pgfplots/tikz), boîtes pédagogiques
% (définition, propriété, méthode, attention, le savais-tu, exercice résolu,
% exercice, TP, bilan), page de garde et signature OnBuch+ ancrée en bas.
%
% Macros attendues dans main.tex AVANT \input{preamble} :
%   \DOCMATIERE  \DOCNIVEAU  \DOCMODULE  \DOCLECON (numéro)  \DOCTITRE
\documentclass[11pt]{article}

\usepackage{fontspec}
\usepackage[english,french]{babel} % français = langue principale ; l'anglais via \eng{..} / textebox
\usepackage[a4paper,margin=2cm,top=3.4cm,bottom=2.8cm,headheight=1.4cm,footskip=1.3cm]{geometry}
\usepackage{amsmath,amssymb,amsfonts}
\usepackage{mathtools} % \xmapsto, \coloneqq... souvent utilisés par les modèles
\usepackage{cancel} % simplifications barrées (bilans de forces)
% \abs{z} (module, valeur absolue) et \norm{u} (norme), utilisés par les modèles
\providecommand{\abs}{}\let\abs\relax\DeclarePairedDelimiter{\abs}{\lvert}{\rvert}
\providecommand{\norm}{}\let\norm\relax\DeclarePairedDelimiter{\norm}{\lVert}{\rVert}
\usepackage[table]{xcolor} % [table] : fournit aussi \rowcolors (alternance de lignes)
\usepackage{pagecolor}
\usepackage{tikz}
\usetikzlibrary{arrows.meta,positioning,calc,shapes.geometric,shapes.misc,shapes.arrows,decorations.pathmorphing,decorations.pathreplacing,decorations.markings,patterns,shadows,mindmap,trees,fit,backgrounds,matrix,intersections,angles,quotes}
\usepackage{pgfplots}
\usepackage[version=4]{mhchem} % équations nucléaires \ce{^{235}_{92}U}
\usepackage[european,siunitx]{circuitikz} % schémas électriques
\pgfplotsset{compat=1.18}
\usepgfplotslibrary{fillbetween,groupplots} % aires entre courbes (calcul intégral)
\usepackage{enumitem}
\usepackage{fancyhdr}
% [most] charge toutes les bibliothèques tcolorbox (listings, minted, raster,
% documentation...) et gonfle inutilement la mémoire TeX sur un document long
% (~300 boîtes) : on ne charge que ce qui est réellement utilisé.
\usepackage[skins,breakable]{tcolorbox}
\usepackage{graphicx}
\usepackage{colortbl}
\usepackage{booktabs}
\usepackage{tabularx}
\usepackage{diagbox} % case d'en-tête diagonale des tableaux à double entrée
% Colonnes extensibles centrée / gauche / droite (C, L, R), souvent utilisées par les modèles
\newcolumntype{C}{>{\centering\arraybackslash}X}
\newcolumntype{L}{>{\raggedright\arraybackslash}X}
\newcolumntype{R}{>{\raggedleft\arraybackslash}X}
\usepackage{array}
\usepackage{multicol}
\usepackage{multirow}
\usepackage{siunitx}
% parse-numbers=false : accepte aussi \SI{2,5 \times 10^{-3}}{...}
\sisetup{locale=FR,output-decimal-marker={,},group-minimum-digits=4,parse-numbers=false}
\DeclareSIUnit\ohm{\ensuremath{\symup{\Omega}}} % U+2126 absent de la police
\DeclareSIUnit\division{div} % réglages d'oscilloscope (ms/div, V/div)
\DeclareSIUnit\tr{tr} % tours (vitesse de rotation, ex. \SI{1500}{\tr\per\minute})
\DeclareSIUnit\molar{M} % concentration molaire (ex. \SI{3}{\molar}, \SI{1}{\micro\molar})
\DeclareSIUnit\an{an} % année (le modèle confond parfois avec \year, un compteur TeX) — ex. \SI{5}{\kilo\meter\per\an}
\DeclareSIUnit\FCFA{FCFA} % franc CFA (ex. \SI{500}{\FCFA\per\kilo\gram})
\DeclareSIUnit\degreeBrix{\textdegree Brix} % teneur en sucre d'un sirop/jus (réfractométrie)
\DeclareSIUnit\month{mois} % le modèle utilise parfois \month, un compteur TeX, comme unité de durée
\usepackage{qrcode}
% \degree : commande fréquemment utilisée par les modèles pour un angle ou une
% température en dehors de \SI{}{} (non fournie par les paquets chargés).
\providecommand{\degree}{^{\circ}}
% \qed (fin de démonstration), utilisé par les modèles sans amsthm
\providecommand{\qed}{\hfill$\square$}
% \barre : double barre des valeurs interdites dans un tableau de variations (tkz-tab)
\providecommand{\barre}{\ensuremath{\|}}
% environnement proof (amsthm non chargé) utilisé par les modèles
\ifdefined\proof\else\newenvironment{proof}[1][Démonstration]{\par\noindent\textit{#1.}\ }{\qed\par}\fi
\usepackage{lastpage}
\usepackage{hyperref}
\hypersetup{hidelinks}

% ============================================================
% Palette « Pop » OnBuch+ — mêmes hex que la classe P (plus_ui.dart)
% ============================================================
\definecolor{poporange}{HTML}{F59321}
\definecolor{popdark}{HTML}{E07A0C}
\definecolor{popcream}{HTML}{FFF6E8}
\definecolor{popink}{HTML}{1C1714}
\definecolor{poppurple}{HTML}{7A5AE0}
\definecolor{popgreen}{HTML}{2E9E6A}
\definecolor{poppink}{HTML}{E0457B}
\definecolor{popblue}{HTML}{2F7DE1}
\definecolor{popgold}{HTML}{C98A12}
\definecolor{popmuted}{HTML}{8A7F76}
\definecolor{popline}{HTML}{EADFD2}
\definecolor{popgray}{HTML}{8A7F76}
\colorlet{popgrey}{popgray}
% Alias tolérants (le modèle nomme parfois les couleurs différemment)
\colorlet{popwhite}{white}
\colorlet{popblack}{popink}
\colorlet{popred}{poppink}
\colorlet{popyellow}{popgold}
% Teintes claires (fonds de tableaux/encadrés) — le modèle nomme parfois
% les couleurs avec un suffixe L (ex. popblueL) sans les avoir définies.
\colorlet{poporangeL}{poporange!20}
\colorlet{popdarkL}{popdark!20}
\colorlet{poppurpleL}{poppurple!20}
\colorlet{popgreenL}{popgreen!20}
\colorlet{poppinkL}{poppink!20}
\colorlet{popblueL}{popblue!20}
\colorlet{popgoldL}{popgold!20}
\colorlet{popredL}{popred!20}
\colorlet{popgrayL}{popgray!20}
% Teintes claires pour fonds de tableaux / zones
\colorlet{popblueL}{popblue!10!white}
\colorlet{poporangeL}{poporange!14!white}
\colorlet{popgreenL}{popgreen!12!white}
\colorlet{poppurpleL}{poppurple!12!white}
\colorlet{poppinkL}{poppink!12!white}
\colorlet{popgoldL}{popgold!14!white}

\pagecolor{popcream}

% ============================================================
% Typographie
% ============================================================
\defaultfontfeatures{Ligatures=TeX}
\setmainfont{PlusJakartaSans-Medium.ttf}[Path=../../../fonts/,BoldFont=PlusJakartaSans-Bold.ttf]
% Symboles, API et emojis absents de la police principale : repli sur DejaVu Sans (caractères actifs)
\newfontface\symfont{DejaVuSans.ttf}[Path=../../../fonts/]
\makeatletter
\newcommand\symactive[1]{\catcode#1=13 \begingroup\lccode`\~=#1 \lowercase{\endgroup\def~}{{\symfont\char#1}}}
\newcommand\symblank[1]{\catcode#1=13 \begingroup\lccode`\~=#1 \lowercase{\endgroup\def~}{}}
\newcommand\symhyph[1]{\catcode#1=13 \begingroup\lccode`\~=#1 \lowercase{\endgroup\def~}{\mbox{-}}}
\newcount\symc
\newcommand\symrange[2]{\symc=#1 \@whilenum\symc<#2 \do{\expandafter\symactive\expandafter{\the\symc}\advance\symc by 1 }}
\newcommand\symblankrange[2]{\symc=#1 \@whilenum\symc<#2 \do{\expandafter\symblank\expandafter{\the\symc}\advance\symc by 1 }}
\makeatother
\symrange{"0250}{"02FF}\symactive{"2713}\symactive{"2714}\symactive{"2717}\symactive{"2718}\symactive{"2610}\symactive{"2611}\symactive{"2612}
\symhyph{"2011}\symhyph{"2E1A}\symblankrange{"1F300}{"1FAFF}\symblank{"FE0F}
% Maths en sans-serif (Fira Math) avec chiffres et lettres droites de Plus
% Jakarta, pour que les formules se fondent dans le texte.
\usepackage[math-style=ISO,bold-style=ISO]{unicode-math}
\setmathfont{FiraMath-Regular.otf}[Path=../../../fonts/]
\setmathfont{PlusJakartaSans-Medium.ttf}[Path=../../../fonts/,range={up/num,up/{Latin,latin},\mathup}]
\usepackage{icomma} % virgule décimale sans espace en mode math
% Caractères mathématiques Unicode absents des polices de texte (Plus Jakarta,
% Archivo Black) : rendus par leur équivalent mathématique, en texte comme en
% formule — sinon ils s'affichent en carré vide (ex. « Arithmétique dans ℤ »).
\usepackage{newunicodechar}
\newunicodechar{ℝ}{\ensuremath{\mathbb{R}}}
\newunicodechar{ℤ}{\ensuremath{\mathbb{Z}}}
\newunicodechar{ℂ}{\ensuremath{\mathbb{C}}}
\newunicodechar{ℕ}{\ensuremath{\mathbb{N}}}
\newunicodechar{ℚ}{\ensuremath{\mathbb{Q}}}
\newunicodechar{𝔻}{\ensuremath{\mathbb{D}}}
\newunicodechar{α}{\ensuremath{\alpha}}
\newunicodechar{β}{\ensuremath{\beta}}
\newunicodechar{γ}{\ensuremath{\gamma}}
\newunicodechar{δ}{\ensuremath{\delta}}
\newunicodechar{ε}{\ensuremath{\varepsilon}}
\newunicodechar{θ}{\ensuremath{\theta}}
\newunicodechar{λ}{\ensuremath{\lambda}}
\newunicodechar{μ}{\ensuremath{\mu}}
\newunicodechar{σ}{\ensuremath{\sigma}}
\newunicodechar{τ}{\ensuremath{\tau}}
\newunicodechar{φ}{\ensuremath{\varphi}}
\newunicodechar{ψ}{\ensuremath{\psi}}
\newunicodechar{ω}{\ensuremath{\omega}}
\newunicodechar{Σ}{\ensuremath{\Sigma}}
% majuscules produites par \MakeUppercase dans les titres (α-aminés -> Α)
\newunicodechar{Α}{\ensuremath{\alpha}}
\newunicodechar{Β}{\ensuremath{\beta}}
\newunicodechar{Γ}{\ensuremath{\Gamma}}
\newunicodechar{Δ}{\ensuremath{\Delta}}
\newunicodechar{Ε}{\ensuremath{\varepsilon}}
\newunicodechar{Θ}{\ensuremath{\Theta}}
\newunicodechar{Λ}{\ensuremath{\Lambda}}
\newunicodechar{Μ}{\ensuremath{\mu}}
\newunicodechar{Τ}{\ensuremath{\tau}}
\newunicodechar{Φ}{\ensuremath{\Phi}}
\newunicodechar{Ψ}{\ensuremath{\Psi}}
\newunicodechar{Ω}{\ensuremath{\Omega}}
\newunicodechar{↦}{\ensuremath{\mapsto}}
\newunicodechar{∈}{\ensuremath{\in}}
\newunicodechar{∉}{\ensuremath{\notin}}
\newunicodechar{∧}{\ensuremath{\wedge}}
\newunicodechar{⇔}{\ensuremath{\Leftrightarrow}}
\newunicodechar{⟺}{\ensuremath{\Longleftrightarrow}}
\newunicodechar{⇒}{\ensuremath{\Rightarrow}}
\newunicodechar{⟹}{\ensuremath{\Longrightarrow}}
\newunicodechar{∘}{\ensuremath{\circ}}
\newunicodechar{∪}{\ensuremath{\cup}}
\newunicodechar{∩}{\ensuremath{\cap}}
\newunicodechar{≡}{\ensuremath{\equiv}}
\newunicodechar{⊂}{\ensuremath{\subset}}
\newunicodechar{⊥}{\ensuremath{\perp}}
\newunicodechar{∃}{\ensuremath{\exists}}
\newunicodechar{∀}{\ensuremath{\forall}}
\newunicodechar{∛}{\ensuremath{\sqrt[3]{\,}}}
\newunicodechar{∜}{\ensuremath{\sqrt[4]{\,}}}
\newunicodechar{⃗}{\ensuremath{{}^{\to}}}
\newunicodechar{ⁿ}{\ifmmode^{n}\else\textsuperscript{\ensuremath{n}}\fi}
\newunicodechar{ˣ}{\ifmmode^{x}\else\textsuperscript{\ensuremath{x}}\fi}
\newunicodechar{ᵇ}{\ifmmode^{b}\else\textsuperscript{\ensuremath{b}}\fi}
\newunicodechar{ʳ}{\ifmmode^{r}\else\textsuperscript{\ensuremath{r}}\fi}
\newunicodechar{ᵘ}{\ifmmode^{u}\else\textsuperscript{\ensuremath{u}}\fi}
\newunicodechar{ʸ}{\ifmmode^{y}\else\textsuperscript{\ensuremath{y}}\fi}
\newunicodechar{ᵃ}{\ifmmode^{a}\else\textsuperscript{\ensuremath{a}}\fi}
\newunicodechar{ᵉ}{\ifmmode^{e}\else\textsuperscript{\ensuremath{e}}\fi}
\newunicodechar{ᵏ}{\ifmmode^{k}\else\textsuperscript{\ensuremath{k}}\fi}
\newunicodechar{ᵖ}{\ifmmode^{p}\else\textsuperscript{\ensuremath{p}}\fi}
\newunicodechar{ᴺ}{\ifmmode^{N}\else\textsuperscript{\ensuremath{N}}\fi}
\newunicodechar{ᶜ}{\ifmmode^{c}\else\textsuperscript{\ensuremath{c}}\fi}
\newunicodechar{ʰ}{\ifmmode^{h}\else\textsuperscript{\ensuremath{h}}\fi}
\newunicodechar{ᵠ}{\ifmmode^{q}\else\textsuperscript{\ensuremath{q}}\fi}
\newunicodechar{ₙ}{\ifmmode_{n}\else\textsubscript{\ensuremath{n}}\fi}
\newunicodechar{ₐ}{\ifmmode_{a}\else\textsubscript{\ensuremath{a}}\fi}
\newunicodechar{ᵢ}{\ifmmode_{i}\else\textsubscript{\ensuremath{i}}\fi}
\newunicodechar{ᵣ}{\ifmmode_{r}\else\textsubscript{\ensuremath{r}}\fi}
\newunicodechar{ₖ}{\ifmmode_{k}\else\textsubscript{\ensuremath{k}}\fi}
\newunicodechar{ᵦ}{\ifmmode_{\beta}\else\textsubscript{\ensuremath{\beta}}\fi}
\newunicodechar{ₓ}{\ifmmode_{x}\else\textsubscript{\ensuremath{x}}\fi}
\newfontfamily\popdisplay{ArchivoBlack-Regular.ttf}[Path=../../../fonts/]
\newfontfamily\poplogo{SpaceGrotesk-Bold.ttf}[Path=../../../fonts/]
\newfontfamily\popsemi{PlusJakartaSans-SemiBold.ttf}[Path=../../../fonts/]
\newfontfamily\popxbold{PlusJakartaSans-ExtraBold.ttf}[Path=../../../fonts/]
\newfontfamily\popxboldls{PlusJakartaSans-ExtraBold.ttf}[Path=../../../fonts/,LetterSpace=3.0]

\setlist[enumerate]{leftmargin=1.7em,itemsep=5pt,topsep=4pt}
\setlist[itemize]{leftmargin=1.7em,itemsep=4pt,topsep=4pt,label={\color{poporange}\textbullet}}
\setlength{\parindent}{0pt}
\setlength{\parskip}{5pt}
\renewcommand{\arraystretch}{1.25}

% pgfplots : style Pop par défaut
\pgfplotsset{
  every axis/.append style={axis line style={popink,line width=1pt},tick style={popink},
    grid style={popline,line width=0.5pt},label style={font=\small},tick label style={font=\footnotesize}},
  popcurve/.style={poporange,line width=1.8pt,no markers,smooth},
}
% Même style disponible dans un \draw TikZ simple (hors pgfplots)
\tikzset{popcurve/.style={poporange,line width=1.8pt}}
\tikzset{popgrid/.style={popline,very thin}}
\colorlet{popcurve}{poporange} % le nom du style est parfois utilisé comme couleur

% ============================================================
% Pill / squiggle
% ============================================================
\newcommand{\poppill}[3]{%
  \tikz[baseline=-0.55ex]\node[draw=popink,line width=1.2pt,rounded corners=2.6pt,%
    fill=#2,inner xsep=6.5pt,inner ysep=3pt]{%
    \popxboldls\fontsize{7}{7}\selectfont\color{#3}\MakeUppercase{#1}};%
}
\newcommand{\popsquiggle}[1]{%
  \tikz[baseline=-0.4ex]{
    \draw[#1,line width=2.6pt,line cap=round]
      plot [smooth] coordinates {(0,0.09) (0.34,0.01) (0.68,0.21) (1.02,0.01) (1.36,0.10)};
  }%
}
% Mot-clé surligné (marqueur orange sous le texte)
\newcommand{\cle}[1]{{\popxbold\color{popdark}#1}}

% ============================================================
% En-tête / pied
% ============================================================
\pagestyle{fancy}
\fancyhf{}
\renewcommand{\headrulewidth}{0pt}
\renewcommand{\footrulewidth}{0pt}
\lhead{\raisebox{-0.15ex}{\poplogo\fontsize{13}{13}\selectfont\color{popink}OnBuch\color{poporange}+}}
\rhead{\poppill{\DOCMATIERE\ \textbullet\ \DOCNIVEAU\ \textbullet\ Leçon \DOCLECON}{white}{popink}}
\lfoot{\popxbold\fontsize{7.6}{7.6}\selectfont\color{popmuted}\MakeUppercase{Cours signé OnBuch+}\ \textbullet\ {\popsemi\DOCTITRE}}
\rfoot{\tikz[baseline=-0.6ex]\node[rounded corners=8.5pt,fill=popink,minimum height=17pt,minimum width=17pt,inner xsep=5pt,inner sep=0]{\popxbold\fontsize{8}{8}\selectfont\color{popcream}\thepage\,/\,\pageref*{LastPage}};}

% ============================================================
% Titres
% ============================================================
\newcommand{\chaptitre}[2]{%
  \begin{tcolorbox}[enhanced,colback=poporange,colframe=popink,boxrule=2.6pt,arc=11pt,
      drop shadow=popink,left=15pt,right=15pt,top=12pt,bottom=11pt,before skip=3pt,after skip=18pt]
    \poppill{#2}{popink}{popcream}\par\vspace{9pt}
    {\raggedright\hyphenpenalty=10000\exhyphenpenalty=10000\popdisplay\fontsize{22}{24}\selectfont\color{popcream}\MakeUppercase{#1}\par}
  \end{tcolorbox}
}
% Section numérotée : pastille orange avec le numéro + titre Archivo
\newcounter{popsec}
\newcommand{\coursec}[1]{%
  \stepcounter{popsec}\setcounter{popsub}{0}%
  \par\vspace{16pt}\noindent
  \begin{minipage}{\linewidth}
  \tikz[baseline=-0.7ex]\node[draw=popink,line width=1.6pt,fill=poporange,rounded corners=4pt,minimum size=22pt,inner sep=0]{\popdisplay\fontsize{12}{12}\selectfont\color{popcream}\thepopsec};%
  \hspace{8pt}{\popdisplay\fontsize{15}{17}\selectfont\color{popink}\MakeUppercase{#1}}\par
  \vspace{4pt}\noindent{\color{poporange}\rule{36pt}{3.6pt}}
  \end{minipage}\par\nopagebreak\vspace{8pt}
}
\newcounter{popsub}
\newcommand{\courssub}[1]{%
  \stepcounter{popsub}%
  \par\vspace{9pt}\noindent{\popxbold\fontsize{11.5}{13}\selectfont\color{popink}{\color{poporange}\thepopsec.\thepopsub}\hspace{6pt}#1}\par\nopagebreak\vspace{4pt}
}

% ============================================================
% Boîtes pédagogiques — même recette : fond blanc, bordure épaisse colorée,
% ombre dure encre, badge plein. Titre optionnel [#1].
% ============================================================
\tcbset{popbase/.style={breakable,enhanced,colback=white,boxrule=2.3pt,arc=9pt,
  shadow={3pt}{-3pt}{0pt}{popink},left=14pt,right=14pt,top=11pt,bottom=11pt,before skip=11pt,after skip=15pt,
  fontupper=\popsemi\fontsize{10.6}{14.5}\selectfont\color{popink},
  fontlower=\popsemi\fontsize{10.6}{14.5}\selectfont\color{popink}}}
% Coche dessinée (le glyphe ✓ n'existe pas dans les polices du document)
\newcommand{\popcheck}[1]{\tikz[baseline=-0.5ex]\draw[#1,line width=1.6pt,line cap=round,line join=round](-0.13,0.0)--(-0.04,-0.09)--(0.13,0.1);}
\providecommand{\coche}{\popcheck{popgreen}}
\providecommand{\resultat}[1]{\par\smallskip\noindent\fcolorbox{popgreen}{popgreenL}{\parbox{\dimexpr\linewidth-2\fboxsep-2\fboxrule\relax}{\textbf{Résultat.} #1}}\par\smallskip}
\newcommand{\popboxhead}[3]{\poppill{#1}{#2}{white}\ifx\relax#3\relax\else\hspace{6pt}{\popxbold\fontsize{10.6}{12}\selectfont\color{popink}#3}\fi\par\vspace{7pt}}

\newtcolorbox{aretenir}[1][]{popbase,colframe=popgreen,before upper={\popboxhead{À retenir}{popgreen}{#1}}}
% Texte en anglais (document de lecture, dialogue, extrait) : cadre
% d'étude. Un titre optionnel (source, auteur) s'affiche dans l'en-tête.
\newtcolorbox{textebox}[1][]{popbase,colframe=popblue,before upper={\popboxhead{Text}{popblue}{#1}\begin{otherlanguage}{english}},after upper={\end{otherlanguage}}}
% Marques ✓ / ✗ (forme correcte / fautive) souvent utilisées par les modèles
\providecommand{\cross}{\textcolor{poppink}{\ensuremath{\times}}}
\providecommand{\xmark}{\cross}\providecommand{\crossmark}{\cross}\providecommand{\cmark}{\ensuremath{\checkmark}}
% Ligne de réponse / justification à compléter (inventée par les modèles)
\providecommand{\justification}[1]{\par\noindent\textit{Justification~:} \dotfill\par}
% \textipa (package tipa non chargé) : la police ne couvre pas l'API -> simple passage
\providecommand{\textipa}[1]{#1}
% Soulignement (ulem non chargé) : \uline, \ul
\providecommand{\uline}[1]{\underline{#1}}\providecommand{\ul}[1]{\underline{#1}}
% \glossaire{\item ...} inventée par les modèles : liste de glossaire
\providecommand{\glossaire}[1]{\begin{itemize}#1\end{itemize}}
% \alert (beamer) inventée par les modèles : texte mis en évidence
\providecommand{\alert}[1]{\textbf{\textcolor{poppink}{#1}}}
% variantes ASCII d'ordinaux écrites en macro
\providecommand{\iere}{\textsuperscript{re}}\providecommand{\ieme}{\textsuperscript{e}}\providecommand{\ere}{\textsuperscript{re}}\providecommand{\eme}{\textsuperscript{e}}\providecommand{\er}{\textsuperscript{er}}\providecommand{\ie}{\textsuperscript{e}}
% Ordinaux français écrits en macro par les modèles : 1\ière, 3\ième
\newcommand{\ière}{\textsuperscript{re}}\newcommand{\ième}{\textsuperscript{e}}
% \exemple (inventée par les modèles) : étiquette « Exemple : »
\providecommand{\exemple}{\textbf{Exemple~:}\ }
% \square utilisable aussi hors mode math (cases à cocher)
\AtBeginDocument{\let\squareMath\square\renewcommand{\square}{\ensuremath{\squareMath}}}
% Mot ou phrase en anglais à l'intérieur d'un paragraphe français
\newcommand{\eng}[1]{\ifmmode\text{\textit{#1}}\else\begingroup\selectlanguage{english}\itshape #1\endgroup\fi}
\let\esp\eng % alias toléré
% flèches d'intonation inventées par les modèles
\providecommand{\textsearrow}{}\renewcommand{\textsearrow}{\ensuremath{\searrow}}\providecommand{\textnearrow}{}\renewcommand{\textnearrow}{\ensuremath{\nearrow}}
% \fr{..} : mot français (inventé par les modèles)
\providecommand{\fr}[1]{\textit{#1}}
\newtcolorbox{exemplebox}[1][]{popbase,colframe=poporange,before upper={\popboxhead{Exemple}{poporange}{#1}}}
\newtcolorbox{definition}[1][]{popbase,colframe=popblue,before upper={\popboxhead{Définition}{popblue}{#1}}}
\newtcolorbox{propriete}[1][]{popbase,colframe=poppurple,before upper={\popboxhead{Propriété}{poppurple}{#1}}}
\newtcolorbox{methode}[1][]{popbase,colframe=popgold,before upper={\popboxhead{Méthode}{popgold}{#1}}}
\newtcolorbox{attention}[1][]{popbase,colframe=poppink,before upper={\popboxhead{Attention}{poppink}{#1}}}
\newtcolorbox{experience}[1][]{popbase,colframe=popblue,colback=popblueL,before upper={\popboxhead{Expérience}{popblue}{#1}}}
% « Le savais-tu ? » : carte violette pleine (contexte, culture, Cameroun)
\newtcolorbox{savaistu}[1][]{popbase,colback=poppurple,colframe=popink,
  fontupper=\popsemi\fontsize{10.4}{14.2}\selectfont\color{white},
  before upper={\poppill{Le savais-tu\,?}{white}{poppurple}\ifx\relax#1\relax\else\hspace{6pt}{\popxbold\color{white}#1}\fi\par\vspace{7pt}}}
% Exercice résolu : énoncé en haut, \tcblower puis la solution sur fond crème
\newcounter{popexo}\newcounter{popexr}
\newtcolorbox{exoresolu}[1][]{popbase,colframe=poporange,colbacklower=popcream,
  segmentation style={popink,line width=1.2pt,dashed},
  before upper={\stepcounter{popexr}\popboxhead{Exercice résolu \thepopexr}{poporange}{#1}},
  before lower={\poppill{Solution}{popink}{popcream}\par\vspace{6pt}}}
% Exercice (non résolu en place) — la correction est dans la partie Corrigés
\newtcolorbox{exercice}[1][]{popbase,colframe=popink,
  before upper={\stepcounter{popexo}\popboxhead{Exercice \thepopexo}{popink}{#1}}}
% Corrigé
\newtcolorbox{corrige}[1][]{popbase,colframe=popgreen,colback=popgreenL,
  before upper={\popboxhead{Corrigé}{popgreen}{#1}}}
% Objectifs / prérequis / bilan
\newtcolorbox{objectifs}{popbase,colframe=popink,colback=poporangeL,
  before upper={\popboxhead{Objectifs de la leçon}{popink}{}}}
\newtcolorbox{prerequis}{popbase,colframe=popink,colback=popblueL,
  before upper={\popboxhead{Prérequis}{popblue}{}}}
\newtcolorbox{bilan}{popbase,colframe=popink,colback=white,boxrule=2.8pt,
  before upper={\popboxhead{Fiche bilan}{poporange}{L'essentiel à connaître pour l'examen}}}
% Figure encadrée avec légende
\newtcolorbox{popfigure}[1][]{enhanced,breakable=false,colback=white,colframe=popline,boxrule=1.4pt,arc=8pt,
  left=10pt,right=10pt,top=10pt,bottom=8pt,before skip=10pt,after skip=12pt,halign=center,
  lower separated=false,halign lower=center,
  fontlower=\popsemi\fontsize{9}{11.5}\selectfont\color{popmuted},
  #1}
\newcommand{\legende}[1]{\par\vspace{6pt}{\popsemi\fontsize{9}{11.5}\selectfont\color{popmuted}#1}}

% ============================================================
% Signature OnBuch+
% ============================================================
\newcommand{\poplogoinline}[1]{{\poplogo\fontsize{#1}{#1}\selectfont\color{popink}OnBuch\color{poporange}+}}
\newcommand{\signatureonbuch}{%
  \begin{tcolorbox}[enhanced,colback=popink,colframe=popink,boxrule=2.2pt,arc=10pt,
      drop shadow=poporange,left=14pt,right=14pt,top=10pt,bottom=10pt,before skip=0pt,after skip=0pt]
    \tikz[baseline=-0.7ex]\node[circle,fill=popgreen,draw=popcream,line width=1.3pt,minimum size=22pt,inner sep=0]{\popcheck{white}};%
    \hspace{9pt}\begin{minipage}[c]{0.62\linewidth}
      {\popxbold\fontsize{12}{13}\selectfont\color{popcream}Signé\ \textbullet\ {\poplogo OnBuch\color{poporange}+}}\par\vspace{2pt}
      {\raggedright\popsemi\fontsize{8}{10.5}\selectfont\color{popline}\MakeUppercase{Équipe pédagogique OnBuch+}\\\MakeUppercase{Conforme au programme officiel MINESEC}\par}
    \end{minipage}\hfill
    {\poplogo\fontsize{22}{22}\selectfont\color{popcream}OnBuch\color{poporange}+}
  \end{tcolorbox}%
}

% ============================================================
% Page de garde
% #1 titre · #2 matière · #3 niveau · #4 module · #5 n° leçon · #6 durée ·
% #7 description / objectifs (texte ou itemize)
% ============================================================
\newcommand{\pagegarde}[7]{%
  \thispagestyle{empty}
  \newgeometry{a4paper,margin=1.7cm,top=1.6cm,bottom=1.4cm}%
  \begin{tikzpicture}[remember picture,overlay]
    \fill[poporange!18] ([xshift=-3.2cm,yshift=-3.6cm]current page.north east) circle (4.6cm);
    \fill[poppurple!14] ([xshift=2.4cm,yshift=7.6cm]current page.south west) circle (3.8cm);
  \end{tikzpicture}%
  {\poplogo\fontsize{30}{30}\selectfont\color{popink}OnBuch\color{poporange}+}\hfill
  \raisebox{0.6ex}{\poppill{Cours signé \textbullet\ Premium}{popink}{popcream}}\par
  \vspace{1pt}\popsquiggle{poppurple}\par
  \vspace{8pt}
  \begin{tcolorbox}[enhanced,colback=poporange,colframe=popink,boxrule=2.8pt,arc=13pt,
      drop shadow=popink,left=18pt,right=18pt,top=12pt,bottom=13pt,after skip=10pt]
    \poppill{#2 \textbullet\ #3}{popink}{popcream}\hspace{5pt}\poppill{#4}{white}{popink}\par\vspace{8pt}
    {\popdisplay\fontsize{14}{14}\selectfont\color{popink}LEÇON #5}\par\vspace{4pt}
    {\raggedright\hyphenpenalty=10000\exhyphenpenalty=10000\popdisplay\fontsize{25}{27}\selectfont\color{popcream}\MakeUppercase{#1}\par}
    \vspace{8pt}
    {\popxbold\fontsize{9.5}{9.5}\selectfont\color{popink}\MakeUppercase{Durée conseillée : #6}}
  \end{tcolorbox}
  \begin{tcolorbox}[enhanced,colback=white,colframe=popink,boxrule=2.2pt,arc=10pt,
      drop shadow=popink,left=16pt,right=16pt,top=10pt,bottom=10pt,after skip=8pt]
    \poppill{Dans cette leçon}{poppurple}{white}\par\vspace{5pt}
    \popsemi\fontsize{9.3}{12.6}\selectfont\color{popink}#7
  \end{tcolorbox}
  \noindent\begin{minipage}[t]{0.487\linewidth}
    \begin{tcolorbox}[enhanced,colback=popgreen,colframe=popink,boxrule=2.2pt,arc=10pt,
        drop shadow=popink,left=11pt,right=11pt,top=10pt,bottom=10pt,height=3.1cm,valign=center]
      \tcbox[colback=white,colframe=popink,boxrule=1.3pt,arc=5pt,boxsep=0pt,nobeforeafter,
        left=3pt,right=3pt,top=3pt,bottom=3pt,valign=center]{\qrcode[height=1.9cm]{https://play.google.com/store/apps/details?id=com.luvvix.onbuch&hl=fr}}%
      \hspace{8pt}\begin{minipage}[c]{0.5\linewidth}
        {\popdisplay\fontsize{11}{12.5}\selectfont\color{white}\MakeUppercase{Télécharge OnBuch}}\par\vspace{4pt}
        {\raggedright\popsemi\fontsize{8.4}{11}\selectfont\color{white}Gratuit \textbullet\ Google Play\\Tuteur IA, annales, concours, résultats\par}
      \end{minipage}
    \end{tcolorbox}
  \end{minipage}\hfill
  \begin{minipage}[t]{0.487\linewidth}
    \begin{tcolorbox}[enhanced,colback=poppurple,colframe=popink,boxrule=2.2pt,arc=10pt,
        drop shadow=popink,left=11pt,right=11pt,top=10pt,bottom=10pt,height=3.1cm,valign=center]
      \tikz[baseline=-0.7ex]\node[circle,fill=white,draw=popink,line width=1.3pt,minimum size=1.6cm,inner sep=0]{\popdisplay\fontsize{24}{24}\selectfont\color{poppurple}+};%
      \hspace{8pt}\begin{minipage}[c]{0.6\linewidth}
        {\popdisplay\fontsize{11}{12.5}\selectfont\color{white}\MakeUppercase{Passe à OnBuch+}}\par\vspace{4pt}
        {\raggedright\popsemi\fontsize{8.4}{11}\selectfont\color{white}Tous les cours signés, Tuteur IA illimité, fiches PDF — onglet OnBuch+ de l'app\par}
      \end{minipage}
    \end{tcolorbox}
  \end{minipage}
  \vspace{8pt}
  \popcontenu
  \vfill
  \signatureonbuch
  \restoregeometry
  \newpage
}

% Rangée « Ce cours contient » (page de garde)
\newcommand{\poptile}[3]{% #1 couleur #2 gros texte #3 légende
  \begin{tcolorbox}[enhanced,colback=#1,colframe=popink,boxrule=2pt,arc=9pt,shadow={2.5pt}{-2.5pt}{0pt}{popink},
      left=6pt,right=6pt,top=9pt,bottom=9pt,halign=center,valign=center,height=1.75cm,width=\linewidth,nobeforeafter]
    {\popdisplay\fontsize{12.5}{13}\selectfont\color{white}\MakeUppercase{#2}}\par\vspace{3pt}
    {\centering\popsemi\fontsize{7.8}{9.5}\selectfont\color{white}#3\par}
  \end{tcolorbox}}
\newcommand{\popcontenu}{%
  \par\noindent{\popxboldls\fontsize{7.5}{7.5}\selectfont\color{popmuted}CE COURS SIGNÉ CONTIENT}\par\vspace{7pt}
  \noindent\begin{minipage}[t]{0.235\linewidth}\poptile{popblue}{Cours}{complet, illustré, pas à pas}\end{minipage}\hfill
  \begin{minipage}[t]{0.235\linewidth}\poptile{poporange}{Méthodes}{et exercices résolus}\end{minipage}\hfill
  \begin{minipage}[t]{0.235\linewidth}\poptile{poppink}{Exercices}{type examen + corrigés}\end{minipage}\hfill
  \begin{minipage}[t]{0.235\linewidth}\poptile{popgreen}{Fiche bilan}{l'essentiel à retenir}\end{minipage}\par}

% Sommaire
\newtcolorbox{sommaire}{enhanced,colback=white,colframe=popink,boxrule=2.2pt,arc=10pt,
  drop shadow=popink,left=18pt,right=18pt,top=13pt,bottom=13pt,before skip=6pt,after skip=20pt,
  before upper={\poppill{Sommaire}{poppurple}{white}\par\vspace{9pt}\popsemi\fontsize{10.6}{14.5}\selectfont\color{popink}}}

% Signature de fin de cours — poussée tout en bas de la dernière page
\newcommand{\findecours}{%
  \par\vfill\nopagebreak
  \begin{center}\popsquiggle{poporange}\end{center}\vspace{4pt}
  \signatureonbuch
}
\AtBeginDocument{\def\llbracket{\mathopen{[\mkern-3.5mu[}}\def\rrbracket{\mathclose{]\mkern-3.5mu]}}} % intervalles d'entiers (glyphe ⟦ absent de la police)
% Glyphes absents de Fira Math : redéfinis à partir de glyphes disponibles
\AtBeginDocument{%
  \def\setminus{\mathbin{\backslash}}\let\smallsetminus\setminus
  \def\vdots{\mathord{\vbox{\baselineskip3.2pt\lineskiplimit0pt\kern4pt\hbox{.}\hbox{.}\hbox{.}}}}%
  \def\ddots{\mathinner{\mkern1mu\raise6.5pt\hbox{.}\mkern2mu\raise3.6pt\hbox{.}\mkern2mu\raise0.7pt\hbox{.}\mkern1mu}}%
  \def\triangle{\mathord{\scalebox{0.9}{$\symup{\Delta}$}}}%
  \def\nabla{\mathord{\raisebox{\depth}{\scalebox{1}[-1]{$\symup{\Delta}$}}}}%
  \def\checkmark{\popcheck{popdark}}%
  \def\star{\mathbin{\ast}}\def\bigstar{\mathord{\ast}}%
  \def\diamond{\mathbin{\scalebox{0.6}{$\Diamond$}}}%
  \def\bigcup{\mathop{\mathchoice{\raisebox{-0.3em}{\scalebox{1.7}{$\cup$}}}{\scalebox{1.25}{$\cup$}}{\cup}{\cup}}\displaylimits}%
  \def\bigcap{\mathop{\mathchoice{\raisebox{-0.3em}{\scalebox{1.7}{$\cap$}}}{\scalebox{1.25}{$\cap$}}{\cap}{\cap}}\displaylimits}%
  \def\lesseqqgtr{\mathrel{\lessgtr}}\def\bigoplus{\mathop{\oplus}\displaylimits}%
}
\providecommand{\ssi}{si et seulement si} % abréviation fréquente
\AtBeginDocument{\providecommand{\wideparen}{\overparen}} % arc de cercle

%% FIGFIX-BEGIN v5 — correctifs globaux des figures (docs/onbuch-generation/tools/figfix)
\usepackage{adjustbox}\usepackage{etoolbox}\tcbuselibrary{raster}
% toute figure TikZ/pgfplots plus large que la ligne est réduite à la largeur disponible
\BeforeBeginEnvironment{tikzpicture}{\begin{adjustbox}{max width=\linewidth}}
\AfterEndEnvironment{tikzpicture}{\end{adjustbox}}
% légendes pgfplots lisibles par-dessus les courbes
\pgfplotsset{every axis/.append style={legend style={fill=white,fill opacity=0.92,text opacity=1,draw=black!15,inner sep=3pt}}}
% étiquettes d'arêtes (syntaxe edge["..."]) avec fond blanc
\tikzset{every edge quotes/.append style={fill=white,inner sep=1.2pt}}
%% FIGFIX-END
\begin{document}

\pagegarde{Lesson 11 — Words and expressions related to anticipating one's second language needs at work}{Anglais}{Tle A}{Chapter 1 : Family and Social Life}{EN11}{2 h}{Cette leçon de vocabulaire présente les mots, collocations et expressions liés à l'anticipation de ses besoins en langue seconde au travail : évaluer son niveau, demander une formation, préparer un entretien ou une mobilité professionnelle. Elle propose des champs lexicaux, des faux amis, des familles de mots et des mises en contexte pour un réemploi oral et écrit.
\vspace{4pt}
\begin{multicols}{2}\small
\begin{itemize}
\item À la fin de cette leçon, je sais identifier et définir le vocabulaire lié aux besoins linguistiques au travail (language needs, proficiency, training, fluency...).
\item À la fin de cette leçon, je sais employer les collocations professionnelles correctes (attend a training course, brush up one's English, meet the requirements...).
\item À la fin de cette leçon, je sais évaluer et décrire mon niveau de langue avec les termes appropriés (beginner, intermediate, fluent, proficient).
\item À la fin de cette leçon, je sais éviter les faux amis et les calques du français (formation ≠ training, assister ≠ to assist...).
\item À la fin de cette leçon, je sais former des mots de la même famille (employ, employer, employee, employment, unemployment).
\item À la fin de cette leçon, je sais réemployer ce lexique dans des phrases, des dialogues et de courts textes professionnels.
\end{itemize}\end{multicols}}

\begin{sommaire}
\begin{multicols}{2}
\begin{enumerate}
\item Évaluer son niveau et ses besoins linguistiques
\item Formation et développement des compétences
\item L'anglais dans la vie professionnelle
\item Faux amis et pièges des francophones
\item Familles de mots, prononciation et orthographe
\item Réemploi en contexte et production guidée
\item Activité d'intégration
\item Méthodes et exercices résolus
\item Exercices
\item Corrigés des exercices
\item Fiche bilan
\end{enumerate}
\end{multicols}
\end{sommaire}

\begin{objectifs}
\begin{itemize}
\item À la fin de cette leçon, je sais identifier et définir le vocabulaire lié aux besoins linguistiques au travail (language needs, proficiency, training, fluency...).
\item À la fin de cette leçon, je sais employer les collocations professionnelles correctes (attend a training course, brush up one's English, meet the requirements...).
\item À la fin de cette leçon, je sais évaluer et décrire mon niveau de langue avec les termes appropriés (beginner, intermediate, fluent, proficient).
\item À la fin de cette leçon, je sais éviter les faux amis et les calques du français (formation ≠ training, assister ≠ to assist...).
\item À la fin de cette leçon, je sais former des mots de la même famille (employ, employer, employee, employment, unemployment).
\item À la fin de cette leçon, je sais réemployer ce lexique dans des phrases, des dialogues et de courts textes professionnels.
\item À la fin de cette leçon, je sais rédiger un court paragraphe expliquant comment j'anticipe mes besoins en anglais pour ma future carrière.
\end{itemize}
\end{objectifs}

\begin{prerequis}
\begin{itemize}
\item Vocabulaire général du travail et des professions vu en Première (job, career, skill, qualification)
\item Les verbes modaux de nécessité et de conseil (must, need to, should, have to)
\item La construction des comparatifs et superlatifs pour évaluer un niveau (better, the most fluent)
\item Les structures de base de la lettre formelle et du dialogue d'entretien
\item La notion de famille de mots (préfixes et suffixes : -er, -ee, -ment, -tion)
\end{itemize}
\end{prerequis}

\newpage

\coursec{Évaluer son niveau et ses besoins linguistiques}

\textbf{Situation d'accroche.} Imagine : tu viens d'être recruté par une entreprise de Douala qui travaille avec des partenaires du Nigeria et du Ghana. Le premier jour, ton chef te demande : \eng{“Can you draft a report and attend the briefing at 10?”} Tu réalises soudain que l'anglais scolaire ne suffit pas : il faut \cle{anticiper} tes besoins en langue seconde au travail — savoir demander une formation, évaluer ton niveau, préparer un entretien. Cette leçon te donne les mots et expressions pour \emph{identifier}, \emph{exprimer} et \emph{planifier} tes besoins en anglais professionnel, une compétence clé pour le Bac et pour ta vie active au Cameroun bilingue.

\textbf{Problématique.} Comment dire, en anglais correct et naturel, quel est mon niveau, quelles sont mes forces et mes faiblesses, et ce que je dois faire pour progresser avant un entretien ou un nouveau poste ?

\courssub{Observer : un texte pour découvrir le vocabulaire}

\begin{textebox}[Sandra assesses her English — reading text]
Before applying for the job in Bamenda, Sandra assessed her English level with a placement test at a language centre. The results showed that she had an intermediate level: her listening and speaking were fairly good, but her writing skills were weak. “I can get by in English when I speak,” she told the trainer, “but my written English is a bit rusty, and I make too many grammar mistakes.” The trainer explained that fluency is not the same as accuracy: Sandra spoke easily, but not always correctly. Together, they identified her language needs for the new job — writing reports, answering emails and attending meetings in English. Sandra then set clear goals: she decided to upgrade her writing skills with an evening course and to practise every day. Three months later, she took the test again and reached an upper-intermediate level. “Now I have a good command of English,” she said with a smile, “and I feel ready for the interview.”
\end{textebox}

\textbf{Glossaire.}

\begin{center}
\begin{tabularx}{\linewidth}{|l|l|X|}
\hline
\rowcolor{popblueL} \textbf{Mot anglais} & \textbf{Nature} & \textbf{Traduction} \\
\hline
to assess & verbe & évaluer \\
\hline
placement test & nom & test de niveau (de positionnement) \\
\hline
weak & adjectif & faible \\
\hline
to get by & verbe & se débrouiller \\
\hline
rusty & adjectif & rouillé (par manque de pratique) \\
\hline
fluency & nom & aisance (orale) \\
\hline
accuracy & nom & exactitude, correction \\
\hline
to set goals & expression & se fixer des objectifs \\
\hline
to upgrade & verbe & améliorer, mettre à niveau \\
\hline
command (of English) & nom & maîtrise (de l'anglais) \\
\hline
\end{tabularx}
\end{center}

\textbf{Questions de compréhension.}
\begin{enumerate}
\item What was Sandra's level at the beginning of the text?
\item Which skill was weak, and which skills were fairly good?
\item What is the difference between fluency and accuracy, according to the trainer?
\item What did Sandra do to improve her English?
\end{enumerate}

\courssub{Le champ lexical de l'évaluation du niveau}

\begin{definition}[Proficiency]
\eng{Proficiency} (nom, se prononce « pro-fi-chen-si ») : \emph{the ability to use a language well and effectively} — la maîtrise, la compétence globale dans une langue. On dit \eng{English proficiency} (la maîtrise de l'anglais) et \eng{a proficiency test} (un test de compétence linguistique).
\end{definition}

Observons les mots du texte : \eng{level}, \eng{intermediate}, \eng{upper-intermediate}, \eng{fluent}, \eng{command of English}. En anglais, le niveau d'une langue se décrit avec une \emph{échelle d'adjectifs} bien établie, qu'il faut connaître dans l'ordre.

\begin{center}
\begin{tabularx}{\linewidth}{|l|X|X|}
\hline
\rowcolor{popblueL} \textbf{English} & \textbf{Français} & \textbf{Exemple} \\
\hline
beginner & débutant & \eng{She is a beginner in English.} \\
\hline
elementary & élémentaire & \eng{He has an elementary level.} \\
\hline
intermediate & intermédiaire & \eng{I am an intermediate learner.} \\
\hline
upper-intermediate & intermédiaire supérieur & \eng{She reached an upper-intermediate level.} \\
\hline
advanced & avancé & \eng{He is an advanced speaker of English.} \\
\hline
fluent & courant & \eng{She is fluent in English.} \\
\hline
native speaker & locuteur natif & \eng{He speaks like a native speaker.} \\
\hline
\end{tabularx}
\end{center}

\begin{popfigure}
\begin{tikzpicture}
\draw[-{Stealth}, thick, popblue] (0,0) -- (0,6.6);
\foreach \y/\lev/\col in {0.4/beginner/popmuted, 1.5/elementary/popgold, 2.6/intermediate/popgreen, 3.7/upper-intermediate/popblue, 4.8/advanced/poppurple, 5.9/fluent — native speaker/poporange}{
  \fill[\col] (0,\y) circle (0.09);
  \node[anchor=west, align=left, font=\small] at (0.25,\y) {\textbf{\lev}};
}
\node[rotate=90, anchor=south, font=\small\itshape, popmuted] at (-0.5,3.3) {increasing level};
\end{tikzpicture}
\legende{L'échelle des niveaux d'anglais, du débutant au locuteur natif.}
\end{popfigure}

\begin{attention}[Ne dis pas « I have a good level in English »]
C'est un calque du français « j'ai un bon niveau en anglais ». En anglais naturel, dis plutôt :
\begin{itemize}
\item \eng{I have a good command of English.} « J'ai une bonne maîtrise de l'anglais. »
\item \eng{My English is good.} « Mon anglais est bon. »
\item \eng{I speak English well.} « Je parle bien anglais. »
\end{itemize}
De même, on dit \eng{fluent \textbf{in} English} (et jamais « fluent of English »).
\end{attention}

\courssub{Décrire son niveau : les expressions clés}

\begin{exemplebox}[Expressions pour parler de son niveau]
\begin{itemize}
\item \eng{I have a good command of English.} « J'ai une bonne maîtrise de l'anglais. »
\item \eng{I am fluent in English.} « Je parle anglais couramment. »
\item \eng{She is fluent in English.} « Elle parle anglais couramment. »
\item \eng{My English is a bit rusty.} « Mon anglais est un peu rouillé. »
\item \eng{His English is rusty.} « Son anglais est rouillé. »
\item \eng{I can get by in English.} « Je peux me débrouiller en anglais. »
\item \eng{I have an intermediate level.} « J'ai un niveau intermédiaire. »
\end{itemize}
\end{exemplebox}

\begin{methode}[Décrire son niveau en trois étapes]
\begin{enumerate}
\item \textbf{Nommer le niveau} : \eng{I am an intermediate learner.} / \eng{My level is intermediate.}
\item \textbf{Préciser la compétence} : \eng{My speaking is good, but my writing is weak.}
\item \textbf{Indiquer le besoin} : \eng{I need to improve my writing.} / \eng{I want to upgrade my listening skills.}
\end{enumerate}
Modèle complet : \eng{I am an intermediate learner. My speaking is fairly good, but my writing is weak, so I need to improve it.} « J'ai un niveau intermédiaire. Mon expression orale est assez bonne, mais mon expression écrite est faible, donc je dois l'améliorer. »
\end{methode}

\courssub{Fluency, accuracy, proficiency : trois mots à ne pas confondre}

\begin{definition}[Trois notions distinctes]
\begin{itemize}
\item \eng{Fluency} (« flou-en-si ») : l'\emph{aisance}, la facilité à parler sans hésiter — même avec des erreurs.
\item \eng{Accuracy} (« a-kiou-ra-si ») : l'\emph{exactitude}, le fait de parler correctement (grammaire, vocabulaire, prononciation).
\item \eng{Proficiency} : la \emph{maîtrise globale} de la langue, qui combine fluency et accuracy dans les quatre compétences (listening, speaking, reading, writing).
\end{itemize}
\end{definition}

\begin{exemplebox}[En contexte]
\eng{Sandra speaks with fluency, but her accuracy is poor: she makes grammar mistakes.} « Sandra parle avec aisance, mais son exactitude est faible : elle fait des fautes de grammaire. »

\eng{The job requires a high level of English proficiency.} « Le poste exige un haut niveau de maîtrise de l'anglais. »
\end{exemplebox}

\coursec{Formation et développement des compétences}

\courssub{Le champ lexical des besoins et les verbes d'anticipation}

\begin{definition}[Les besoins linguistiques]
\begin{itemize}
\item \eng{language needs} : les besoins en langue ; \eng{language requirements} : les exigences linguistiques (d'un poste, d'un concours).
\item \eng{language skills} : les compétences linguistiques (les quatre : listening, speaking, reading, writing) ; \eng{communication skills} : les compétences en communication.
\item \eng{a strength} : un point fort ; \eng{a weakness} : un point faible ; \eng{a gap} : une lacune, un écart à combler.
\end{itemize}
\eng{She identified a gap between her level and the job requirements.} « Elle a identifié un écart entre son niveau et les exigences du poste. »
\end{definition}

\begin{propriete}[Les verbes d'anticipation et de progression]
Pour parler de sa stratégie de progression, l'anglais utilise des verbes précis :
\begin{itemize}
\item \eng{to assess one's level} : évaluer son niveau — \eng{I assessed my level with a test.}
\item \eng{to identify one's needs} : identifier ses besoins — \eng{She identified her language needs.}
\item \eng{to plan ahead} : planifier à l'avance (intransitif) — \eng{You should plan ahead for the interview.}
\item \eng{to anticipate} : anticiper (transitif) — \eng{We must anticipate our future needs.}
\item \eng{to set goals} : se fixer des objectifs — \eng{He set clear goals for the year.}
\item \eng{to improve} : améliorer — \eng{I want to improve my writing.}
\item \eng{to upgrade one's skills} : mettre ses compétences à niveau — \eng{She upgraded her computer skills.}
\item \eng{to upskill} : se perfectionner, acquérir de nouvelles compétences — \eng{The company encourages staff to upskill.}
\item \eng{to reskill} : se reconvertir, acquérir de nouvelles compétences pour un autre métier.
\end{itemize}
\textbf{Constructions :} \eng{to improve} et \eng{to upgrade} sont suivis d'un complément direct ; \eng{to plan} est suivi de l'infinitif avec \eng{to} : \eng{I plan to take a course.} ; \eng{to anticipate} est suivi d'un nom ou d'un gérondif : \eng{We anticipate needing English.}
\end{propriete}

\begin{popfigure}
\begin{tikzpicture}[mindmap, grow cyclic,
  every node/.style={concept, font=\small},
  concept color=popblue,
  level 1/.append style={level distance=3.2cm, sibling angle=90},
  level 2/.append style={level distance=2.2cm, sibling angle=45}]
\node[concept, text=white, align=center]{\textbf{Language}\\\textbf{needs}}
  child[concept color=popgreen] { node {assess} child { node[font=\scriptsize] {placement test} } child { node[font=\scriptsize] {self-assessment} } }
  child[concept color=poporange] { node {identify} child { node[font=\scriptsize] {strengths} } child { node[font=\scriptsize] {weaknesses} } }
  child[concept color=poppurple] { node {plan} child { node[font=\scriptsize] {set goals} } child { node[font=\scriptsize] {plan ahead} } }
  child[concept color=popgold] { node {improve} child { node[font=\scriptsize] {take a course} } child { node[font=\scriptsize] {practise daily} } };
\end{tikzpicture}
\legende{Carte mentale : anticiper ses besoins linguistiques en quatre verbes.}
\end{popfigure}

\courssub{Le vocabulaire de la formation professionnelle}

\begin{definition}[Mots et expressions de la formation continue]
\begin{itemize}
\item \eng{on-the-job training} : formation en poste, formation sur le tas.
\item \eng{a workshop} : un atelier (de travail pratique).
\item \eng{a seminar} : un séminaire.
\item \eng{in-service training} : formation continue (pour employés en poste).
\item \eng{professional development} : développement professionnel.
\item \eng{lifelong learning} : apprentissage tout au long de la vie.
\item \eng{a mentor} / \eng{a mentee} : un mentor / un mentoré.
\item \eng{an appraisal} / \eng{a performance review} : une évaluation annuelle / un entretien d'évaluation.
\item \eng{a personal development plan (PDP)} : un plan de développement personnel.
\end{itemize}
\eng{During my appraisal, I discussed my language needs and asked for a workshop on business writing.} « Lors de mon entretien d'évaluation, j'ai évoqué mes besoins en langue et demandé un atelier sur la rédaction professionnelle. »
\end{definition}

\begin{methode}[Rédiger un plan de développement personnel (PDP) en anglais]
\begin{enumerate}
\item \textbf{Current level} : \eng{My current level is intermediate (B1).}
\item \textbf{Target level} : \eng{I aim to reach upper-intermediate (B2) in six months.}
\item \textbf{Actions} : \eng{I will attend a weekly business English course and practise writing emails daily.}
\item \textbf{Resources} : \eng{I will use the British Council website and a grammar app.}
\item \textbf{Timeline} : \eng{I will review my progress at the end of each month.}
\end{enumerate}
\end{methode}

\courssub{Les outils d'évaluation}

\begin{definition}[Tests, certificats et grilles]
\begin{itemize}
\item \eng{a placement test} : un test de niveau, passé avant d'entrer dans une formation pour être placé dans le bon groupe.
\item \eng{a language certificate} : un certificat de langue, document officiel attestant le niveau.
\item \eng{TOEFL} (\eng{Test of English as a Foreign Language}) et \eng{IELTS} (\eng{International English Language Testing System}) : deux examens internationaux qui mesurent l'\eng{English proficiency} ; ils sont souvent exigés pour étudier ou travailler dans un pays anglophone.
\item \eng{a self-assessment grid} : une grille d'auto-évaluation (ex. : grille CECRL — « I can introduce myself », « I can write a formal letter » : tu coches ce que tu sais faire).
\end{itemize}
\eng{She took the IELTS to prove her English proficiency.} « Elle a passé l'IELTS pour prouver sa maîtrise de l'anglais. »
\end{definition}

\begin{savaistu}[L'anglais, une stratégie de carrière au Cameroun]
Au Cameroun, de nombreux concours professionnels — comme l'ENAM ou l'IRIC — comportent une épreuve d'anglais, même pour les candidats francophones. Anticiper ses besoins linguistiques (évaluer son niveau tôt, combler ses lacunes, passer un certificat) est donc une véritable stratégie de carrière, pas seulement une affaire d'école. Le bilinguisme officiel du pays fait de l'anglais un atout décisif sur le marché du travail.
\end{savaistu}

\coursec{L'anglais dans la vie professionnelle}

\courssub{Écouter : un entretien d'évaluation (script d'écoute)}

\begin{textebox}[Appraisal interview — listening script]
\textbf{Manager:} Good morning, Paul. Thanks for coming. Let's start with your self-assessment. How would you rate your English level today?
\textbf{Paul:} Good morning. I'd say I'm at an intermediate level. I can get by in meetings, but I struggle with writing formal reports.
\textbf{Manager:} I see. The new project requires a good command of written English for the monthly newsletter. What are your language needs?
\textbf{Paul:} I've identified a gap in business vocabulary and report structure. I plan to take an online course on business writing next month.
\textbf{Manager:} That's a good initiative. We can also arrange a workshop with a language trainer. Let's set a goal: by the end of the quarter, you should be able to draft the newsletter independently.
\textbf{Paul:} That sounds realistic. I'll upgrade my skills and keep you posted.
\end{textebox}

\courssub{Parler : jeu de rôle — l'entretien d'évaluation}

\begin{experience}[Role-play: Appraisal Meeting]
\textbf{Consigne.} En binôme. L'un joue le manager, l'autre l'employé (Paul ou un autre prénom). Utilisez le vocabulaire de la leçon.
\begin{enumerate}
\item \textbf{Manager :} Demandez l'auto-évaluation (\eng{“How would you rate your English?”}).
\item \textbf{Employé :} Donnez votre niveau, citez une force et une faiblesse (\eng{“I'm intermediate. My speaking is fluent but my writing is rusty.”}).
\item \textbf{Manager :} Annoncez un besoin du service (\eng{“We need someone to write reports in English.”}).
\item \textbf{Employé :} Identifiez l'écart (\eng{“I have a gap in business vocabulary.”}) et proposez un plan (\eng{“I plan to take a course / upskill.”}).
\item \textbf{Manager :} Validez et fixez un objectif SMART (\eng{“Let's set a goal: draft a report by June.”}).
\end{enumerate}
\textbf{Durée :} 3 minutes par binôme. Changez de rôle. Le professeur note l'usage de \eng{command of}, \eng{fluent in}, \eng{gap}, \eng{upgrade}, \eng{set goals}.
\end{experience}

\coursec{Faux amis et pièges des francophones}

\courssub{Les faux amis du monde professionnel}

\begin{attention}[Faux amis fréquents — Attention danger !]
\begin{center}
\begin{tabularx}{\linewidth}{|l|l|l|X|}
\hline
\rowcolor{popblueL} \textbf{Mot anglais} & \textbf{Ne signifie PAS} & \textbf{Vrai sens} & \textbf{Français correct} \\
\hline
\eng{actually} & actuellement & en fait, réellement & \eng{currently} = actuellement \\
\hline
\eng{eventually} & éventuellement & finalement, à la fin & \eng{possibly} = éventuellement \\
\hline
\eng{formation} & formation (pro) & formation (géologie, forme) & \eng{training} = formation pro \\
\hline
\eng{stage} & stage (entreprise) & scène, étape & \eng{internship} (US) / \eng{work placement} (UK) = stage \\
\hline
\eng{resume} (US) / \eng{CV} (UK) & résumé (court) & curriculum vitæ & \eng{summary} = résumé \\
\hline
\eng{salary} & salaire (terme générique) & salaire mensuel/annuel (cadre) & \eng{wages} = salaire (horaire, ouvrier) \\
\hline
\eng{competition} & compétition (sport) & concurrence (économie), concours & \eng{contest} = compétition sportive \\
\hline
\eng{opportunity} & opportunité (hasard) & occasion favorable, chance & \eng{opportunism} = opportunisme \\
\hline
\end{tabularx}
\end{center}
\end{attention}

\courssub{Pièges grammaticaux et prépositionnels}

\begin{attention}[Prépositions et collocations : les erreurs classiques]
\begin{itemize}
\item \eng{fluent \textbf{in} English} (pas \eng{of}, pas \eng{at}).
\item \eng{good \textbf{at} speaking} / \eng{weak \textbf{in} writing} (pas \eng{in speaking} pour \eng{good}, pas \eng{at writing} pour \eng{weak} — bien que \eng{weak at} soit possible, \eng{weak in} est plus courant pour une matière).
\item \eng{a good command \textbf{of} English} (pas \eng{in}, pas \eng{on}).
\item \eng{proficient \textbf{in} English} (pas \eng{at}, pas \eng{of}).
\item \eng{to apply \textbf{for} a job} (pas \eng{to}).
\item \eng{to attend \textbf{a meeting}} (verbe transitif direct, pas de préposition : \eng{attend to} = s'occuper de).
\item \eng{to take \textbf{a test}} / \eng{sit \textbf{an exam}} (UK) (pas \eng{pass an exam} pour « passer » = présenter ; \eng{pass} = réussir).
\end{itemize}
\end{attention}

\coursec{Familles de mots, prononciation et orthographe}

\courssub{Dérivation morphologique (word formation)}

\begin{propriete}[Familles de mots clés du thème]
Maîtriser la dérivation permet de passer du verbe au nom, à l'adjectif ou à l'adverbe sans faute d'orthographe.
\begin{center}
\begin{tabular}{|l|l|l|l|l|}
\hline
\rowcolor{popblueL} \textbf{Verbe} & \textbf{Nom (action/état)} & \textbf{Nom (agent)} & \textbf{Adjectif} & \textbf{Adverbe} \\
\hline
assess & assessment & assessor & assessable & — \\
\hline
improve & improvement & — & improved / improving & — \\
\hline
require & requirement & — & required / requiring & — \\
\hline
identify & identification & identifier & identifiable & identifiably \\
\hline
anticipate & anticipation & — & anticipated / anticipatory & anticipatorily \\
\hline
communicate & communication & communicator & communicative & communicatively \\
\hline
evaluate & evaluation & evaluator & evaluative & evaluatively \\
\hline
proficiency (n) & — & — & proficient & proficiently \\
\hline
fluency (n) & — & — & fluent & fluently \\
\hline
accuracy (n) & — & — & accurate & accurately \\
\hline
\end{tabular}
\end{center}
\textbf{Rappel orthographe (UK) :} \eng{practise} (verbe) vs \eng{practice} (nom) ; \eng{programme} (UK, hors informatique) vs \eng{program} (US / informatique) ; \eng{centre} (UK) vs \eng{center} (US).
\end{propriete}

\courssub{Prononciation : accent tonique et sons pièges}

\begin{definition}[Points de vigilance à l'oral]
\begin{itemize}
\item \eng{Proficiency} /prəˈfɪʃnsi/ → accent sur \textbf{FI} (« pro-\textbf{FI}-chen-si »). Ne pas dire « pro-fi-CIEN-cy ».
\item \eng{Assessment} /əˈsesmənt/ → accent sur \textbf{SES} (« a-\textbf{SES}-ment »). Le \eng{ss} se prononce /s/.
\item \eng{Requirement} /rɪˈkwɪəmənt/ → accent sur \textbf{KWI} (« ri-\textbf{KWI}-re-ment »). Attention au \eng{ui} = /waɪ/ ou /wɪ/.
\item \eng{Anticipate} /ænˈtɪsɪpeɪt/ → accent sur \textbf{TI} (« an-\textbf{TI}-ci-peit »). Le \eng{ci} = /sɪ/.
\item \eng{Upgrade} /ˌʌpˈɡreɪd/ → accent secondaire sur \textbf{UP}, principal sur \textbf{GRADE}.
\item \eng{Rusty} /ˈrʌsti/ → \eng{u} = /ʌ/ (comme dans \eng{cut}), pas /y/.
\item \eng{Command} /kəˈmɑːnd/ → accent sur \textbf{MAND} (« co-\textbf{MAND} »).
\end{itemize}
\end{definition}

\coursec{Réemploi en contexte et production guidée}

\courssub{Exercice résolu : consolidation du vocabulaire}

\begin{exoresolu}[Complète avec le vocabulaire de la leçon]
Énoncé. Complète chaque phrase avec le mot ou l'expression qui convient : \eng{fluent}, \eng{placement test}, \eng{rusty}, \eng{command}, \eng{weakness}, \eng{set goals}, \eng{upskill}, \eng{internship}.
\begin{enumerate}
\item Before the course, every student takes a \ldots\ to find the right group.
\item « I haven't spoken English for ten years, so my English is a bit \ldots »
\item She is \ldots\ in English: she speaks easily and correctly.
\item He has a good \ldots\ of English, especially in writing.
\item Grammar is my main \ldots ; I must work on it.
\item To progress, you should \ldots\ clear \ldots\ for each term.
\item Many graduates do an \ldots\ to gain professional experience.
\item The company pays for employees to \ldots\ through online certifications.
\end{enumerate}
\tcblower
Solution détaillée.
\begin{enumerate}
\item \eng{placement test} — le test de niveau place l'élève dans le bon groupe.
\item \eng{rusty} — « rouillé » : on n'a pas pratiqué depuis longtemps.
\item \eng{fluent} — construction : \eng{fluent in} + langue.
\item \eng{command} — expression figée : \eng{a good command of English}.
\item \eng{weakness} — un point faible ; le contraire est \eng{a strength}.
\item \eng{set goals} — expression figée : \eng{to set (clear) goals}.
\item \eng{internship} (US) / \eng{work placement} (UK) — faux ami : \eng{stage} $\neq$ stage pro.
\item \eng{upskill} — verbe intransitif : se perfectionner, monter en compétences.
\end{enumerate}
\end{exoresolu}

\courssub{Exercices d'entraînement}

\begin{exercice}[À toi de jouer]
\begin{enumerate}
\item Traduis en anglais : « Mon anglais est un peu rouillé, mais je peux me débrouiller. »
\item Décris ton propre niveau en trois phrases, en suivant la méthode : niveau, compétence, besoin.
\item Trouve les deux mots de la leçon désignant des tests internationaux mesurant la maîtrise de l'anglais.
\item Corrige la phrase : \eng{I have a good level in English.}
\item Transforme les verbes entre parenthèses en la forme nominale (nom d'action) adaptée : \eng{(assess)}, \eng{(improve)}, \eng{(identify)}, \eng{(require)}, \eng{(communicate)}.
\item \textbf{Production écrite :} Rédige ton \eng{Personal Development Plan (PDP)} pour l'anglais sur 6 mois (80–100 mots). Utilise la méthode de la leçon : niveau actuel, niveau visé, actions, ressources, échéancier.
\end{enumerate}
\end{exercice}

\begin{corrige}[Exercice « À toi de jouer »]
\begin{enumerate}
\item \eng{My English is a bit rusty, but I can get by.}
\item Réponse libre ; modèle : \eng{I am an intermediate learner. My listening is good, but my writing is weak, so I need to improve it.}
\item \eng{TOEFL} et \eng{IELTS}.
\item \eng{I have a good command of English.} ou \eng{My English is good.} ou \eng{I speak English well.}
\item \eng{assessment}, \eng{improvement}, \eng{identification}, \eng{requirement}, \eng{communication}.
\item \textbf{Modèle de PDP :}
\eng{My current level is intermediate (B1). I aim to reach upper-intermediate (B2) in six months. To achieve this, I will attend a weekly business English workshop at the British Council centre in Yaoundé. I will also practise writing formal emails every day using online templates. My resources are the “Business English Pod” podcast and the “Grammar in Use” app. I will review my progress at the end of each month with a self-assessment grid.}
\end{enumerate}
\end{corrige}

\begin{aretenir}[L'essentiel de la leçon — Lesson 11]
\begin{itemize}
\item \textbf{Échelle des niveaux :} \eng{beginner} → \eng{elementary} → \eng{intermediate} → \eng{upper-intermediate} → \eng{advanced} → \eng{fluent} / \eng{native speaker}.
\item \textbf{Décrire son niveau :} \eng{I have a good command of English}, \eng{I am fluent in English}, \eng{My English is rusty}, \eng{I can get by in English}.
\item \textbf{Trois notions clés :} \eng{Fluency} (aisance) $\neq$ \eng{Accuracy} (exactitude) ; \eng{Proficiency} = maîtrise globale (4 compétences).
\item \textbf{Verbes d'anticipation :} \eng{assess}, \eng{identify}, \eng{plan ahead}, \eng{anticipate}, \eng{set goals}, \eng{improve}, \eng{upgrade}, \eng{upskill}, \eng{reskill}.
\item \textbf{Outils :} \eng{placement test}, \eng{language certificate}, \eng{TOEFL}, \eng{IELTS}, \eng{self-assessment grid}, \eng{PDP}.
\item \textbf{Formation pro :} \eng{training}, \eng{workshop}, \eng{internship} / \eng{work placement}, \eng{appraisal}, \eng{mentor}.
\item \textbf{Pièges majeurs (francophones) :}
\begin{itemize}
\item \eng{fluent in} / \eng{command of} / \eng{good at} / \eng{weak in} / \eng{proficient in}.
\item \eng{actually} = en fait ; \eng{currently} = actuellement.
\item \eng{eventually} = finalement ; \eng{possibly} = éventuellement.
\item \eng{training} = formation pro ; \eng{formation} = géologie/structure.
\item \eng{internship} = stage pro ; \eng{stage} = scène/étape.
\item \eng{apply for a job} ; \eng{attend a meeting} (sans préposition) ; \eng{sit an exam} (UK) / \eng{take a test}.
\end{itemize}
\item \textbf{Orthographe UK :} \eng{practise} (v), \eng{practice} (n) ; \eng{programme} (général), \eng{program} (informatique) ; \eng{centre}.
\end{itemize}
\end{aretenir}

\coursec{Formation et développement des compétences}

Dans la section précédente, nous avons appris à évaluer notre niveau d'anglais et à identifier nos besoins linguistiques au travail. Une fois le besoin identifié (\eng{need identified}), l'étape suivante est logique : se former. C'est tout l'objet de cette deuxième section : découvrir le vocabulaire anglais de la \cle{formation professionnelle}, apprendre les collocations qui vont avec, et savoir demander une formation à son employeur dans un anglais correct et naturel.

\courssub{Texte support : Mr Etoundi goes back to school}

Lisons d'abord ce court texte. Soulignez mentalement tous les mots qui appartiennent au champ lexical de la formation.

\begin{textebox}[Mr Etoundi's training journey]
Mr Etoundi is an accountant in a big import-export company in Yaoundé. Last year, his director announced that the company was about to sign a contract with a Ghanaian firm. All the negotiations would be in English, and Mr Etoundi realised that his school English was not good enough.

He decided to act quickly. First, he enrolled in an online English course offered by a well-known language school. The course included two live workshops every week and a lot of listening exercises. His trainer, Mrs Fon, was very patient and always corrected his pronunciation. There were twelve trainees in his virtual class, including a nurse from Bamenda and a lawyer from Douala.

After three months, Mr Etoundi felt he needed more practice, so he also signed up for evening classes at a language centre near his office. Twice a week, from 6 p.m. to 8 p.m., he attended lessons with a tutor who helped him brush up his grammar and enrich his professional vocabulary. Some of his colleagues preferred distance learning because they travelled a lot, but Mr Etoundi liked meeting other learners face to face.

After six months of training, he completed the programme and obtained a certificate in business English. His director was impressed: Mr Etoundi could now chair meetings in English, write clear reports and welcome Ghanaian partners at the airport. Thanks to this skills upgrade, he was promoted a few months later.

“Investing in training is investing in your career,” he often tells young employees. “Never wait until you need English tomorrow. Start learning it today.”
\end{textebox}

\textbf{Glossaire}

\begin{center}
\begin{tabularx}{\linewidth}{|l|l|X|}
\hline
\rowcolor{popblueL} Mot anglais & Nature & Traduction française \\
\hline
accountant & nom & comptable \\
\hline
to be about to & expression & être sur le point de \\
\hline
firm & nom & entreprise, cabinet \\
\hline
to enrol in & verbe & s'inscrire à \\
\hline
workshop & nom & atelier (de formation) \\
\hline
trainee & nom & stagiaire, personne en formation \\
\hline
to sign up for & verbe & s'inscrire à \\
\hline
evening classes & nom & cours du soir \\
\hline
to brush up & verbe & réviser, réactiver \\
\hline
distance learning & nom & formation à distance \\
\hline
certificate & nom & certificat, attestation \\
\hline
to chair a meeting & expression & présider une réunion \\
\hline
skills upgrade & nom & montée en compétences \\
\hline
to be promoted & verbe & être promu \\
\hline
\end{tabularx}
\end{center}

\textbf{Questions de compréhension}

\begin{enumerate}
\item Why did Mr Etoundi decide to improve his English?
\item What two types of training did he follow?
\item Who was Mrs Fon?
\item What did Mr Etoundi obtain at the end of the programme?
\item What happened to his career after the training?
\end{enumerate}

\courssub{Le champ lexical de la formation}

Le texte contient presque tout le vocabulaire essentiel du thème. Observons-le et organisons-le.

\begin{definition}[Les types de formation]
\begin{itemize}
\item \eng{training} : la formation (mot général, indénombrable). \eng{Staff training is a priority.} « La formation du personnel est une priorité. »
\item \eng{a training course} : un stage, un cours de formation. \eng{She is following a training course in accounting.} « Elle suit une formation en comptabilité. »
\item \eng{a workshop} : un atelier pratique, souvent court et interactif.
\item \eng{a seminar} : un séminaire, réunion de travail avec exposés et discussions.
\item \eng{a refresher course} : un cours de remise à niveau.
\item \eng{evening classes} : les cours du soir.
\item \eng{an online course} : un cours en ligne.
\item \eng{distance learning} : la formation à distance.
\item \eng{a language school} : une école de langues.
\item \eng{on-the-job training} : la formation en poste, formation interne.
\end{itemize}
\end{definition}

\begin{definition}[Refresher course]
A \cle{refresher course} is a short course to practise and update skills you have not used for a long time. C'est un \textbf{cours de remise à niveau} : on ne part pas de zéro, on « rafraîchit » des connaissances oubliées ou rouillées.\\
\eng{After ten years without speaking English, she took a refresher course.} « Après dix ans sans parler anglais, elle a suivi une remise à niveau. »
\end{definition}

\begin{definition}[Le personnel de formation et les résultats]
\textbf{Les personnes :} \eng{trainer} (formateur), \eng{trainee} (stagiaire, personne formée), \eng{tutor} (tuteur), \eng{coach} (coach, accompagnateur), \eng{instructor} (instructeur, moniteur), \eng{mentor} (mentor, conseiller expérimenté).\\
\textbf{Les résultats :} \eng{certificate} (certificat), \eng{diploma} (diplôme), \eng{qualification} (qualification, titre), \eng{skills upgrade} (montée en compétences), \eng{professional development} (développement professionnel), \eng{enrolment} (inscription).
\end{definition}

\begin{attention}[Le mot « formation » : faux ami célèbre]
En français, « suivre une formation » se dit \eng{to follow a training course}, \eng{to take a course} ou simplement \eng{training}. Le mot anglais \eng{formation} existe, mais il signifie « formation géologique » (\eng{rock formation}) ou « formation militaire » (\eng{the planes flew in formation}). Ne dites jamais \emph{I need a formation in English} : dites \eng{I need training in English} ou \eng{I need to take an English course}.
\end{attention}

\courssub{Les collocations essentielles}

En anglais, un mot ne voyage jamais seul : il a ses partenaires habituels. Ce sont les \cle{collocations}. Les apprendre par cœur, c'est parler un anglais naturel au lieu d'un français traduit mot à mot.

\begin{propriete}[Collocations figées avec « course »]
On dit \eng{to attend a course} (assister à un cours) et \textbf{jamais} \emph{to assist a course} : le verbe \eng{to assist} signifie « aider » en anglais (faux ami avec « assister à »). On dit aussi \eng{to take a course} (suivre un cours), tout à fait correct, et \eng{to enrol in a course} (s'inscrire à un cours). À la fin, on \eng{completes a training programme} (on termine un programme de formation).
\end{propriete}

\begin{center}
\begin{tabularx}{\linewidth}{|X|X|}
\hline
\rowcolor{popblueL} Collocation anglaise & Traduction française \\
\hline
to attend a training course & suivre / assister à un stage de formation \\
\hline
to enrol in a course & s'inscrire à un cours \\
\hline
to take a course & suivre un cours \\
\hline
to sign up for evening classes & s'inscrire aux cours du soir \\
\hline
to brush up one's English & rafraîchir / réactiver son anglais \\
\hline
to complete a training programme & terminer un programme de formation \\
\hline
to obtain a certificate & obtenir un certificat \\
\hline
to improve one's skills & améliorer ses compétences \\
\hline
to fund / to sponsor training & financer une formation \\
\hline
\end{tabularx}
\end{center}

\begin{definition}[Trois expressions idiomatiques à connaître]
\begin{itemize}
\item \eng{to brush up (one's English)} : réviser, réactiver une langue qu'on a apprise mais peu pratiquée. \eng{I need to brush up my English before the conference.} « Je dois rafraîchir mon anglais avant la conférence. »
\item \eng{to pick up a language} : apprendre une langue naturellement, sans cours, par immersion. \eng{She picked up English in Lagos.} « Elle a appris l'anglais naturellement à Lagos. »
\item \eng{to keep up with} : se maintenir à niveau, suivre le rythme. \eng{He reads the news in English to keep up with new vocabulary.} « Il lit les actualités en anglais pour rester à niveau. »
\end{itemize}
\end{definition}

\courssub{Le parcours de formation : du besoin au diplôme}

Voici le cheminement typique d'un employé qui anticipe ses besoins en anglais, comme Mr Etoundi.

\begin{popfigure}
\begin{center}
\begin{tikzpicture}[node distance=0.55cm and 1.1cm,
etape/.style={rectangle, rounded corners, draw=popblue, fill=popblueL, align=center, minimum height=0.9cm, text width=3.1cm, font=\small},
fleche/.style={-{Stealth[length=3mm]}, thick, poporange}]
\node[etape, align=center] (n1){Need identified\\ \emph{besoin identifié}};
\node[etape, right=of n1, align=center] (n2){Choose a course\\ evening classes, online course, workshop};
\node[etape, right=of n2, align=center] (n3){Attend training\\ \emph{suivre la formation}};
\node[etape, right=of n3, align=center] (n4){Obtain a certificate\\ \emph{obtenir un certificat}};
\node[etape, right=of n4, align=center] (n5){Career progress\\ \emph{évolution de carrière}};
\draw[fleche] (n1) -- (n2);
\draw[fleche] (n2) -- (n3);
\draw[fleche] (n3) -- (n4);
\draw[fleche] (n4) -- (n5);
\end{tikzpicture}
\end{center}
\legende{Le parcours de formation professionnelle en cinq étapes}
\end{popfigure}

\courssub{Demander une formation à son employeur}

Savoir formuler une demande de formation est une compétence professionnelle précieuse. Voici des phrases types, du plus direct au plus formel.

\begin{exemplebox}[Phrases types pour demander une formation]
\begin{itemize}
\item \eng{I would like to attend a refresher course in business English.} « J'aimerais suivre une remise à niveau en anglais des affaires. »
\item \eng{Could I enrol in the evening classes offered by the language centre?} « Pourrais-je m'inscrire aux cours du soir proposés par le centre de langues ? »
\item \eng{I think an online course would help me communicate better with our Ghanaian partners.} « Je pense qu'un cours en ligne m'aiderait à mieux communiquer avec nos partenaires ghanéens. »
\item \eng{Would the company be willing to pay for this training programme?} « L'entreprise accepterait-elle de financer ce programme de formation ? »
\item \eng{This certificate would be useful for my professional development.} « Ce certificat serait utile pour mon développement professionnel. »
\end{itemize}
\end{exemplebox}

\begin{methode}[Rédiger une demande de formation en quatre étapes]
\begin{enumerate}
\item \textbf{Présenter le besoin} : \eng{Our company is working more and more with English-speaking clients.}
\item \textbf{Proposer une solution précise} : \eng{I would like to take an online course in business English.}
\item \textbf{Montrer l'avantage pour l'entreprise} : \eng{This training would help me write reports and chair meetings in English.}
\item \textbf{Formuler la demande poliment} : \eng{I would be grateful if the company could support this training.}
\end{enumerate}
\end{methode}

\begin{exoresolu}[Compléter avec la bonne collocation]
Énoncé — Complétez chaque phrase avec le verbe correct : \eng{attend}, \eng{enrol in}, \eng{brush up}, \eng{pick up}, \eng{complete}.\\
1. My sister wants to \ldots\ a course in computer science.\\
2. Before the interview, I must \ldots\ my English.\\
3. All employees can \ldots\ the workshop on Friday.\\
4. He \ldots\ some Spanish while working in Equatorial Guinea.\\
5. She will \ldots\ her training programme in December.
\tcblower
Solution détaillée :\\
1. \eng{enrol in} — on s'inscrit \emph{à} un cours : \eng{enrol in a course}.\\
2. \eng{brush up} — réactiver son anglais avant un entretien : collocation idiomatique.\\
3. \eng{attend} — assister à un atelier : \eng{attend the workshop} (jamais \emph{assist}).\\
4. \eng{picked up} — il a appris naturellement, par immersion, au travail ; attention au prétérit : \eng{pick} devient \eng{picked}.\\
5. \eng{complete} — terminer un programme : \eng{complete a training programme}.
\end{exoresolu}

\begin{attention}[Erreurs fréquentes des francophones]
\begin{itemize}
\item \emph{I assist to a course.} Double erreur ! Dites \eng{I attend a course} (sans \eng{to}, et avec \eng{attend}).
\item \emph{I made a formation.} Dites \eng{I did some training} ou \eng{I took a course}.
\item \emph{I learned English in a school of languages.} Calque du français : dites \eng{a language school}.
\item \emph{She followed a stage.} \eng{Stage} n'existe pas avec ce sens en anglais : dites \eng{training course} ou \eng{internship} (stage pratique en entreprise).
\item \emph{He passed a training.} Dites \eng{He completed a training course} ou \eng{He passed an exam} (réussir un examen).
\end{itemize}
\end{attention}

\courssub{Familles de mots, prononciation et orthographe}

Travailler sur les familles de mots permet d'enrichir son vocabulaire de manière exponentielle. Voici les dérivations principales du champ lexical de la formation.

\begin{definition}[Familles de mots clés]
\begin{center}
\begin{tabular}{|l|l|l|l|}
\hline
\rowcolor{popblueL} Verbe & Nom (action/état) & Nom (personne) & Adjectif / Participe \\
\hline
to train & training & trainer / trainee & trained \\
\hline
to enrol & enrolment & enrollee & enrolled \\
\hline
to attend & attendance & attendee & — \\
\hline
to complete & completion & — & complete / completed \\
\hline
to qualify & qualification & — & qualified \\
\hline
to develop & development & developer & developed / developing \\
\hline
to certify & certification / certificate & — & certified \\
\hline
to assess & assessment & assessor & assessed \\
\hline
\end{tabular}
\end{center}
\end{definition}

\begin{propriete}[Points d'orthographe et de prononciation à surveiller]
\begin{itemize}
\item \textbf{Enrol (UK) vs Enroll (US)} : Au Cameroun (anglais britannique), on écrit \eng{enrol}, \eng{enrolment}, \eng{enrolling} (un seul \eng{l} à l'infinitif et au nom, double \eng{l} en \eng{-ing}).
\item \textbf{Programme (UK) vs Program (US)} : On écrit \eng{a training programme} (nom), mais \eng{a computer program} (informatique).
\item \textbf{Certificate} : Se prononce « sèr-ti-fi-kèt » (\eng{/səˈtɪfɪkət/}), pas *certifi-keït. L'accent est sur la 2e syllabe.
\item \textbf{Qualification} : Se prononce « quo-li-fi-ké-chn » (\eng{/ˌkwɒlɪfɪˈkeɪʃn/}). Attention au \eng{ti} = \eng{/ʃ/}.
\item \textbf{Trainee} : Accent sur la dernière syllabe : « traï-ni » (\eng{/ˌtreɪˈniː/}). Ne pas confondre avec \eng{trainer} (« traï-né », \eng{/ˈtreɪnə/}).
\item \textbf{Refresher} : Se prononce « ri-fre-cha » (\eng{/rɪˈfreʃə/}), le \eng{sh} s'entend.
\item \textbf{Workshop} : Mot composé, accent sur le premier élément : « worck-chop » (\eng{/ˈwɜːkʃɒp/}).
\end{itemize}
\end{propriete}

\begin{exemplebox}[Exercices de dérivation]
Complétez le tableau en trouvant la forme manquante (vérifiez l'orthographe UK).
\begin{center}
\begin{tabular}{|l|l|l|l|}
\hline
\rowcolor{popblueL} Verbe & Nom (chose/action) & Personne & Adjectif \\
\hline
to assess & \ldots & \ldots & assessed \\
\hline
\ldots & certification & \ldots & certified \\
\hline
to develop & development & \ldots & \ldots \\
\hline
\ldots & \ldots & trainee & trained \\
\hline
to qualify & \ldots & \ldots & qualified \\
\hline
\end{tabular}
\end{center}
\tcblower
\textbf{Corrigé :} assessment / assessor ; to certify / — / developer / developing ; to train / training / trainer / — ; qualification / —.
\end{exemplebox}

\courssub{Réemploi en contexte et production guidée}

\begin{exercice}[Mise en situation : La demande de formation]
Vous êtes \eng{Marketing Assistant} dans une entreprise à Douala. Votre directeur général a annoncé l'ouverture d'un nouveau marché au Nigeria (pays anglophone). Vous avez besoin de suivre une formation en \eng{Business English} pour gérer les emails et les appels clients.
\begin{enumerate}
\item Rédigez un email formel (120--150 mots) à votre responsable RH (\eng{HR Manager}) pour demander le financement d'une formation.
\item Utilisez au moins 6 mots ou expressions du vocabulaire de cette leçon (\eng{enrol in}, \eng{brush up}, \eng{training course}, \eng{certificate}, \eng{professional development}, \eng{skills upgrade}, \eng{workshop}, \eng{distance learning}, etc.).
\item Structurez votre email selon la méthode en 4 étapes vue plus haut.
\end{enumerate}
\end{exercice}

\begin{corrige}[Exercice : Mise en situation]
\textbf{Modèle de réponse :}
\eng{Subject: Request for Business English Training Funding}

\eng{Dear Mr Mbarga,}

\eng{I am writing to request your support for my professional development. As you know, our company is about to sign a contract with a Nigerian firm and all communications will be in English.}

\eng{I would like to enrol in an intensive Business English training course offered by the British Council in Douala. The programme includes weekly workshops and distance learning modules. It would allow me to brush up my grammar and acquire specific vocabulary for negotiations and report writing.}

\eng{Obtaining a recognised certificate would be a major skills upgrade for me and a real asset for the department. I would be grateful if the company could fund this training.}

\eng{I am available to discuss this further at your convenience.}

\eng{Yours sincerely,}

\eng{Jean Dupont}
\eng{Marketing Assistant}
\end{corrige}

\begin{savaistu}[L'anglais en ligne au Cameroun]
De nombreux jeunes Camerounais suivent aujourd'hui des cours d'anglais en ligne gratuits — par exemple ceux du \eng{British Council} ou les \eng{MOOCs} (\eng{Massive Open Online Courses}) — pour anticiper les exigences du marché du travail. Avec un simple téléphone et une connexion, on peut \eng{brush up one's English}, obtenir des certificats et rendre son CV plus attractif, sans quitter Yaoundé, Bafoussam ou Maroua. C'est le \eng{distance learning} au service de la carrière !
\end{savaistu}

\begin{aretenir}[L'essentiel de la section]
\begin{itemize}
\item « Formation » se dit \eng{training}, jamais \emph{formation}.
\item On \eng{attends} ou \eng{takes} a course ; on \eng{enrols in} a course ; on \eng{signs up for} evening classes ; on \eng{completes} a programme.
\item \eng{To brush up} = réactiver ; \eng{to pick up} = apprendre naturellement ; \eng{to keep up with} = se maintenir à niveau.
\item Le formateur est le \eng{trainer}, le stagiaire le \eng{trainee} ; à la fin, on obtient un \eng{certificate}.
\item Orthographe UK : \eng{enrol / enrolment}, \eng{programme}.
\item Prononciation : \eng{certificate} (sèr-ti-fi-kèt), \eng{trainee} (traï-ni), \eng{workshop} (worck-chop).
\item Pour demander une formation : \eng{I would like to attend a refresher course in business English.}
\end{itemize}
\end{aretenir}

\coursec{L'anglais dans la vie professionnelle}

\courssub{Évaluer son niveau et ses besoins linguistiques}

Avant de se former, il faut savoir où l'on en est. Cette première étape consiste à identifier son niveau actuel (selon le CECRL) et à cibler les compétences à renforcer pour son poste ou son projet professionnel.

\begin{textebox}[Auto-évaluation : identifier ses besoins]
\textbf{HR Officer:} Before we validate your training plan, please complete this self-assessment grid. What is your current CEFR level?

\textbf{Employee:} I would say B1 overall. I can understand the main points of clear standard input on familiar matters. But I struggle with \eng{technical jargon} during \eng{video conferences} with our Lagos office.

\textbf{HR Officer:} Which skills need priority? Writing reports? Speaking in meetings?

\textbf{Employee:} Definitely \eng{oral fluency} and \eng{drafting formal emails}. I also need to prepare for a \eng{proficiency test} (TOEIC) next quarter.
\end{textebox}

\begin{center}
\begin{tabularx}{\linewidth}{|l|X|l|}\hline
\rowcolor{popblueL} \textbf{English} & \textbf{Nature} & \textbf{Français} \\ \hline
self-assessment grid & nom & grille d'auto-évaluation \\ \hline
CEFR level & nom & niveau CECRL (A1--C2) \\ \hline
technical jargon & nom & jargon technique \\ \hline
video conference & nom & visioconférence \\ \hline
oral fluency & nom & aisance à l'oral \\ \hline
drafting formal emails & nom & rédaction de courriels formels \\ \hline
proficiency test & nom & test de certification (niveau) \\ \hline
to prioritize & verbe & prioriser \\ \hline
target language & nom & langue cible \\ \hline
\end{tabularx}
\end{center}

\begin{methode}[Décrire son niveau et ses besoins en entretien]
Pour répondre à \eng{“What is your English level?”} ou \eng{“What are your language needs?”}, structure ta réponse en trois temps :
\begin{enumerate}
\item \textbf{Niveau CECRL} : \eng{I am at a B2 level (upper-intermediate).}
\item \textbf{Points forts / Faiblesses} : \eng{I am comfortable writing reports but I lack fluency in spontaneous conversation.}
\item \textbf{Objectif concret} : \eng{I need to chair meetings in English next semester, so I am focusing on meeting language.}
\end{enumerate}
\end{methode}

\courssub{Formation et développement des compétences}

Une fois les besoins identifiés, place à la formation. Le vocabulaire suivant couvre les modalités d'apprentissage et les certifications reconnues.

\begin{definition}[Training pathways]
\eng{Continuous professional development (CPD)} : \emph{formation continue}. \eng{To enrol in a course} / \eng{to sign up for a workshop} : s'inscrire. \eng{Blended learning} : formation mixte (présentiel + distanciel). \eng{Self-paced learning} : apprentissage en autonomie. \eng{Language coaching} : coaching linguistique individuel.
\end{definition}

\begin{center}
\begin{tabularx}{\linewidth}{|X|X|}\hline
\rowcolor{poporangeL} \textbf{English} & \textbf{Français} \\ \hline
to enrol in a course / to sign up for a course & s'inscrire à une formation \\ \hline
a workshop / a seminar & un atelier / un séminaire \\ \hline
an online course / e-learning & une formation en ligne \\ \hline
blended learning & formation mixte (présentiel + distanciel) \\ \hline
a language coach & un coach linguistique \\ \hline
a proficiency test / a certification & un test de niveau / une certification \\ \hline
TOEIC / IELTS / TOEFL / Linguaskill & certifications internationales \\ \hline
CEFR (Common European Framework of Reference) & CECRL \\ \hline
\end{tabularx}
\end{center}

\begin{attention}[Faux ami : \eng{formation} $\neq$ \eng{formation}]
\eng{Formation} en anglais signifie « formation (géologique, militaire, musicale) » ou « constitution ». Pour « formation professionnelle », on dit \eng{training}, \eng{course}, \eng{programme} ou \eng{professional development}. \eng{To follow a training} se dit \eng{to take a course} ou \eng{to undergo training}. De même, \eng{un stage} (étudiant) = \eng{an internship} / \eng{a work placement} ; \eng{un stage} (formation courte) = \eng{a workshop} / \eng{a short course}.
\end{attention}

\courssub{Le champ lexical du recrutement}

Le recrutement est la première marche : il faut le vocabulaire pour postuler, passer l'entretien et négocier son contrat.

\begin{definition}[Job requirements]
\eng{Job requirements} : \emph{the skills and qualifications needed for a job} — les compétences et qualifications exigées pour un poste. On parle aussi de \eng{language requirements} (exigences linguistiques). \eng{Fluent English required} = « anglais courant exigé ».
\end{definition}

\begin{center}
\begin{tabularx}{\linewidth}{|X|X|}\hline
\rowcolor{popblueL} \textbf{English} & \textbf{Français} \\ \hline
a job interview & un entretien d'embauche \\ \hline
a CV (UK) / a résumé (US) & un curriculum vitæ \\ \hline
a cover letter & une lettre de motivation \\ \hline
an application & une candidature \\ \hline
an applicant / a candidate & un candidat \\ \hline
a recruiter / a hiring manager & un recruteur / un responsable du recrutement \\ \hline
job requirements & les conditions du poste \\ \hline
to apply for a job & postuler à un emploi \\ \hline
to meet the requirements & remplir les conditions requises \\ \hline
to shortlist & présélectionner \\ \hline
to hire / to recruit & embaucher / recruter \\ \hline
\end{tabularx}
\end{center}

\begin{exemplebox}[Le recrutement en contexte]
\eng{She meets all the job requirements.} « Elle remplit toutes les conditions du poste. »

\eng{Over two hundred applicants applied for the position.} « Plus de deux cents candidats ont postulé à ce poste. »

\eng{Send your CV and a cover letter before Friday.} « Envoyez votre CV et une lettre de motivation avant vendredi. »

\eng{The recruiter asked me about my language skills during the interview.} « Le recruteur m'a interrogé sur mes compétences linguistiques pendant l'entretien. »

\eng{He was shortlisted for the final interview.} « Il a été présélectionné pour l'entretien final. »
\end{exemplebox}

\begin{attention}[Prépositions piégeuses]
On dit \eng{to apply \textbf{for} a job} (postuler \emph{à} un emploi) mais \eng{to apply \textbf{to} a company} (s'adresser \emph{à} une entreprise). Ne dites jamais \eng{to apply a job}. De même : \eng{an interview \textbf{for} a position}, pas \eng{an interview of a position}. \eng{To be hired \textbf{by} a company}, \eng{to work \textbf{for} an employer}.
\end{attention}

\courssub{Le champ lexical des tâches quotidiennes au travail}

Une fois embauché, l'anglais devient un outil de travail quotidien. Voici les collocations \emph{verbe + nom} essentielles — apprends-les par couples indissociables.

\begin{center}
\begin{tabularx}{\linewidth}{|X|X|}\hline
\rowcolor{popgreenL} \textbf{English} & \textbf{Français} \\ \hline
to draft a report / an email & rédiger un rapport / un courriel (première version) \\ \hline
to attend a briefing / a meeting & assister à une réunion d'info / une réunion \\ \hline
to chair a meeting & présider une réunion \\ \hline
to write (the) minutes & rédiger le compte rendu / le procès-verbal \\ \hline
to make / to give a presentation & faire une présentation \\ \hline
to negotiate a contract / a deal & négocier un contrat / un accord \\ \hline
to deal with clients / complaints & s'occuper de clients / réclamations \\ \hline
to handle phone calls / enquiries & gérer les appels / demandes \\ \hline
to schedule an appointment & fixer un rendez-vous \\ \hline
to meet a deadline & respecter une échéance \\ \hline
\end{tabularx}
\end{center}

\begin{propriete}[Collocations : le bon verbe avec le bon nom]
En anglais, chaque nom appelle son verbe : on \eng{makes} a presentation mais on \eng{writes} minutes ; on \eng{chairs} a meeting mais on \eng{attends} a briefing ; on \eng{handles} calls mais on \eng{deals with} clients. Ces couples figés sont des \cle{collocations}. Le français utilise souvent « faire » ; l'anglais est précis. Apprends toujours le verbe \emph{avec} son nom.
\end{propriete}

\begin{exemplebox}[Les tâches en contexte]
\eng{Could you draft the report before the briefing?} « Pourriez-vous rédiger le rapport avant la réunion d'information ? »

\eng{Mrs Etame chaired the meeting and Mr Ndi wrote the minutes.} « Mme Etame a présidé la réunion et M. Ndi a rédigé le compte rendu. »

\eng{She deals with English-speaking clients every day.} « Elle s'occupe de clients anglophones tous les jours. »

\eng{He negotiated the contract in English.} « Il a négocié le contrat en anglais. »

\eng{We must meet the deadline for the Lagos project.} « Nous devons respecter l'échéance pour le projet de Lagos. »
\end{exemplebox}

\begin{attention}[Faux ami majeur : \eng{to assist} $\neq$ assister à]
\eng{To assist} signifie \emph{aider} : \eng{She assisted the manager} = « Elle a aidé le directeur. » Le français « assister à » (être présent) se traduit par \eng{to attend}. \eng{I assisted the meeting} est une \textbf{erreur} ; dis \eng{I attended the meeting}. De même, \eng{an assistant} est un « assistant/adjoint », pas un « participant » (\eng{an attendee} / \eng{a participant}).
\end{attention}

\courssub{Le champ lexical de la mobilité professionnelle}

L'anglais est souvent la clé qui permet de partir travailler ailleurs : mutations, expatriation, missions.

\begin{center}
\begin{tabularx}{\linewidth}{|X|X|}\hline
\rowcolor{popgoldL} \textbf{English} & \textbf{Français} \\ \hline
a posting abroad & une affectation à l'étranger \\ \hline
to be posted abroad & être muté / affecté à l'étranger \\ \hline
a transfer & une mutation, un transfert \\ \hline
expatriation & l'expatriation \\ \hline
an expatriate / an expat & un expatrié \\ \hline
a business trip & un déplacement professionnel \\ \hline
an international assignment & une mission internationale \\ \hline
to relocate & déménager (pour le travail) \\ \hline
\end{tabularx}
\end{center}

\begin{exemplebox}[La mobilité en contexte]
\eng{He was posted abroad.} « Il a été muté à l'étranger. »

\eng{She went on a business trip to Accra last month.} « Elle est partie en déplacement professionnel à Accra le mois dernier. »

\eng{After two years of expatriation in Nairobi, he returned to Yaoundé.} « Après deux ans d'expatriation à Nairobi, il est rentré à Yaoundé. »

\eng{Her transfer to the Lagos office starts in January.} « Sa mutation au bureau de Lagos commence en janvier. »

\eng{The company offered him an international assignment in London.} « L'entreprise lui a proposé une mission internationale à Londres. »
\end{exemplebox}

\begin{attention}[Constructions piégeuses : \eng{travel} vs \eng{trip}]
\eng{To go \textbf{on} a business trip} (jamais \eng{to make a business travel}). Le nom \eng{travel} est indénombrable : \eng{Travel broadens the mind} (« Les voyages forment la jeunesse »), jamais \eng{a travel}. Pour un déplacement précis : \eng{a trip} (court, aller-retour) ou \eng{a journey} (trajet, insiste sur la route).
\end{attention}

\courssub{Faux amis et pièges des francophones — Synthèse}

\begin{aretenir}[Les 5 faux amis « travail » à ne plus confondre]
\begin{center}
\begin{tabularx}{\linewidth}{|l|X|X|}\hline
\rowcolor{poppinkL} \textbf{Mot anglais} & \textbf{Sens réel} & \textbf{Piège français} \\ \hline
\eng{to assist} & aider & assister à $\rightarrow$ \eng{to attend} \\ \hline
\eng{an assistant} & un adjoint, un assistant & un participant $\rightarrow$ \eng{an attendee} \\ \hline
\eng{a resume} (US) & un CV & un résumé (texte) $\rightarrow$ \eng{a summary} \\ \hline
\eng{formation} & constitution, formation géologique & formation pro $\rightarrow$ \eng{training / a course} \\ \hline
\eng{actually} & en fait, en réalité & actuellement $\rightarrow$ \eng{currently} \\ \hline
\end{tabularx}
\end{center}
\end{aretenir}

\begin{attention}[Autres pièges fréquents]
\begin{itemize}
\item \eng{Master English} : correct mais très fort (niveau expert). Sur un CV, préfère \eng{Have a good command of English} ou \eng{Be proficient in English}.
\item \eng{Fluent \textbf{in} English} (préposition \textbf{in}), pas \eng{fluent \textbf{at/on}}.
\item \eng{Effective} (qui produit l'effet voulu) vs \eng{Efficient} (qui ne gaspille pas de ressources). Collocation : \eng{communicate effectively}.
\item \eng{Salary} (salaire mensuel/annuel) vs \eng{Wage} (salaire horaire/journalier).
\item \eng{Career} (carrière, parcours pro) vs \eng{Carriage} (wagon, carrosse) — attention à l'orthographe et à la prononciation.
\end{itemize}
\end{attention}

\courssub{Familles de mots, prononciation et orthographe}

Travailler par familles permet de multiplier le vocabulaire et de maîtriser l'accentuation tonique (stress), cruciale en anglais.

\begin{propriete}[Dérivation nominale et accentuation (Word stress)]
En anglais, le suffixe déplace souvent l'accent tonique. Règle générale : les suffixes \eng{-tion}, \eng{-sion}, \eng{-ic}, \eng{-ity}, \eng{-ial} attirent l'accent sur la syllabe qui les précède.
\end{propriete}

\begin{center}
\begin{tabular}{|l|l|l|l|l|}\hline
\rowcolor{poppurpleL} \textbf{Verbe} & \textbf{Nom (action/état)} & \textbf{Nom (personne)} & \textbf{Adjectif} & \textbf{Prononciation (approx.)} \\ \hline
\eng{apply} /əˈplaɪ/ & \eng{application} /ˌæplɪˈkeɪʃn/ & \eng{applicant} /ˈæplɪkənt/ & \eng{applicable} /ˈæplɪkəbl/ & a-pli-KEI-chen \\ \hline
\eng{recruit} /rɪˈkruːt/ & \eng{recruitment} /rɪˈkruːtmənt/ & \eng{recruiter} /rɪˈkruːtə/ & \eng{recruited} /rɪˈkruːtɪd/ & ri-KRUT-ment \\ \hline
\eng{employ} /ɪmˈplɔɪ/ & \eng{employment} /ɪmˈplɔɪmənt/ & \eng{employer/employee} /ɪmˈplɔɪə/ /ˌɛmplɔɪˈiː/ & \eng{employable} /ɪmˈplɔɪəbl/ & im-PLOI-ment \\ \hline
\eng{train} /treɪn/ & \eng{training} /ˈtreɪnɪŋ/ & \eng{trainer} /ˈtreɪnə/ & \eng{trained} /treɪnd/ & TRÉ-ning \\ \hline
\eng{assist} /əˈsɪst/ & \eng{assistance} /əˈsɪstəns/ & \eng{assistant} /əˈsɪstənt/ & \eng{assistive} /əˈsɪstɪv/ & as-SIS-tens \\ \hline
\eng{attend} /əˈtɛnd/ & \eng{attendance} /əˈtɛndəns/ & \eng{attendee} /ˌætɛnˈdiː/ & \eng{attentive} /əˈtɛntɪv/ & at-TEN-dens \\ \hline
\eng{present} (v) /prɪˈzɛnt/ & \eng{presentation} /ˌprɛznˈteɪʃn/ & \eng{presenter} /prɪˈzɛntə/ & \eng{presentable} /prɪˈzɛntəbl/ & pré-zen-TEI-chen \\ \hline
\eng{negotiate} /nɪˈɡəʊʃieɪt/ & \eng{negotiation} /nɪˌɡəʊʃiˈeɪʃn/ & \eng{negotiator} /nɪˈɡəʊʃieɪtə/ & \eng{negotiable} /nɪˈɡəʊʃiəbl/ & né-go-chi-EI-chen \\ \hline
\eng{mobilise} /ˈməʊbɪlaɪz/ & \eng{mobility} /məʊˈbɪlɪti/ & --- & \eng{mobile} /ˈməʊbaɪl/ & mo-BI-li-ti \\ \hline
\eng{expatriate} /ɛksˈpætrieɪt/ & \eng{expatriation} /ɛksˌpætriˈeɪʃn/ & \eng{expatriate/expat} /ɛksˈpætrieɪt/ & \eng{expatriate} (adj) & eks-pa-tri-EI-chen \\ \hline
\end{tabular}
\end{center}

\begin{aretenir}[Règles d'orthographe utiles]
\begin{itemize}
\item \eng{-ise} / \eng{-ize} : les deux s'écrivent (UK préfère \eng{-ise}, US \eng{-ize}) : \eng{organise/organize}, \eng{realise/realize}. Sois cohérent.
\item Double consonne avant \eng{-ing/-ed} si verbe court accentué : \eng{travel} $\rightarrow$ \eng{travelling} (UK), \eng{traveling} (US) ; \eng{control} $\rightarrow$ \eng{controlling}.
\item \eng{Programme} (UK, nom) vs \eng{Program} (US, ou informatique UK).
\item \eng{Licence} (nom, UK) vs \eng{License} (verbe, UK ; nom+verbe US). \eng{Practice} (nom) / \eng{Practise} (verbe) en UK.
\end{itemize}
\end{aretenir}

\courssub{Carte mentale : English at work}

\begin{popfigure}
\begin{center}
\begin{tikzpicture}[mindmap, grow cyclic, every node/.style={concept, align=center, font=\small},
root concept/.append style={concept color=popblue, text=white, font=\small\bfseries},
level 1/.append style={level distance=3.4cm, sibling angle=120},
level 2/.append style={level distance=2.2cm, sibling angle=45, concept color=popblueL, font=\scriptsize}]
\node[root concept]{English at work}
  child[concept color=poporange] { node {Recruitment}
    child { node {CV, cover letter} }
    child { node {job interview} }
    child { node {apply for, meet requirements} }
  }
  child[concept color=popgreen] { node {Daily tasks}
    child { node {draft reports, write minutes} }
    child { node {chair meetings, handle calls} }
    child { node {deal with, negotiate} }
  }
  child[concept color=poppurple] { node {Mobility}
    child { node {posting abroad} }
    child { node {business trips} }
    child { node {expatriation, transfer} }
  };
\end{tikzpicture}
\end{center}
\legende{Carte mentale du vocabulaire de l'anglais au travail : trois branches — recrutement, tâches quotidiennes, mobilité.}
\end{popfigure}

\courssub{Expressions pour parler de l'utilité de l'anglais}

Pour argumenter — à l'oral comme à l'écrit — sur l'importance de l'anglais, trois expressions idiomatiques sont à connaître par cœur :

\begin{aretenir}[Trois expressions clés]
\begin{itemize}
\item \eng{English is a must for this position.} « L'anglais est indispensable pour ce poste. » (\eng{a must} = une nécessité absolue)
\item \eng{A good command of English opens doors.} « Une bonne maîtrise de l'anglais ouvre des portes. »
\item \eng{English is the working language of the company.} « L'anglais est la langue de travail de l'entreprise. »
\end{itemize}
Ajoute les collocations : \eng{to be fluent in English} (parler couramment anglais), \eng{to work in an English-speaking environment} (travailler dans un environnement anglophone), \eng{to communicate effectively} (communiquer efficacement).
\end{aretenir}

\begin{savaistu}[L'anglais, passeport régional]
Au Nigeria et au Ghana, voisins du Cameroun, l'anglais est la langue de travail officielle : administration, entreprises, universités et médias fonctionnent en anglais. Pour un Camerounais francophone, maîtriser l'anglais multiplie donc les opportunités régionales : un cadre de Douala peut négocier avec Lagos, suivre une formation à Accra ou obtenir un \eng{posting} dans une filiale anglophone. Le bilinguisme français-anglais est l'un des atouts les plus recherchés sur le marché du travail d'Afrique centrale et de l'Ouest.
\end{savaistu}

\courssub{Réemploi en contexte et production guidée}

\begin{experience}[Jeu de rôle : Entretien annuel / Plan de formation]
\textbf{Contexte :} Employé (vous) et Manager. L'employé demande une formation anglais pour un projet avec le Ghana.
\textbf{Rôles :}
\begin{itemize}
\item \textbf{Employé} : Présente ton niveau (B1), tes tâches actuelles en anglais (\eng{draft reports, attend briefings}), tes difficultés (\eng{fluency, technical vocab}), ton objectif (\eng{chair a meeting in Accra}). Demande un \eng{workshop} ou \eng{coaching}.
\item \textbf{Manager} : Valide le besoin, questionne sur le \eng{ROI} (retour sur investissement), propose un \eng{blended learning} + \eng{TOEIC} en fin d'année.
\end{itemize}
\textbf{Contraintes :} Utilise au moins 5 collocations de la leçon (\eng{draft a report}, \eng{handle calls}, \eng{meet requirements}\ldots) et 2 expressions d'utilité (\eng{English is a must}\ldots).
\end{experience}

\begin{exoresolu}[Complète avec le mot ou la collocation correcte]
Énoncé — Complète chaque phrase avec : \eng{attended}, \eng{draft}, \eng{meet}, \eng{posted}, \eng{handle}, \eng{applied}, \eng{chaired}.
\begin{enumerate}
\item She \ldots\ all the job requirements, so the recruiter called her for an interview.
\item Could you \ldots\ the report before Friday's briefing?
\item I \ldots\ a meeting with our Nigerian partners yesterday.
\item He was \ldots\ abroad, in the Accra office, for two years.
\item My job is to \ldots\ phone calls from English-speaking clients.
\item Last month, I \ldots\ for a position in a multinational company.
\item The director \ldots\ the meeting because the CEO was absent.
\end{enumerate}
\tcblower
Solution détaillée :
\begin{enumerate}
\item \eng{meets} — collocation \eng{to meet the requirements} ; présent simple, 3e personne : \emph{meets}.
\item \eng{draft} — collocation \eng{to draft a report} ; après \eng{could}, base verbale.
\item \eng{attended} — faux ami : « assister à » = \eng{to attend}, jamais \eng{to assist} ; ici au prétérit (\eng{yesterday}).
\item \eng{posted} — \eng{to be posted abroad} = être muté à l'étranger ; passif au prétérit.
\item \eng{handle} — collocation \eng{to handle phone calls} ; après \eng{to}, base verbale.
\item \eng{applied} — \eng{to apply for a position} ; prétérit.
\item \eng{chaired} — \eng{to chair a meeting} = présider ; prétérit.
\end{enumerate}
\end{exoresolu}

\begin{exercice}[À toi de jouer]
\begin{enumerate}
\item Traduis en anglais :
      \begin{enumerate}
      \item Une bonne maîtrise de l'anglais ouvre des portes.
      \item Elle a présidé la réunion et il a rédigé le compte rendu.
      \item J'ai postulé à un emploi au Ghana.
      \item Nous devons respecter l'échéance de vendredi.
      \item Il est en déplacement professionnel à Nairobi cette semaine.
      \end{enumerate}
\item Corrige l'erreur dans chaque phrase :
      \begin{enumerate}
      \item I assisted the briefing last Monday.
      \item She has a good command in English.
      \item He made a business travel to Lagos.
      \item The recruiter asked me to send my formation.
      \item We need to make the minutes of the meeting.
      \end{enumerate}
\item Production écrite (80 mots) : Tu es responsable RH à Douala. Écris un courriel formel à ton Directeur Général pour justifier le budget d'une formation anglais \eng{blended learning} pour 10 cadres. Utilise : \eng{language requirements}, \eng{business trips}, \eng{posting abroad}, \eng{proficiency test}, \eng{ROI}.
\end{enumerate}
\end{exercice}

\begin{corrige}[Exercice — À toi de jouer]
\begin{enumerate}
\item
      \begin{enumerate}
      \item \eng{A good command of English opens doors.}
      \item \eng{She chaired the meeting and he wrote the minutes.}
      \item \eng{I applied for a job in Ghana.}
      \item \eng{We must meet Friday's deadline.} / \eng{We have to meet the deadline for Friday.}
      \item \eng{He is on a business trip to Nairobi this week.}
      \end{enumerate}
\item
      \begin{enumerate}
      \item \eng{I attended the briefing last Monday.} (\eng{to assist} = aider ; « assister à » = \eng{to attend}.)
      \item \eng{She has a good command \textbf{of} English.} (Préposition \textbf{of} après \eng{command}.)
      \item \eng{He went on a business trip to Lagos.} (\eng{travel} indénombrable ; \eng{a trip} pour un déplacement.)
      \item \eng{The recruiter asked me to send my CV / training certificates.} (\eng{formation} = \eng{training / a course} ; \eng{formation} seul ne s'emploie pas ainsi.)
      \item \eng{We need to write the minutes of the meeting.} (\eng{write minutes}, pas \eng{make}.)
      \end{enumerate}
\item Production libre. Vérifier : structure formelle (\eng{Dear Sir/Madam}, \eng{Subject:}, \eng{Yours faithfully}), vocabulaire imposé employé correctement, arguments clairs (besoins clients anglophones, mobilité, certification).
\end{enumerate}
\end{corrige}

Tu sais désormais évaluer tes besoins, nommer les étapes du recrutement, décrire les tâches quotidiennes en anglais, parler de mobilité internationale, éviter les faux amis majeurs et manipuler les familles de mots avec leur prononciation. Dans le prochain module, nous consoliderons ces acquis par une révision générale et des sujets d'examen type Baccalauréat.

\coursec{Faux amis et pièges des francophones}

Dans la section précédente, nous avons découvert comment l'anglais fonctionne dans la vie professionnelle : rédiger un CV, passer un entretien, décrire ses compétences. Mais un danger guette le candidat francophone : les mots qui \emph{ressemblent} au français mais ne \emph{signifient} pas la même chose en anglais. Imaginez un candidat de Douala qui écrit sur son CV : \eng{“I followed a formation last year.”} Le recruteur anglophone fronce les sourcils : en anglais, \eng{formation} désigne la disposition de joueurs sur un terrain de football ou la création d'une entreprise, jamais un cours de formation ! Cette section vous apprend à repérer et à éviter ces pièges.

\courssub{Qu'est-ce qu'un faux ami ?}

\begin{definition}[False friend — faux ami]
A \cle{false friend} (faux ami) is a word that looks or sounds similar in two languages but has a different meaning in each. French and English share thousands of words of Latin origin, which makes false friends especially dangerous for francophone learners. Exemple : \eng{actually} ne veut pas dire « actuellement » mais « en fait, en réalité ».
\end{definition}

\begin{attention}[Le piège fondamental]
Parce que l'anglais et le français partagent une même racine latine, environ un mot anglais sur trois ressemble à un mot français. La plupart sont de vrais amis (\eng{nation} = nation, \eng{important} = important), mais les faux amis sont assez nombreux pour ruiner la crédibilité d'un CV ou d'une lettre de motivation. Règle d'or : \textbf{la ressemblance ne garantit jamais le sens}.
\end{attention}

\courssub{Observation : un paragraphe truffé d'erreurs}

Lisez ce court texte écrit par un élève de Terminale. Chaque phrase contient un faux ami ou un calque du français. Saurez-vous les repérer avant de lire la correction ?

\begin{textebox}[A cover letter full of traps — lettre de motivation piégée]
Dear Sir,

I am writing to demand information about your company. I have a good level in English and I would like to ameliorate it. Last year, I followed a formation in computer science and I did a stage in a bank in Douala. I assisted to many meetings with the manager. I am actually available for an interview at any time. I am very motivated and I can assist your team with pleasure.

Yours faithfully,

Jean-Paul N.
\end{textebox}

Comptez les erreurs : il y en a \textbf{sept}. Les voici, corrigées une à une dans les sous-sections suivantes.

\courssub{Les six faux amis essentiels du monde du travail}

\textbf{1. \eng{formation} (fr.) $\neq$ \eng{formation} (ang.)}\par
En français, une « formation » est un apprentissage. En anglais, \eng{formation} signifie « formation » au sens de \emph{disposition, structure} (a rock formation = une formation rocheuse ; the formation of a company = la création d'une entreprise). Pour parler d'apprentissage, on dit \cle{training} ou \cle{training course}.

\eng{I attended a training course in accounting last year.} « J'ai suivi une formation en comptabilité l'année dernière. »

\textbf{2. « assister à » $\neq$ \eng{to assist}}\par
En français, « assister à une réunion » signifie y être présent. En anglais, \eng{to assist} signifie \emph{aider} (to assist a customer = aider un client). Pour « assister à », on dit \cle{to attend}.

\eng{She attended the meeting in Yaoundé.} « Elle a assisté à la réunion à Yaoundé. »
\eng{He assisted an old lady with her luggage.} « Il a aidé une dame âgée avec ses bagages. »

\textbf{3. « une formation professionnelle » $\rightarrow$ \eng{vocational training}}\par
Ne dites jamais \eng{professional formation}. L'expression correcte est \cle{vocational training} (formation professionnelle) ; on parle aussi de \cle{vocational school} (lycée technique / école de formation professionnelle).

\eng{Many young Cameroonians choose vocational training after the BEPC.} « Beaucoup de jeunes Camerounais choisissent une formation professionnelle après le BEPC. »

\textbf{4. « actuel » $\neq$ \eng{actual}}\par
\eng{Actual} signifie \emph{réel, véritable}. Pour « actuel », on dit \cle{current} ; pour « actuellement », \cle{currently} ou \eng{at present}.

\eng{The actual cost is higher than we expected.} « Le coût réel est plus élevé que prévu. »
\eng{My current job is in Buea.} « Mon emploi actuel est à Buéa. »
\eng{I am currently looking for a job.} « Je cherche actuellement un emploi. »

\textbf{5. « demander » $\neq$ \eng{to demand}}\par
\eng{To demand} signifie \emph{exiger} (souvent avec agressivité). Pour « demander », on dit \cle{to ask} ou \cle{to ask for} (demander quelque chose).

\eng{I am writing to ask for information about the job.} « Je vous écris pour demander des informations sur le poste. »
\eng{The workers demanded higher salaries.} « Les ouvriers ont exigé des salaires plus élevés. »

\textbf{6. « un stage / un stagiaire » $\rightarrow$ \eng{an internship / an intern ou a trainee}}\par
Le mot anglais \eng{stage} signifie « scène » (théâtre) ou « étape ». Un stage professionnel se dit \cle{internship} (surtout en anglais américain) ou \cle{work placement} (anglais britannique) ; le stagiaire est \cle{an intern} ou \cle{a trainee}.

\eng{He is doing an internship at a bank.} « Il fait un stage dans une banque. »
\eng{The trainees will start on Monday.} « Les stagiaires commenceront lundi. »

\begin{popfigure}
\begin{center}
\begin{tikzpicture}[font=\small]
\fill[popblueL] (0,0) rectangle (3.1,0.8);
\node[align=center] at (1.55,0.4) {\textbf{Français}};
\fill[poporangeL] (3.2,0) rectangle (6.3,0.8);
\node[align=center] at (4.75,0.4) {\textbf{Faux ami (à éviter)}};
\fill[popgreenL] (6.4,0) rectangle (9.5,0.8);
\node[align=center] at (7.95,0.4) {\textbf{Bon anglais}};
\foreach \y/\fra/\faux/\bon in {
 -0.7/une formation/formation/training (course),
 -1.5/assister à/to assist/to attend,
 -2.3/formation professionnelle/professional formation/vocational training,
 -3.1/actuel(lement)/actual(ly)/current(ly),
 -3.9/demander/to demand/to ask (for),
 -4.7/un stage ou un stagiaire/a stage/an internship or an intern}{
 \draw[popline] (0,\y-0.4) -- (9.5,\y-0.4);
 \node[align=center] at (1.55,\y) {\fra};
 \node[align=center, text=poporange] at (4.75,\y) {\emph{\faux}};
 \node[align=center, text=popgreen] at (7.95,\y) {\textbf{\bon}};
}
\draw[popline] (0,0.8) rectangle (9.5,-5.1);
\draw[popline] (3.15,0.8) -- (3.15,-5.1);
\draw[popline] (6.35,0.8) -- (6.35,-5.1);
\end{tikzpicture}
\end{center}
\legende{Six faux amis à connaître par cœur pour le monde du travail}
\end{popfigure}

\courssub{Les calques : des traductions mot à mot à bannir}

Au-delà des faux amis, les francophones traduisent souvent leurs expressions mot à mot. Voici deux calques très fréquents dans les lettres de motivation :

\begin{propriete}[Calques fréquents à corriger]
\begin{itemize}
\item \eng{I have a good level in English.} (calque de « j'ai un bon niveau ») $\rightarrow$ dites : \eng{I have a good command of English.} ou \eng{I speak English fluently.}
\item \eng{I want to ameliorate my English.} (calque de « améliorer ») $\rightarrow$ dites : \eng{I want to improve my English.} Le verbe \eng{to ameliorate} existe mais est rare et soutenu ; \eng{to improve} est le mot naturel.
\item \eng{I am agree.} (calque de « je suis d'accord ») $\rightarrow$ dites : \eng{I agree.}
\item \eng{I have twenty years.} $\rightarrow$ dites : \eng{I am twenty (years old).}
\end{itemize}
\end{propriete}

\begin{exemplebox}[La lettre corrigée]
Voici la lettre de Jean-Paul après correction :
\eng{Dear Sir, I am writing to ask for information about your company. I have a good command of English and I would like to improve it. Last year, I attended a training course in computer science and I did an internship in a bank in Douala. I attended many meetings with the manager. I am currently available for an interview at any time. I am very motivated and I would be happy to assist your team. Yours faithfully, Jean-Paul N.}
\end{exemplebox}

\begin{methode}[Comment débusquer un faux ami en quatre étapes]
\begin{enumerate}
\item \textbf{Méfiance} : si le mot anglais ressemble exactement au mot français, allumez un signal d'alarme.
\item \textbf{Contexte} : lisez la phrase entière. Le sens supposé est-il logique ? (\eng{The actual situation} = la situation réelle, pas actuelle.)
\item \textbf{Vérification} : consultez un dictionnaire bilingue ou votre liste personnelle avant d'écrire le mot dans une production.
\item \textbf{Mémorisation} : notez la paire (faux ami / bon mot) avec un exemple dans votre carnet de vocabulaire.
\end{enumerate}
\end{methode}

\begin{exoresolu}[Corrige les erreurs]
Énoncé : corrigez les faux amis et calques dans les phrases suivantes.
\begin{enumerate}
\item \eng{I followed a formation in marketing last year.}
\item \eng{The actual unemployment rate worries the government.}
\item \eng{She assisted to a conference in Bamenda.}
\item \eng{I demand you to send me the application form.}
\item \eng{My brother found a stage in a travel agency.}
\end{enumerate}
\tcblower
Solution détaillée :
\begin{enumerate}
\item \eng{I attended a training course in marketing last year.} — \eng{formation} est un faux ami ; pour un apprentissage, on utilise \eng{training course} avec le verbe \eng{to attend}.
\item \eng{The current unemployment rate worries the government.} — \eng{actual} signifie « réel » ; « actuel » se traduit \eng{current}.
\item \eng{She attended a conference in Bamenda.} — \eng{to assist} signifie « aider » ; « assister à » se dit \eng{to attend} (sans préposition).
\item \eng{I would like you to send me the application form.} ou \eng{I am writing to ask for the application form.} — \eng{to demand} signifie « exiger » et se construit avec un nom, jamais avec \eng{somebody to do}.
\item \eng{My brother found an internship in a travel agency.} — \eng{stage} signifie « scène » ou « étape » ; un stage professionnel est \eng{an internship}.
\end{enumerate}
\end{exoresolu}

\begin{attention}[Le calque de la lettre de motivation]
La phrase \eng{“I followed a formation last year”} est un double calque : le verbe \eng{follow} (suivre quelqu'un ou quelque chose physiquement) et le nom \eng{formation}. La phrase correcte est \eng{“I attended a training course last year.”} Retenez le couple verbe + nom : \cle{to attend a training course}, \cle{to do an internship}, \cle{to improve one's English}.
\end{attention}

\begin{center}
\begin{tabularx}{\linewidth}{|X|X|X|}
\hline
\rowcolor{popblueL} \textbf{Français} & \textbf{Piège (à éviter)} & \textbf{Bon anglais + exemple} \\
\hline
une formation & \emph{formation} & \eng{training course} — \eng{a computer training course} \\
\hline
assister à & \emph{to assist} (= aider) & \eng{to attend} — \eng{to attend a meeting} \\
\hline
formation professionnelle & \emph{professional formation} & \eng{vocational training} \\
\hline
actuel / actuellement & \emph{actual / actually} (= réel / en fait) & \eng{current / currently} — \eng{my current job} \\
\hline
demander & \emph{to demand} (= exiger) & \eng{to ask (for)} — \eng{to ask for information} \\
\hline
un stage / un stagiaire & \emph{a stage} (= scène, étape) & \eng{an internship / an intern, a trainee} \\
\hline
j'ai un bon niveau & \emph{I have a good level} & \eng{I have a good command of English} \\
\hline
améliorer mon anglais & \emph{to ameliorate my English} & \eng{to improve my English} \\
\hline
\end{tabularx}
\end{center}

\begin{savaistu}[Les faux amis coûtent des points au Bac]
À l'épreuve d'anglais du Baccalauréat, l'usage répété de faux amis fait perdre des points dans le critère d'\emph{accuracy} (correction de la langue). Les correcteurs repèrent immédiatement des phrases comme \eng{I am actually unemployed} pour « je suis actuellement au chômage ». La meilleure stratégie de révision : tenez tout au long de l'année une \textbf{liste personnelle de faux amis}, avec pour chacun un exemple correct et un exemple incorrect. Relisez-la la veille de l'examen !
\end{savaistu}

\begin{exercice}[À toi de jouer]
Traduis en anglais correct, en évitant les faux amis :
\begin{enumerate}
\item Je suis actuellement en formation professionnelle à Douala.
\item Elle a assisté à une réunion importante hier.
\item Il fait un stage dans une entreprise de cacao.
\item Je vous écris pour demander un emploi.
\item Le coût réel de la formation est élevé.
\end{enumerate}
\end{exercice}

\begin{corrige}[Exercice — Traduction]
\begin{enumerate}
\item \eng{I am currently doing vocational training in Douala.}
\item \eng{She attended an important meeting yesterday.}
\item \eng{He is doing an internship in a cocoa company.}
\item \eng{I am writing to ask for a job.} (ou \eng{to apply for a job})
\item \eng{The actual cost of the training course is high.} — ici \eng{actual} est correct, car « réel » !
\end{enumerate}
\end{corrige}

\begin{aretenir}[L'essentiel de la section]
Un faux ami ressemble au français mais signifie autre chose. Les six pièges majeurs du monde professionnel : \eng{formation} $\rightarrow$ \cle{training course} ; \eng{to assist} $\rightarrow$ \cle{to attend} ; \eng{actual} $\rightarrow$ \cle{current} ; \eng{to demand} $\rightarrow$ \cle{to ask for} ; \eng{stage} $\rightarrow$ \cle{internship} ; \eng{professional formation} $\rightarrow$ \cle{vocational training}. En cas de doute, vérifie toujours le sens dans le contexte ou dans un dictionnaire avant d'écrire.
\end{aretenir}

\coursec{Familles de mots, prononciation et orthographe}

Dans la section précédente, nous avons traqué les faux amis et les calques du français qui guettent le candidat au Bac. Nous allons maintenant adopter une stratégie plus positive : au lieu d'apprendre les mots un par un, nous allons les apprendre \textbf{par familles}. Connaître \eng{to employ}, c'est pouvoir deviner et construire \eng{employer}, \eng{employee}, \eng{employment}, \eng{unemployment} et \eng{employable}. C'est multiplier par cinq ou six votre vocabulaire actif sans effort supplémentaire — à condition de maîtriser les suffixes, la prononciation et l'orthographe de chaque dérivé.

\courssub{Observer avant de mémoriser : un texte pour découvrir les familles}

\begin{textebox}[Building a career in Yaoundé]
When Amina finished her A-Levels in Yaoundé, she was worried about \textbf{unemployment}. Many of her friends were \textbf{unemployed}, and finding a job seemed difficult. She decided to improve her \textbf{employability}: she enrolled in a computer \textbf{training} centre where a patient \textbf{trainer} taught her office software. As a \textbf{trainee}, she worked hard and obtained a \textbf{qualification} that made her \textbf{qualified} for administrative posts.

Her first \textbf{employer}, an import-export company in Douala, appreciated her \textbf{fluency} in English. She spoke so \textbf{fluently} that clients thought she had studied abroad. In fact, her \textbf{proficiency} came from years of regular practice. After two years as an \textbf{employee}, she was promoted. Today she says: “A good \textbf{qualification} opens the door, but language skills keep you inside.”
\end{textebox}

\textbf{Glossaire :} \emph{unemployment} (n.) chômage ; \emph{employability} (n.) employabilité ; \emph{trainee} (n.) stagiaire ; \emph{fluency} (n.) aisance, fluidité ; \emph{proficiency} (n.) maîtrise, compétence ; \emph{to be promoted} (v.) être promu.

\textbf{Questions de compréhension :}
\begin{enumerate}
\item What did Amina do to fight unemployment?
\item Who teaches in a training centre, and who learns?
\item According to Amina, what is the difference between a qualification and language skills?
\end{enumerate}

\courssub{La famille de EMPLOY}

\begin{definition}[Famille de EMPLOY]
\eng{To employ} (verbe) signifie « employer, embaucher ». De cette racine, l'anglais construit toute une famille :
\begin{itemize}
\item \cle{employer} (n.) : l'employeur, celui qui embauche ;
\item \cle{employee} (n.) : le salarié, celui qui est embauché ;
\item \cle{employment} (n.) : l'emploi, le travail ;
\item \cle{unemployment} (n.) : le chômage ;
\item \cle{employable} (adj.) : employable, apte à être embauché.
\end{itemize}
\end{definition}

\begin{popfigure}
\begin{tikzpicture}[grow=down, level distance=1.4cm,
  every node/.style={draw=popblue, rounded corners, fill=popblueL, align=center, font=\small},
  level 1/.style={sibling distance=4.2cm},
  level 2/.style={sibling distance=3.4cm},
  edge from parent/.style={draw=popdark, -{Stealth}}]
\node[fill=poporangeL, draw=poporange, font=\bfseries, align=center]{employ\\(racine)}
  child { node[align=center]{employer\\(employeur)} }
  child { node[align=center]{employee\\(salarié)} }
  child { node[align=center]{employment\\(emploi)}
    child { node[align=center]{unemployment\\(chômage)} } }
  child { node[align=center]{employable\\(employable)} };
\end{tikzpicture}
\legende{Arbre de la famille de \eng{employ} : une racine, cinq dérivés.}
\end{popfigure}

\begin{exemplebox}[La famille EMPLOY en contexte]
\eng{The employer interviewed the employee about his working conditions.} « L'employeur a passé un entretien avec le salarié au sujet de ses conditions de travail. »

\eng{Unemployment among young graduates is a serious problem.} « Le chômage chez les jeunes diplômés est un problème grave. »

\eng{Computer skills make you more employable.} « Les compétences informatiques vous rendent plus employable. »
\end{exemplebox}

\courssub{Les familles de TRAIN, QUALIFY, FLUENT et PROFICIENT}

\begin{center}
\begin{tabularx}{\linewidth}{|l|l|X|X|}
\hline
\rowcolor{popblueL} Racine & Nature & Mot dérivé & Traduction \\
\hline
train & verbe & to train & former, entraîner \\
\hline
train & nom (-er) & trainer & le formateur \\
\hline
train & nom (-ee) & trainee & le stagiaire, le formé \\
\hline
train & nom (-ing) & training & la formation \\
\hline
qualify & verbe & to qualify & qualifier, obtenir un diplôme \\
\hline
qualify & nom & qualification & le diplôme, la qualification \\
\hline
qualify & adjectif & qualified / unqualified & qualifié / non qualifié \\
\hline
qualify & adjectif & qualifying & qualificatif (examen) \\
\hline
fluent & adjectif & fluent & courant (parler couramment) \\
\hline
fluent & adverbe & fluently & couramment \\
\hline
fluent & nom & fluency & l'aisance, la fluidité \\
\hline
proficient & adjectif & proficient & compétent, maîtrisant \\
\hline
proficient & nom & proficiency & la maîtrise, la compétence \\
\hline
\end{tabularx}
\end{center}

\begin{exemplebox}[Exemples traduits]
\eng{She is a trainee in a law firm in Buea.} « Elle est stagiaire dans un cabinet d'avocats à Buea. »

\eng{The trainer gave the trainees a practical exercise.} « Le formateur a donné aux stagiaires un exercice pratique. »

\eng{He is fluent in English and speaks it fluently.} « Il parle couramment anglais et le parle avec aisance. »

\eng{A certificate of proficiency in English is required for this job.} « Un certificat de maîtrise de l'anglais est exigé pour ce poste. »

\eng{Only qualified candidates will be invited for the interview.} « Seuls les candidats qualifiés seront convoqués à l'entretien. »
\end{exemplebox}

\courssub{La règle des suffixes -er et -ee}

\begin{propriete}[Suffixes -er et -ee : celui qui agit / celui qui subit]
En anglais, deux suffixes forment des noms de personnes à partir d'un verbe :
\begin{itemize}
\item \cle{-er} (ou \emph{-or}) désigne \textbf{celui qui fait l'action} : \eng{the trainer teaches} (le formateur enseigne), \eng{the employer pays} (l'employeur paie) ;
\item \cle{-ee} désigne \textbf{celui qui reçoit ou subit l'action} : \eng{the trainee learns} (le stagiaire apprend), \eng{the employee works} (le salarié travaille).
\end{itemize}
Autres paires utiles : \eng{interviewer / interviewee} (l'intervieweur / la personne interviewée), \eng{examiner / examinee} (l'examinateur / le candidat), \eng{trainer / trainee}, \eng{employer / employee}.
\end{propriete}

\begin{attention}[Prononciation de employee]
Ne prononcez jamais \eng{employee} « à la française » (\emph{an-plo-ia-yé}). En anglais, l'accent tonique tombe sur la \textbf{dernière syllabe} : « èm-ploï-\textbf{Ï} ». De même : \eng{trainee} « trèï-\textbf{NI} », \eng{interviewee} « in-teur-viou-\textbf{I} ». À l'inverse, \eng{employer} s'accentue sur la deuxième syllabe : « èm-\textbf{PLO}-yeur ».
\end{attention}

\courssub{Prononciation des mots clés de la leçon}

\begin{center}
\begin{tabularx}{\linewidth}{|l|X|l|}
\hline
\rowcolor{popblueL} Mot & Prononciation (lettres françaises) & Traduction \\
\hline
employer & « èm-plo-yeur » & l'employeur \\
\hline
employee & « èm-ploï-ï » (accent final) & le salarié \\
\hline
unemployment & « an-èm-ploï-mente » & le chômage \\
\hline
fluency & « flou-en-ssi » & l'aisance \\
\hline
proficiency & « pro-fi-chèn-ssi » & la maîtrise \\
\hline
certificate & « seur-ti-fi-kète » & le certificat \\
\hline
qualification & « kouo-li-fi-ké-chène » & le diplôme \\
\hline
trainee & « trèï-ni » (accent final) & le stagiaire \\
\hline
\end{tabularx}
\end{center}

\courssub{Orthographe : les pièges à éviter}

\begin{propriete}[Doublement de la consonne]
Devant un suffixe commençant par une voyelle (\emph{-er}, \emph{-ing}, \emph{-ed}), la consonne finale d'un verbe court (une syllabe accentuée, schéma consonne-voyelle-consonne) \textbf{double} : \eng{begin} $\rightarrow$ \eng{beginner} (débutant), \eng{beginning} ; \eng{run} $\rightarrow$ \eng{runner}. D'où \eng{a beginner's level} « un niveau débutant ». En revanche, pas de doublement si la voyelle est longue ou double : \eng{train} $\rightarrow$ \eng{trainer} (et non \emph{trainner}).
\end{propriete}

\begin{attention}[Mots souvent mal orthographiés]
\begin{itemize}
\item \eng{necessary} : un seul \textbf{c}, deux \textbf{s} (« né-sé-sa-ri ») — jamais \emph{neccessary} ;
\item \eng{professional} : un seul \textbf{f}, deux \textbf{s} — jamais \emph{professionnal} ;
\item \eng{assessment} : deux \textbf{s} au début, deux \textbf{s} au milieu (\emph{a-ss-e-ss-ment}) ;
\item \eng{fluency} : pas de \emph{o} après \emph{u} — jamais \emph{fluensy} ;
\item \eng{beginner} : deux \textbf{n} — jamais \emph{beginer}.
\end{itemize}
\end{attention}

\begin{methode}[Apprendre le vocabulaire par familles]
\begin{enumerate}
\item Choisis un verbe ou un adjectif de base (ex. \eng{employ}).
\item Construis tous ses dérivés à l'aide des suffixes : \emph{-er}, \emph{-ee}, \emph{-ment}, \emph{-able}, \emph{-ity}, \emph{-ly}.
\item Note chaque mot dans une phrase complète, pas isolément.
\item Vérifie la prononciation et l'accent tonique de chaque dérivé (ils changent souvent !).
\item Réemploie la famille entière dans un court paragraphe écrit.
\end{enumerate}
\end{methode}

\begin{exoresolu}[Compléter avec le bon dérivé]
Complétez chaque phrase avec un mot de la famille indiquée entre parenthèses.

1. The \ldots\ldots offered the job to the best candidate. (employ)
2. She attended a six-month \ldots\ldots course in accounting. (train)
3. His \ldots\ldots in English helped him get the job. (fluent)
4. An \ldots\ldots person cannot find work. (employ — négatif)
5. You need a language \ldots\ldots to apply for this post. (qualify)
\tcblower
\textbf{Solutions :}
1. \eng{employer} — c'est celui qui embauche (suffixe \emph{-er}, celui qui agit).
2. \eng{training} — nom d'action en \emph{-ing} : « une formation de six mois en comptabilité ».
3. \eng{fluency} — il faut un nom après le possessif \eng{his} ; \eng{fluent} est un adjectif.
4. \eng{unemployed} — adjectif négatif avec le préfixe \emph{un-} : « une personne au chômage ».
5. \eng{qualification} — nom dérivé du verbe \eng{to qualify} : « un diplôme de langue ».
\end{exoresolu}

\begin{savaistu}[Unemployed ou jobless ?]
\eng{Unemployed} et \eng{jobless} sont synonymes (« au chômage »), mais ils n'ont pas le même registre. \eng{Jobless} est plus journalistique et imagé ; \eng{unemployed} est le terme \textbf{neutre et formel} attendu dans une composition du Bac, tout comme le nom \eng{unemployment} (et non \emph{joblessness}) dans une phrase comme : \eng{The government must reduce unemployment.} « Le gouvernement doit réduire le chômage. »
\end{savaistu}

\begin{exercice}[À vous de jouer]
1. Construisez la famille complète de \eng{to develop} (verbe, deux noms, un adjectif).
2. Traduisez : « Le formateur a félicité la stagiaire pour son aisance en anglais. »
3. Corrigez les fautes : \emph{The employe was very proffesional and his fluensy was impressive.}
\end{exercice}

\begin{corrige}[Exercice : familles et orthographe]
1. \eng{to develop} (v.) $\rightarrow$ \eng{developer} (n., développeur), \eng{development} (n., développement), \eng{developed / developing} (adj., développé / en développement).
2. \eng{The trainer congratulated the trainee on her fluency in English.}
3. \eng{The \textbf{employee} was very \textbf{professional} and his \textbf{fluency} was impressive.} (un seul \emph{f} dans \emph{professional}, pas de \emph{s} final à \emph{fluency}, accent tonique final dans \emph{employee}).
\end{corrige}

\coursec{Réemploi en contexte et production guidée}

Dans la section précédente, nous avons enrichi notre vocabulaire grâce aux familles de mots (\emph{employ}, \emph{employer}, \emph{employee}, \emph{employment}…) et nous avons soigné la prononciation et l'orthographe. Il est temps maintenant de \cle{réemployer} tout ce lexique dans des tâches réelles : comprendre un texte, dialoguer avec un employeur et rédiger un paragraphe personnel. C'est exactement le type de production attendu au baccalauréat.

\courssub{Lecture : « Planning my English future »}

Lis le texte suivant. Un élève de Terminale explique comment il \cle{anticipe} ses besoins en anglais pour son futur métier.

\begin{textebox}[Planning my English future]
My name is Arnaud and I am a Terminale student in Yaoundé. My dream is to become a civil engineer and to work on road projects across Africa. I soon realised that English would be essential for my career: most technical documents, software manuals and international contracts are written in English.

Currently, my level is intermediate. I can understand simple conversations and I can write short emails, but my speaking is still weak and I lack technical vocabulary. That is why I have decided to anticipate my language needs now, before I start looking for a job.

First, I signed up for an evening course at a language centre in order to improve my fluency. Moreover, I watch engineering videos in English every week so that I can learn the technical vocabulary of my field. I also practise speaking with my cousin who lives in Buea, so that I can gain confidence. As a result, my English is improving month after month.

When I finish my studies, I would like to take an internationally recognised exam, such as the IELTS, because many companies ask for a certificate. I am convinced that a good command of English will open doors for me, not only in Cameroon but also abroad. Anticipating my second language needs today means building my professional future.
\end{textebox}

\begin{definition}[Glossaire du texte]
\begin{itemize}
\item \eng{civil engineer} (n.) : ingénieur civil
\item \eng{to realise} (v.) : se rendre compte
\item \eng{essential} (adj.) : essentiel, indispensable
\item \eng{weak} (adj.) : faible
\item \eng{to sign up for} (v.) : s'inscrire à
\item \eng{fluency} (n.) : aisance (à l'oral)
\item \eng{to gain confidence} (exp.) : prendre confiance
\item \eng{a good command of English} (exp.) : une bonne maîtrise de l'anglais
\item \eng{to open doors} (exp.) : ouvrir des portes, créer des opportunités
\end{itemize}
\end{definition}

\begin{exoresolu}[Compréhension du texte]
\textbf{A. Vrai ou faux ?} 1. Arnaud wants to become a doctor. 2. His speaking is weak. 3. He has signed up for an evening course. 4. He never practises speaking.
\textbf{B. QCM.} Arnaud watches engineering videos (a) to entertain himself (b) to learn technical vocabulary (c) to teach his cousin.
\textbf{C. Réponse courte.} Why does Arnaud want to take the IELTS?
\tcblower
\textbf{A.} 1. \textbf{False} — he wants to become a civil engineer. 2. \textbf{True} — “my speaking is still weak”. 3. \textbf{True} — “I signed up for an evening course”. 4. \textbf{False} — he practises with his cousin in Buea.
\textbf{B.} (b) — “so that I can learn the technical vocabulary of my field”.
\textbf{C.} Because many companies ask for an internationally recognised certificate.
\end{exoresolu}

\begin{exercice}[À toi de jouer]
Réponds par des phrases complètes : 1. What is Arnaud's current level of English? 2. Name two actions he takes to improve his English. 3. What does “anticipating my second language needs” mean in your own words?
\end{exercice}

\begin{corrige}[Exercice — Compréhension]
1. His current level is intermediate: he can understand simple conversations and write short emails. 2. He attends an evening course and watches engineering videos in English; he also practises speaking with his cousin. 3. It means preparing today the language skills you will need for your future job.
\end{corrige}

\courssub{Dialogue guidé : convaincre un employeur}

Imaginons une situation professionnelle : tu demandes à ton employeur de financer une formation d'anglais. Observe d'abord les formules de politesse et d'argumentation.

\begin{propriete}[Formules pour convaincre]
\begin{itemize}
\item \eng{I would like to improve my English.} — J'aimerais améliorer mon anglais. (\emph{would like to} + base verbale : plus poli que \emph{want})
\item \eng{It would help me to communicate with foreign clients.} — Cela m'aiderait à communiquer avec les clients étrangers.
\item \eng{This training would benefit the company because we could negotiate international contracts.} — Cette formation profiterait à l'entreprise car nous pourrions négocier des contrats internationaux.
\end{itemize}
Remarque l'emploi du conditionnel \cle{would} : il adoucit la demande, comme le conditionnel français (“j'aimerais”, “cela m'aiderait”).
\end{propriete}

\begin{experience}[Jeu de rôle : at the manager's office]
Par groupes de deux : l'élève A est l'employé(e), l'élève B est le directeur ou la directrice d'une entreprise de Douala. A demande le financement d'une formation d'anglais en utilisant au moins trois formules du cours (\emph{I would like to…}, \emph{It would help me to…}, \emph{This training would benefit the company because…}). B pose des questions (\emph{How much does it cost? How long is the course?}) puis accepte ou refuse. Changez ensuite les rôles.
\end{experience}

\courssub{Les connecteurs utiles}

Pour rédiger un paragraphe cohérent, il faut enchaîner les idées. Voici les connecteurs de cette leçon.

\begin{center}
\begin{tabularx}{\linewidth}{|X|X|X|}
\hline
\rowcolor{popblueL} English & Français & Emploi \\
\hline
first, then, finally & d'abord, ensuite, enfin & ordre des étapes \\
\hline
moreover & de plus, en outre & ajout d'un argument \\
\hline
therefore / as a result & donc / par conséquent & conséquence \\
\hline
in order to + BV & afin de + infinitif & but (même sujet) \\
\hline
so that + proposition & afin que + phrase & but (sujet exprimé) \\
\hline
\end{tabularx}
\end{center}

\begin{exemplebox}[Exemples traduits]
\eng{I am learning English in order to work abroad.} « J'apprends l'anglais afin de travailler à l'étranger. »

\eng{She attends evening classes so that she can get promoted.} « Elle suit des cours du soir afin d'obtenir une promotion. »

\eng{My writing is weak; therefore, I practise every day.} « Mon expression écrite est faible ; c'est pourquoi je m'entraîne chaque jour. »

\eng{Moreover, English is required in most job advertisements.} « De plus, l'anglais est exigé dans la plupart des offres d'emploi. »
\end{exemplebox}

\begin{attention}[« in order to » et « so that » : ne pas confondre]
\eng{in order to} est suivi de la \cle{base verbale} : \eng{in order to improve my English}. Ne dis jamais \emph{in order for to improve} ! Si le sujet change, utilise \eng{in order for + nom/pronom + to} : \eng{My parents saved money in order for me to study abroad.}

\eng{so that} est suivi d'une \cle{proposition complète} (sujet + verbe, souvent avec \emph{can/could/will}) : \eng{I study hard so that I can pass the exam.} Ne dis pas \emph{so that to pass}.
\end{attention}

\courssub{Rédaction guidée : « How I am anticipating my second language needs at work »}

Observe d'abord le modèle suivant, rédigé par une élève de Terminale.

\begin{textebox}[A model paragraph]
My name is Larissa. I am in Terminale A in Bafoussam. I want to become a journalist, so English is a must for me. My level is intermediate, but my writing is weak. That is why I have signed up for an online course and I read English newspapers every day. In this way, I am anticipating my future language needs.
\end{textebox}

\begin{methode}[Rédiger le paragraphe en 4 étapes]
\begin{enumerate}
\item \textbf{Présente-toi} et annonce ton projet professionnel : \eng{My name is… I want to become a…}
\item \textbf{Évalue ton niveau actuel} : \eng{My level is… but my speaking/writing is weak.}
\item \textbf{Identifie tes besoins} : \eng{I will need English to… because…}
\item \textbf{Décris ton plan d'action} avec \eng{in order to} et \eng{so that} : \eng{I have signed up for… in order to… I practise every day so that I can…}
\end{enumerate}
Termine par une phrase de conclusion : \eng{In this way, I am anticipating my future language needs.}
\end{methode}

\begin{popfigure}
\begin{center}
\begin{tikzpicture}[node distance=0.9cm,
box/.style={rectangle, rounded corners, draw=popblue, fill=popblueL, text width=4.6cm, align=center, minimum height=1cm, font=\small}]
\node[box, align=center] (a){\textbf{1. My current level}\\\eng{My level is intermediate…}};
\node[box, below=of a, align=center] (b){\textbf{2. My future job}\\\eng{I want to become an engineer…}};
\node[box, below=of b, align=center] (c){\textbf{3. My language needs}\\\eng{I will need English to…}};
\node[box, below=of c, fill=popgreenL, draw=popgreen, align=center] (d){\textbf{4. My action plan}\\\eng{in order to… so that…}};
\draw[-{Stealth}, thick, popblue] (a) -- (b);
\draw[-{Stealth}, thick, popblue] (b) -- (c);
\draw[-{Stealth}, thick, popblue] (c) -- (d);
\end{tikzpicture}
\end{center}
\legende{Organigramme du plan de rédaction : du niveau actuel au plan d'action.}
\end{popfigure}

\begin{exoresolu}[Transformation de phrases]
Réécris avec le connecteur indiqué : 1. I learn English. I want to work abroad. (\emph{in order to}) 2. She attends evening classes. She wants to get promoted. (\emph{so that}) 3. My speaking is weak. I practise with a friend every week. (\emph{therefore})
\tcblower
1. \eng{I learn English in order to work abroad.} 2. \eng{She attends evening classes so that she can get promoted.} 3. \eng{My speaking is weak; therefore, I practise with a friend every week.}
\end{exoresolu}

\begin{exercice}[Production écrite]
Rédige un paragraphe de 8 à 10 lignes intitulé \eng{How I am anticipating my second language needs at work}, en suivant le plan imposé (niveau actuel → métier visé → besoins → plan d'action) et en utilisant au moins quatre connecteurs de la leçon.
\end{exercice}

\begin{methode}[Auto-évaluation : grille de relecture]
Avant de rendre ta copie, vérifie chaque point :
\begin{enumerate}
\item Ai-je utilisé au moins cinq mots du thème (\emph{level, needs, training, fluency, command of English}…) ?
\item Mes collocations sont-elles correctes (\eng{sign up for a course}, \eng{gain confidence}, \eng{a good command of English}) ?
\item Ai-je évité les faux amis (\emph{actually} = en fait, \emph{to achieve} = réaliser, pas « achever ») ?
\item \eng{in order to} + base verbale et \eng{so that} + proposition sont-ils bien construits ?
\item Ai-je varié les connecteurs (\emph{first, moreover, therefore, as a result}) ?
\item Orthographe et prononciation des mots de la même famille sont-elles correctes (\emph{employer}, \emph{employee}, \emph{development}) ?
\end{enumerate}
\end{methode}

\begin{savaistu}[Le paragraphe argumenté au baccalauréat]
Au Bac, les correcteurs attribuent les meilleures notes aux paragraphes qui combinent un \cle{vocabulaire thématique précis} et une \cle{cohérence} visible grâce à des connecteurs variés. Un texte court mais bien enchaîné, sans faux amis et avec des collocations naturelles, vaut toujours mieux qu'un long texte désorganisé. Entraîne-toi à rédiger un paragraphe de ce type chaque semaine !
\end{savaistu}

\newpage

\coursec{Activité d'intégration}

\courssub{Objectifs de l'activité}

À l'issue de cette activité d'intégration, l'élève sera capable de :

\begin{itemize}
\item réemployer le vocabulaire de la leçon (niveaux de langue, formation, besoins professionnels) dans une production orale et écrite authentique ;
\item identifier les besoins linguistiques d'un métier précis et les exprimer avec les collocations correctes (\eng{to attend a workshop}, \eng{to improve fluency}, \eng{to reach an intermediate level}) ;
\item évaluer un niveau de langue en utilisant l'échelle étudiée (\eng{beginner}, \eng{elementary}, \eng{intermediate}, \eng{upper-intermediate}, \eng{advanced}, \eng{fluent}) ;
\item proposer un plan de formation réaliste (\eng{online course}, \eng{evening classes}, \eng{language exchange}, \eng{refresher course}) ;
\item éviter les faux amis et les calques du français (\eng{formation} $\neq$ \eng{training}, \eng{stage} $\neq$ \eng{internship}) ;
\item présenter un projet à l'oral en trois minutes et répondre aux questions de l'auditoire.
\end{itemize}

\courssub{Situation}

\begin{textebox}[An announcement at Government Bilingual High School Yaoundé]
\textbf{Careers Week — English for Work Project}

Dear final-year students,

As part of Careers Week, our school is organising a competition called \textbf{“My Language Needs Action Plan”}. In our bilingual country, employers in Douala, Yaoundé, Buea and Bamenda increasingly require workers who can use English confidently at work: writing reports, answering phone calls, attending meetings and welcoming foreign visitors.

Working in groups of four, you will choose one profession — journalist, nurse, accountant, hotel receptionist or customs officer — and design an action plan answering three questions:

1. What English-language tasks does this job require?
2. What level of English is needed to perform them well?
3. What training plan will help a worker reach that level?

Each group will produce a poster and present it orally in three minutes. The class will ask questions. The best plans will be displayed in the school library and shared with our partner school in Buea.

The English Club Committee
\end{textebox}

\begin{definition}[Glossaire utile]
\begin{itemize}
\item \cle{careers week} (n.) : semaine des métiers.
\item \cle{employer} (n.) : employeur.
\item \cle{to require} (v.) : exiger, demander.
\item \cle{customs officer} (n.) : douanier.
\item \cle{action plan} (n.) : plan d'action.
\item \cle{to display} (v.) : afficher, exposer.
\end{itemize}
\end{definition}

\courssub{Consignes}

\textbf{Tâche 1 — Compréhension et choix du métier (10 min).} Lisez l'annonce en groupe. Soulignez les trois questions auxquelles votre plan doit répondre. Choisissez un métier dans la liste et justifiez votre choix en une phrase : \eng{We chose the job of nurse because hospitals in Yaoundé receive many English-speaking patients.}

\textbf{Tâche 2 — Identification des besoins linguistiques (15 min).} Complétez un tableau à trois colonnes : \eng{Tasks} (tâches en anglais), \eng{Skills} (\eng{listening, speaking, reading, writing}), \eng{Level needed}. Listez au moins cinq tâches réelles du métier. Exemple pour l'hôtelier : \eng{welcoming guests at the front desk} (speaking, intermediate), \eng{answering booking e-mails} (writing, upper-intermediate).

\textbf{Tâche 3 — Évaluation du niveau (10 min).} Pour chaque tâche, indiquez le niveau requis en utilisant l'échelle de la leçon. Rédigez deux phrases de synthèse avec les structures étudiées : \eng{A customs officer needs an upper-intermediate level to question travellers politely.} / \eng{He must be able to understand different accents.}

\textbf{Tâche 4 — Plan de formation (15 min).} Proposez un plan de formation en trois étapes avec un calendrier (\eng{first, then, finally} ; \eng{within six months} ; \eng{by the end of the year}). Utilisez au moins quatre mots de la leçon : \eng{online course}, \eng{workshop}, \eng{evening classes}, \eng{refresher course}, \eng{language exchange}, \eng{to enrol in}, \eng{to improve fluency}, \eng{to brush up on}.

\textbf{Tâche 5 — Vérification linguistique (10 min).} Relisez votre affiche à la chasse aux faux amis. Remplacez tout calque du français : \eng{formation} $\rightarrow$ \eng{training} ; \eng{stage} $\rightarrow$ \eng{internship} ; \eng{assister à une réunion} $\rightarrow$ \eng{to attend a meeting} (et non \eng{to assist}). Vérifiez les familles de mots : \eng{to train / training / trainer / trainee}.

\textbf{Tâche 6 — Affiche et présentation orale (20 min).} Rédigez l'affiche (titre, trois rubriques, phrases courtes). Répartissez la prise de parole : chaque membre parle environ 45 secondes. Présentez en trois minutes, puis répondez à deux questions de la classe.

\courssub{Ressources à mobiliser}

\begin{aretenir}[Boîte à outils linguistique]
\begin{itemize}
\item \textbf{Niveaux} : \eng{beginner, elementary, intermediate, upper-intermediate, advanced, fluent}. Structure : \eng{to have a good command of English}, \eng{to be proficient in English}.
\item \textbf{Besoins} : \eng{to need English for + nom} (\eng{for meetings}) ; \eng{to need to + base verbale} ; \eng{to be able to + base verbale}.
\item \textbf{Formation} : \eng{to enrol in an online course}, \eng{to attend a workshop}, \eng{to take evening classes}, \eng{to brush up on one's grammar}, \eng{to improve one's fluency and accuracy}.
\item \textbf{Familles de mots} : \eng{employ / employer / employee / employment} ; \eng{train / training / trainer / trainee} ; \eng{fluent / fluency / fluently}.
\item \textbf{Collocations} : \eng{to write a report}, \eng{to make a phone call}, \eng{to hold a meeting}, \eng{to meet a requirement}, \eng{to achieve a goal}.
\item \textbf{Faux amis à éviter} : \eng{actually} = en fait ; \eng{to assist} = aider ; \eng{formation} n'existe pas au sens de « formation professionnelle » : dire \eng{training}.
\end{itemize}
\end{aretenir}

\courssub{Proposition de production et corrigé commenté}

\begin{exoresolu}[Modèle d'affiche et de présentation orale — métier : hotel receptionist]
Énoncé : rédigez l'affiche du groupe et le script de la présentation orale de trois minutes.

\tcblower

\textbf{Poster — “English for a Hotel Receptionist: Our Action Plan”}

\eng{\textbf{1. Language needs.} A receptionist in a Douala hotel welcomes foreign guests at the front desk, answers phone calls and booking e-mails, gives information about prices and services, and handles complaints politely.}

\eng{\textbf{2. Level required.} Speaking and listening: upper-intermediate, because guests speak with different accents. Writing: intermediate, for short professional e-mails.}

\eng{\textbf{3. Training plan.} First, within three months, the receptionist will enrol in an online course to brush up on grammar. Then, he or she will attend a conversation workshop twice a week to improve fluency. Finally, by the end of the year, evening classes at the British Council will prepare him or her for an upper-intermediate certificate.}

\textbf{Script de présentation (extrait) :} \eng{Good morning, everyone. Our group chose the job of hotel receptionist because Douala is a business city with many international visitors. A receptionist needs a good command of spoken English to welcome guests and deal with complaints. We think an upper-intermediate level is necessary. Our training plan has three steps: an online course, a conversation workshop, and evening classes. Thank you for listening — do you have any questions?}

\textbf{Commentaires des choix de langue :}
\begin{itemize}
\item \eng{to welcome guests}, \eng{to handle complaints}, \eng{booking e-mails} : collocations authentiques du secteur hôtelier, sans calque du français.
\item \eng{within three months}, \eng{by the end of the year} : marqueurs temporels du plan d'action, avec le futur \eng{will} pour les étapes prévues.
\item \eng{to brush up on grammar}, \eng{to improve fluency}, \eng{to enrol in} : réemploi exact du lexique de formation ; on évite le faux ami \eng{formation}.
\item \eng{he or she} : anglais inclusif correct ; prononciation de \eng{receptionist} : « ri-sèp-cheu-niste », accent sur la deuxième syllabe.
\item Questions possibles de la classe et réponses : \eng{Why not an advanced level? — Because the e-mails are short and routine; intermediate writing is enough.}
\end{itemize}
\end{exoresolu}

\courssub{Bilan de l'activité}

Cette activité a permis de mobiliser l'ensemble des savoir-faire de la leçon dans une tâche finale authentique et motivante. Sur le plan lexical, les élèves ont réemployé l'échelle des niveaux, le vocabulaire de la formation et les collocations professionnelles dans un contexte réel : le monde du travail camerounais bilingue. Sur le plan méthodologique, ils ont appris à construire un raisonnement en trois temps — besoins, niveau, plan — compétence transférable à toute lettre de motivation ou entretien d'embauche. Sur le plan de la correction, la phase de vérification des faux amis a montré que la vigilance face aux calques du français (\eng{training} et non \eng{formation}, \eng{to attend} et non \eng{to assist}) est un réflexe à entretenir. L'évaluation par les pairs, lors des questions-réponses, a enfin développé la compréhension orale et la capacité à réagir spontanément en anglais. Pour aller plus loin, chaque élève pourra rédiger individuellement son propre \eng{Personal Language Needs Action Plan} pour son projet d'orientation après le baccalauréat.

\newpage

\coursec{Méthodes et exercices résolus}

\coursec{Lesson 11 — Words and expressions related to anticipating one's second language needs at work}

\courssub{Savoir-faire 1 : Identifier et employer le vocabulaire du thème}

\begin{methode}[Repérer et utiliser le lexique des besoins linguistiques au travail]
Pour maîtriser le vocabulaire de l'anticipation des besoins en anglais au travail (\eng{anticipating one's second language needs at work}), suis ces étapes :
\begin{enumerate}
\item \textbf{Repère le champ lexical} dans le texte ou l'enregistrement : souligne les mots liés à la formation (\eng{training, course, workshop, refresher course}), aux compétences (\eng{fluency, proficiency, language skills}), aux besoins (\eng{needs assessment, requirement, gap}) et à la carrière (\eng{promotion, job interview, career prospects}).
\item \textbf{Classe les mots par famille} : verbe, nom, adjectif. Exemple : \eng{to improve} (verbe) → \eng{improvement} (nom) → \eng{improved} (adjectif).
\item \textbf{Mémorise les collocations} (associations de mots naturelles) : \eng{to attend a training course} (suivre une formation), \eng{to brush up one's English} (rafraîchir son anglais), \eng{to meet the requirements} (répondre aux exigences), \eng{to apply for a job} (postuler à un emploi).
\item \textbf{Vérifie le sens dans le contexte} : ne traduis jamais mot à mot. \eng{Actually} ne veut pas dire « actuellement » mais « en fait ».
\item \textbf{Réemploie le mot dans ta propre phrase} : c'est la seule façon de le retenir durablement.
\end{enumerate}
\end{methode}

\begin{exoresolu}[Vocabulaire en contexte — type Bac]
Read the text and answer the questions.

\begin{textebox}[Extrait original]
“When Mrs Etoundi was promoted to head of the import department in Douala, she realised her English was not good enough to negotiate with foreign partners. She decided to enrol in an evening language course and to brush up her business vocabulary. After six months of intensive training, her fluency improved so much that she could chair meetings in English.”
\end{textebox}

\textbf{Questions :}
\begin{enumerate}
\item Find in the text two expressions related to language training.
\item What does the phrasal verb “to brush up” mean?
\item Give the noun corresponding to the adjective “fluent”.
\end{enumerate}
\tcblower
\textbf{Solution détaillée :}
\begin{enumerate}
\item \textbf{Étape 1 :} on cherche les expressions du champ lexical de la formation. Le texte contient : \eng{to enrol in an evening language course} (s'inscrire à un cours du soir) et \eng{six months of intensive training} (six mois de formation intensive). On peut aussi citer \eng{language course}.
\item \textbf{Étape 2 :} le contexte aide à deviner \eng{to brush up} : Mrs Etoundi reprend son vocabulaire des affaires avant de négocier. Sens : « rafraîchir, réviser ». Réponse rédigée : \eng{“To brush up” means to revise and improve knowledge or a skill that one has partly forgotten.} (« Réviser et améliorer une compétence qu'on a en partie oubliée. »)
\item \textbf{Étape 3 :} famille de mots : adjectif \eng{fluent} → nom \eng{fluency} (attention à l'orthographe : \eng{fluency}, terminaison en \emph{-cy}, pas \emph{-ncy} ni \emph{-ce}). Le texte le confirme : \eng{her fluency improved}.
\end{enumerate}
\textbf{Conclusion :} toujours justifier la réponse par un élément du texte et vérifier l'orthographe des dérivés.
\end{exoresolu}

\courssub{Savoir-faire 2 : Évaluer son niveau et exprimer ses besoins linguistiques}

\begin{methode}[Parler de son niveau et de ses objectifs en anglais]
Pour décrire ton niveau et tes besoins en anglais (à l'oral comme à l'écrit) :
\begin{enumerate}
\item \textbf{Utilise l'échelle des niveaux} : \eng{beginner} (débutant), \eng{elementary}, \eng{intermediate} (intermédiaire), \eng{upper-intermediate}, \eng{advanced} (avancé), \eng{fluent} (courant).
\item \textbf{Exprime un besoin} avec les structures : \eng{I need to improve my writing skills.} / \eng{I have difficulty understanding native speakers.} / \eng{My English is a bit rusty.} (mon anglais est rouillé).
\item \textbf{Exprime un objectif} : \eng{I would like to take a refresher course.} / \eng{My goal is to become fluent within two years.} / \eng{I want to be able to write professional emails.}
\item \textbf{Exprime un manque à combler} : \eng{There is a gap between my reading and speaking skills.} / \eng{I lack confidence when I speak.}
\item \textbf{Relie besoin et solution} avec \eng{so, therefore, that is why} : \eng{My listening is weak, so I have signed up for a conversation club.}
\end{enumerate}
\end{methode}

\begin{exoresolu}[Expression des besoins — production guidée]
\textbf{Énoncé :} You are applying for a job in a bank in Buea. Write five sentences in English in which you (a) state your current level of English, (b) mention one weakness, (c) say what you need, (d) explain your goal, (e) propose a solution.
\tcblower
\textbf{Raisonnement :} on applique les cinq étapes de la méthode, une phrase par fonction, avec des structures variées et un vocabulaire précis.

\textbf{Réponse rédigée (modèle) :}

\eng{(a) I would describe my level of English as intermediate.}

\eng{(b) However, I sometimes have difficulty understanding fast native speakers on the phone.}

\eng{(c) I therefore need to improve my listening skills and my professional vocabulary.}

\eng{(d) My goal is to become fluent enough to deal with English-speaking customers confidently.}

\eng{(e) That is why I intend to enrol in an evening refresher course at the language centre.}

\textbf{Justifications :} \eng{intermediate} est le niveau correct (ni \emph{medium} ni \emph{middle}) ; \eng{to have difficulty + -ing} (et non « to understand » après \eng{difficulty}) ; \eng{to enrol in a course} avec la préposition \eng{in} ; \eng{That is why} (et non le calque « that is for what »).

\textbf{Conclusion :} une fonction = une structure apprise = une phrase claire et correcte.
\end{exoresolu}

\courssub{Savoir-faire 3 : Réemployer le lexique à l'oral et à l'écrit}

\begin{methode}[Réutiliser le vocabulaire dans un dialogue ou un paragraphe]
\begin{enumerate}
\item \textbf{Prépare une banque de 8 à 10 mots} du thème avant de parler ou d'écrire : \eng{fluency, requirement, interview, training, certificate, to apply, skills, promotion}.
\item \textbf{À l'oral (jeu de rôle, entretien)} : utilise des formules d'interaction : \eng{Could you tell me more about...?} / \eng{In my opinion, ...} / \eng{What I mean is...}. Parle en phrases complètes, jamais en mots isolés.
\item \textbf{À l'écrit} : place au moins un mot du thème par phrase, mais sans forcer ; varie les verbes (\eng{to improve, to develop, to acquire, to strengthen}).
\item \textbf{Utilise des connecteurs} pour structurer : \eng{first, moreover, however, as a result, in conclusion}.
\item \textbf{Relis-toi} : vérifie la troisième personne du singulier (\eng{she needs}), les prépositions (\eng{to apply \textbf{for} a job}, \eng{interested \textbf{in}}) et l'orthographe.
\end{enumerate}
\end{methode}

\begin{exoresolu}[Dialogue — réemploi du lexique]
\textbf{Énoncé :} Complete the job interview with words from the list: \emph{fluency, requirements, training, certificate, apply, skills}.

\eng{Interviewer: Why did you decide to (1) \dots\ for this position?}

\eng{Candidate: Because I meet all the (2) \dots\ : I have a degree and a language (3) \dots\ from the British Council.}

\eng{Interviewer: How would you describe your English (4) \dots\ ?}

\eng{Candidate: My writing is excellent, but I followed a (5) \dots\ course to improve my speaking. I would say my (6) \dots\ is now very good.}
\tcblower
\textbf{Solution détaillée :}
\begin{enumerate}
\item \eng{apply} — collocation obligatoire : \eng{to apply for a position} (postuler à un poste). Ne jamais dire « apply a job ».
\item \eng{requirements} — \eng{to meet the requirements} = « répondre aux exigences » ; le pluriel est naturel ici (\emph{all the}).
\item \eng{certificate} — un diplôme ou certificat de langue ; \eng{a language certificate}. Attention au faux ami : \eng{certificate} ≠ « certificat » médical uniquement ; ici c'est bien une attestation de niveau.
\item \eng{skills} — \eng{English skills} = « compétences en anglais » ; le pluriel est la norme.
\item \eng{training} — collocation : \eng{a training course} (une formation).
\item \eng{fluency} — nom dérivé de \eng{fluent} ; orthographe : f-l-u-e-n-c-y.
\end{enumerate}
\textbf{Conclusion :} chaque réponse repose sur une collocation apprise ; c'est la collocation, et non le mot isolé, qu'il faut mémoriser.
\end{exoresolu}

\courssub{Savoir-faire 4 : Éviter les faux amis et les calques du français}

\begin{methode}[Déjouer les faux amis du monde professionnel]
\begin{enumerate}
\item \textbf{Repère les mots « trop faciles »} : si un mot anglais ressemble à un mot français, méfie-toi. Fais une pause et vérifie.
\item \textbf{Apprends les paires par cœur} :
\end{enumerate}
\begin{center}
\begin{tabularx}{\linewidth}{|X|X|X|}
\hline
\rowcolor{popblueL} Faux ami & Vrai sens anglais & Pour dire le mot français \\
\hline
\eng{actually} & en fait & \eng{currently} (actuellement) \\
\hline
\eng{a formation} (inexistant) & --- & \eng{training} (une formation) \\
\hline
\eng{to assist} & aider & \eng{to attend} (assister à) \\
\hline
\eng{a stage} & une étape, une scène & \eng{an internship / a work placement} (un stage) \\
\hline
\eng{a diploma} & diplôme (titre précis, souvent professionnel) & \eng{a degree} (un diplôme universitaire) \\
\hline
\eng{to demand} & exiger & \eng{to ask for} (demander) \\
\hline
\end{tabularx}
\end{center}
\begin{enumerate}
\item[3.] \textbf{Évite les calques de structure} : « I have 25 years » → \eng{I am 25 (years old)} ; « I am agree » → \eng{I agree} ; « to make a formation » → \eng{to take a training course}.
\item[4.] \textbf{En cas de doute à l'examen}, choisis le mot simple et sûr plutôt que le mot savant risqué.
\end{enumerate}
\end{methode}

\begin{exoresolu}[Correction de phrases — faux amis et calques]
\textbf{Énoncé :} Each sentence contains one mistake (false friend or word-for-word translation). Correct it.
\begin{enumerate}
\item \eng{I need to make a formation in business English.}
\item \eng{Actually, I work in a travel agency in Yaoundé.}
\item \eng{She assisted the meeting last Monday.}
\item \eng{I am agree that English is essential for my career.}
\end{enumerate}
\tcblower
\textbf{Solution détaillée :}
\begin{enumerate}
\item \eng{make a formation} est un double calque : \eng{formation} n'existe pas en anglais dans ce sens, et le verbe est \eng{take}. Correction : \eng{I need to take a training course in business English.} (« J'ai besoin de suivre une formation en anglais des affaires. »)
\item \eng{Actually} signifie « en fait ». Pour « actuellement », il faut \eng{currently} ou \eng{at present}. Correction : \eng{At present, I work in a travel agency in Yaoundé.}
\item \eng{To assist} = « aider ». Pour « assister à », le verbe est \eng{to attend} (sans préposition). Correction : \eng{She attended the meeting last Monday.}
\item \eng{Agree} est un verbe en anglais, pas un adjectif : on ne met jamais \eng{am} devant. Correction : \eng{I agree that English is essential for my career.}
\end{enumerate}
\textbf{Conclusion :} la méthode est toujours la même : repérer le mot suspect, vérifier son vrai sens, choisir l'équivalent anglais authentique.
\end{exoresolu}

\courssub{Savoir-faire 5 : Familles de mots, prononciation et orthographe}

\begin{methode}[Exploiter les familles de mots et sécuriser l'orthographe]
\begin{enumerate}
\item \textbf{Construis le tableau de la famille} dès que tu rencontres un mot nouveau :
\end{enumerate}
\begin{center}
\begin{tabular}{|l|l|l|}
\hline
\rowcolor{popblueL} Verbe & Nom & Adjectif \\
\hline
\eng{to employ} & \eng{employment / employer / employee} & \eng{employed / employable} \\
\hline
\eng{to train} & \eng{training / trainer / trainee} & \eng{trained} \\
\hline
\eng{to apply} & \eng{application / applicant} & \eng{applicable} \\
\hline
\eng{to require} & \eng{requirement} & \eng{required} \\
\hline
\eng{to improve} & \eng{improvement} & \eng{improved} \\
\hline
\end{tabular}
\end{center}
\begin{enumerate}
\item[2.] \textbf{Attention aux pièges d'orthographe} :
\begin{itemize}
\item \eng{applicant} / \eng{application} / \eng{applicable} : double \textbf{« p »}, simple \textbf{« l »} partout.
\item \eng{business} : double \textbf{« s »} final (prononcé « biz-nèss »).
\item \eng{professional} : un \textbf{« f »}, double \textbf{« s »} (pro-fes-sion-al).
\item \eng{requirement} : pas de « i » après le « r » (\emph{requirment} est faux).
\item \eng{fluency} : terminaison \textbf{-cy} (pas \emph{-ncy} ni \emph{-ce}).
\end{itemize}
\item[3.] \textbf{Note la prononciation en lettres françaises} : \eng{fluency} « flou-an-si », \eng{requirement} « ri-kouaïr-ment », \eng{certificate} « seur-ti-fi-kèt », \eng{interview} « in-teur-viou ».
\item[4.] \textbf{Choisis la bonne forme selon la fonction dans la phrase} : après un article, un nom ; après un auxiliaire, un verbe ; devant un nom, un adjectif.
\end{enumerate}
\end{methode}

\begin{exoresolu}[Familles de mots — transformation]
\textbf{Énoncé :} Complete each sentence with the correct form of the word in brackets.
\begin{enumerate}
\item \eng{The company is looking for a new \dots\ . (employ)}
\item \eng{Good English is a \dots\ for this job. (require)}
\item \eng{Her \dots\ in spoken English impressed the panel. (improve)}
\item \eng{All \dots\ must send a CV before Friday. (apply)}
\end{enumerate}
\tcblower
\textbf{Solution détaillée :}
\begin{enumerate}
\item Après l'article \eng{a} et l'adjectif \eng{new}, il faut un \textbf{nom de personne}. Le contexte (« the company is looking for ») indique la personne recrutée : \eng{employee} (employé). \eng{Employer} serait le patron — sens contraire.
\item Après l'article \eng{a}, il faut un \textbf{nom de chose} : \eng{requirement} (exigence). Phrase : \eng{Good English is a requirement for this job.} Attention à la préposition : \eng{a requirement \textbf{for} a job}.
\item Après le possessif \eng{Her}, il faut un \textbf{nom} : \eng{improvement} (verbe \eng{improve} + suffixe \emph{-ment}). Orthographe : un seul « m » ajouté au verbe.
\item Après le quantifieur \eng{All}, il faut un \textbf{nom pluriel de personnes} : \eng{applicants} (candidats). Orthographe : double « p », simple « l » (a-p-p-l-i-c-a-n-t-s).
\end{enumerate}
\textbf{Conclusion :} la fonction grammaticale (article + nom, possessif + nom) dicte la forme du mot ; le contexte dicte le sens (employer/employee).
\end{exoresolu}

\begin{attention}[Les 5 erreurs qui coûtent des points au Bac]
\begin{enumerate}
\item \textbf{Le mot « formation » traduit mot à mot.} En anglais, « une formation » se dit \eng{a training course} ou \eng{training}, jamais « a formation ». De même, « un stage » = \eng{an internship}, pas « a stage ».
\item \textbf{Oublier la préposition dans les collocations clés.} On écrit \eng{to apply \textbf{for} a job}, \eng{to enrol \textbf{in} a course}, \eng{interested \textbf{in} English}, \eng{to depend \textbf{on}}, \eng{to listen \textbf{to}}. Une préposition fausse ou absente coûte un point à chaque occurrence.
\item \textbf{Les calques du français.} « I am agree » → \eng{I agree} ; « I have twenty years » → \eng{I am twenty} ; « actually » pour « actuellement » → \eng{currently}. Le correcteur repère immédiatement ces traductions mot à mot.
\item \textbf{L'orthographe des mots du thème.} Les fautes les plus fréquentes : \eng{fluency} (pas « fluancy »), \eng{professional} (un « f », \textbf{deux « s »}), \eng{business} (double « s » final), \eng{requirement} (pas « requirment »), \eng{applicant} (double « p », simple « l »). Une faute dans un mot du vocabulaire de la leçon est doublement pénalisée.
\item \textbf{L'ordre des mots et la troisième personne.} L'adjectif se place \textbf{avant} le nom : \eng{a language course} (pas « a course of language » si l'adjectif existe). Et au présent simple : \eng{she need\textbf{s} to improve her English} — l'oubli du « s » est l'erreur la plus sanctionnée à l'écrit.
\end{enumerate}
\end{attention}

\newpage

\coursec{Exercices}

\courssub{Je vérifie mes connaissances}

\begin{exercice}[QCM : le bon mot]
Choose the correct word to complete each sentence. Circle the letter of your answer.
\begin{enumerate}
\item Our company organises a \ldots\ course every year to improve the staff's English.\\
A. formation \hspace{0.5cm} B. training \hspace{0.5cm} C. information \hspace{0.5cm} D. instruction
\item I need to \ldots\ my English before the conference in Lagos.\\
A. brush up \hspace{0.5cm} B. wash up \hspace{0.5cm} C. clean up \hspace{0.5cm} D. paint up
\item She has a good \ldots\ of English: she can negotiate with foreign partners.\\
A. command \hspace{0.5cm} B. commandment \hspace{0.5cm} C. order \hspace{0.5cm} D. domination
\item The job advertisement says that candidates must be \ldots\ in English and French.\\
A. liquid \hspace{0.5cm} B. fluent \hspace{0.5cm} C. fluid \hspace{0.5cm} D. floating
\item My \ldots\ decided to pay for my evening classes at the language centre.\\
A. employer \hspace{0.5cm} B. employee \hspace{0.5cm} C. employment \hspace{0.5cm} D. employ
\item To work in this bank, you need a language \ldots\ such as IELTS or TOEFL.\\
A. certificate \hspace{0.5cm} B. certitude \hspace{0.5cm} C. certifier \hspace{0.5cm} D. cereal
\end{enumerate}
\end{exercice}

\begin{exercice}[Vrai ou faux ?]
Read the following text, then say whether each statement is True or False. Justify each answer with a quotation from the text.

\begin{textebox}[English needs at Cameroun Export Ltd]
At Cameroun Export Ltd, a cocoa trading company based in Douala, English has become essential. The sales manager, Mr Etoundi, explains: “Most of our buyers are in Nigeria, Ghana and the United Kingdom. Every contract, every e-mail and every phone call is in English.” Last year, the company tested its forty employees. Only twelve could hold a professional conversation in English. The management then decided to organise evening classes twice a week with a certified trainer. Employees who reach the B2 level will receive a bonus and will be allowed to travel abroad for trade fairs. “We are investing in language because language is our working tool,” concludes Mr Etoundi.
\end{textebox}

\begin{enumerate}
\item Cameroun Export Ltd sells cocoa to French-speaking countries only.
\item All the employees of the company speak good English.
\item The company tested its staff's English last year.
\item English classes take place every day.
\item Employees with a B2 level will get extra money.
\item Mr Etoundi thinks English is not important for his company.
\end{enumerate}
\end{exercice}

\begin{exercice}[Familles de mots]
Complete the table with the correct word from the same family, then use each word in a sentence of your own about work and language learning.

\begin{center}
\begin{tabular}{|l|l|l|}
\hline
\rowcolor{popblueL} Verb & Person (noun) & Abstract noun \\
\hline
to employ & \ldots\ldots\ldots\ldots & \ldots\ldots\ldots\ldots \\
\hline
to train & \ldots\ldots\ldots\ldots & \ldots\ldots\ldots\ldots \\
\hline
to apply & \ldots\ldots\ldots\ldots & \ldots\ldots\ldots\ldots \\
\hline
to qualify & \ldots\ldots\ldots\ldots & \ldots\ldots\ldots\ldots \\
\hline
\end{tabular}
\end{center}
\end{exercice}

\begin{exercice}[Complète avec le bon mot]
Complete the sentences with the following words: \eng{level, skills, course, fluent, improve, requirements}.
\begin{enumerate}
\item My English \ldots\ is intermediate, but I want to reach an advanced level.
\item This position has two language \ldots : English and French.
\item Good communication \ldots\ are essential in international trade.
\item She enrolled in a three-month English \ldots\ at the British Council.
\item Reading the press every day helps me \ldots\ my vocabulary.
\item He is \ldots\ in three languages: English, French and Spanish.
\end{enumerate}
\end{exercice}

\courssub{Je m'entraîne}

\begin{exercice}[Traduction]
Translate the following sentences into English. Pay attention to false friends and word order.
\begin{enumerate}
\item Je dois améliorer mon anglais pour mon travail.
\item Mon entreprise organise une formation en anglais tous les ans.
\item Elle a un bon niveau d'anglais : elle peut négocier avec les clients étrangers.
\item J'aimerais suivre des cours du soir pour obtenir un certificat.
\item La maîtrise de l'anglais est indispensable dans le commerce international.
\end{enumerate}
\end{exercice}

\begin{exercice}[Corrige les faux amis]
Each of the following sentences contains a false friend or a calque from French. Underline the mistake and rewrite the sentence correctly.
\begin{enumerate}
\item I assisted a formation to ameliorate my English.
\item She is actually following an English stage in Buea.
\item My society sends me to conferences abroad every year.
\item He has good knowledges of English and computer science.
\item I passed an examination to obtain my linguistic diploma last June.
\item The formation will take place during three months.
\item I am very motivated to perfect my English for my career.
\item She made her studies in a bilingual college in Yaoundé.
\end{enumerate}
\end{exercice}

\begin{exercice}[Texte à trous]
Complete the following text with the words from the box. There are two extra words.

\begin{center}
\eng{training — fluent — needs — certificate — brush up — apply — skills — level — interview — salary — course — improve}
\end{center}

\begin{textebox}[Amina's plan]
Amina works for a bank in Yaoundé. Her manager told her that she must \ldots\ her English if she wants a promotion. Her current \ldots\ is A2, which is not enough to speak with foreign customers. She has identified her \ldots : she must work on her speaking and writing \ldots . She has decided to enrol in an evening \ldots\ at a language centre in Bastos. After six months of \ldots , she hopes to pass an international \ldots\ such as the TOEIC. Then she will be able to \ldots\ for the post of customer relations officer, which requires \ldots\ English.
\end{textebox}
\end{exercice}

\begin{exercice}[Appariement : collocations]
Match the beginning of each expression (1--10) with its correct ending (a--j).

\begin{center}
\begin{tabular}{|l|l|}
\hline
\rowcolor{popblueL} Beginning & Ending \\
\hline
1. to brush up & a. the requirements \\
\hline
2. to meet & b. one's English \\
\hline
3. to attend & c. a language course \\
\hline
4. to improve & d. for a job \\
\hline
5. to apply & e. a certificate \\
\hline
6. to obtain & f. one's skills \\
\hline
7. to have a good command & g. a promotion \\
\hline
8. to get & h. of English \\
\hline
9. to sit & i. an examination \\
\hline
10. to enrol & j. in a training centre \\
\hline
\end{tabular}
\end{center}

Then write five sentences about your own English learning using five of these collocations.
\end{exercice}

\courssub{Je me prépare au Bac}

\begin{exercice}[Compréhension de texte — type Bac]
Read the following text carefully and answer the questions in English, in your own words as far as possible.

\begin{textebox}[Why Cameroonian companies invest in English]
In Douala, the economic capital of Cameroon, more and more companies are investing in English training for their staff. The reason is simple: business in Central and West Africa is increasingly conducted in English. Nigeria, Ghana and Kenya are major trading partners, and English is the language of contracts, e-mails and international meetings.

Take the example of Njomba Logistics, a transport company with one hundred and twenty employees. Three years ago, the director realised that his teams were losing contracts because they could not communicate with Nigerian clients. “We had the trucks, we had the drivers, but we did not have the language,” he recalls. The company first assessed the English level of every employee. Then it designed a training plan: beginners attended general English classes, while intermediate learners followed a specialised course in business English. The company paid eighty per cent of the fees; the employees paid the rest, which motivated them to attend regularly.

The results came quickly. After eighteen months, the company signed three new contracts with partners in Lagos and Accra. Employees who completed the training received a certificate and a salary increase. “Anticipating our language needs was the best investment we ever made,” the director concludes. “Today, when we recruit, we always ask candidates about their English level and their willingness to learn.”
\end{textebox}

\textbf{A. Comprehension questions}
\begin{enumerate}
\item Why are Cameroonian companies investing in English training? Give two reasons.
\item What problem did Njomba Logistics face three years ago?
\item What does the director mean when he says: “We had the trucks, we had the drivers, but we did not have the language”?
\item Describe the two steps of the company's training plan.
\item How did the company share the cost of the training, and why?
\item Mention two benefits the company obtained after the training.
\end{enumerate}

\textbf{B. Language work}
\begin{enumerate}
\item Find in the text the noun corresponding to: \eng{to employ} (paragraph 3) and the verb corresponding to \eng{a certificate} (paragraph 2).
\item Find in the text words or expressions that mean: \emph{évaluer} (paragraph 2), \emph{une augmentation de salaire} (paragraph 3), \emph{la volonté d'apprendre} (paragraph 3).
\item Rewrite in the reported speech: “Anticipating our language needs was the best investment we ever made,” the director concludes.
\end{enumerate}
\end{exercice}

\begin{exercice}[Traduction — type Bac]
Translate the following passage into English.

« De nos jours, l'anglais est devenu indispensable dans le monde du travail. Au Cameroun, de nombreuses entreprises exigent que leurs employés parlent couramment l'anglais. C'est pourquoi j'ai décidé de suivre une formation dans un centre de langues. Je veux améliorer mon niveau, obtenir un certificat international et ainsi décrocher une promotion. Anticiper ses besoins linguistiques, c'est préparer son avenir professionnel. »
\end{exercice}

\begin{exercice}[Expression écrite — type Bac]
Choose \textbf{one} of the following topics and write your composition.

\textbf{Topic 1.} You want to attend an English training course at a language centre. Write a paragraph of \textbf{80 to 100 words} to convince your employer to pay for it. Explain your current level, your needs, your training plan and the benefits for the company.

\textbf{Topic 2.} Your school is organising a debate on the following motion: “Every Cameroonian worker should be trained in English.” Write a composition of \textbf{150 to 180 words} giving your opinion, with at least three arguments and one concrete example.

\textbf{Topic 3 (role play, oral practice).} In pairs, prepare a job interview. The candidate must describe his or her English level, language certificates and training project; the interviewer asks at least five questions about the candidate's language skills. Be ready to perform in front of the class.
\end{exercice}

\newpage

\coursec{Corrigés des exercices}

\begin{corrige}[Exercice 1 — QCM : le bon mot]
\begin{enumerate}
\item \textbf{B. training}. \eng{Our company organises a training course every year to improve the staff's English.} « Notre entreprise organise une formation chaque année pour améliorer l'anglais du personnel. » \eng{Formation} est un faux ami : en anglais, \eng{formation} signifie « formation (géologique, d'une équipe) », jamais « cours de formation professionnelle ».
\item \textbf{A. brush up}. \eng{I need to brush up my English before the conference in Lagos.} « Je dois rafraîchir mon anglais avant la conférence de Lagos. » Les autres verbes à particule n'existent pas dans ce sens (\eng{wash up} = faire la vaisselle).
\item \textbf{A. command}. \eng{She has a good command of English.} « Elle a une bonne maîtrise de l'anglais. » Collocation figée : \eng{to have a good command of a language}.
\item \textbf{B. fluent}. \eng{Candidates must be fluent in English and French.} « Les candidats doivent parler couramment l'anglais et le français. » Attention à la préposition : \eng{fluent \textbf{in}}.
\item \textbf{A. employer}. \eng{My employer decided to pay for my evening classes.} « Mon employeur a décidé de payer mes cours du soir. » \eng{Employee} = employé (celui qui reçoit le salaire, pas celui qui paie).
\item \textbf{A. certificate}. \eng{You need a language certificate such as IELTS or TOEFL.} « Il faut un certificat de langue comme l'IELTS ou le TOEFL. » \eng{Certitude} est un faux ami (\eng{certainty} en anglais) ; \eng{cereal} = céréale.
\end{enumerate}
\end{corrige}

\begin{corrige}[Exercice 2 — Vrai ou faux ?]
\begin{enumerate}
\item \textbf{False.} The buyers are in Nigeria, Ghana and the United Kingdom, which are English-speaking countries: “Most of our buyers are in Nigeria, Ghana and the United Kingdom.”
\item \textbf{False.} Only twelve employees out of forty could speak English professionally: “Only twelve could hold a professional conversation in English.”
\item \textbf{True.} “Last year, the company tested its forty employees.”
\item \textbf{False.} The classes take place twice a week, not every day: “evening classes twice a week”.
\item \textbf{True.} A bonus is extra money: “Employees who reach the B2 level will receive a bonus.”
\item \textbf{False.} He considers English as essential: “language is our working tool”.
\end{enumerate}
\textbf{Point de méthode :} au Bac, la justification par une citation exacte du texte est obligatoire ; une réponse \eng{True/False} sans citation ne rapporte que la moitié des points.
\end{corrige}

\begin{corrige}[Exercice 3 — Familles de mots]
\begin{center}
\begin{tabular}{|l|l|l|}
\hline
\rowcolor{popblueL} Verb & Person (noun) & Abstract noun \\
\hline
to employ & employer / employee & employment \\
\hline
to train & trainer / trainee & training \\
\hline
to apply & applicant & application \\
\hline
to qualify & qualified person / qualifier & qualification \\
\hline
\end{tabular}
\end{center}
\textbf{Exemples de phrases possibles :}
\begin{itemize}
\item \eng{My employer pays for my English training.} « Mon employeur paie ma formation d'anglais. »
\item \eng{The trainer gives us useful speaking exercises.} « Le formateur nous donne des exercices d'expression orale utiles. »
\item \eng{The applicant sent her application letter last week.} « La candidate a envoyé sa lettre de candidature la semaine dernière. »
\item \eng{A language qualification is a real advantage on the job market.} « Un diplôme de langue est un vrai atout sur le marché du travail. »
\end{itemize}
\textbf{Points de vigilance :} \eng{employer} (celui qui emploie) et \eng{employee} (celui qui est employé) ; \eng{trainer} (formateur) et \eng{trainee} (stagiaire en formation). Ne pas écrire \eng{applier} : le nom de personne est \eng{applicant}.
\end{corrige}

\begin{corrige}[Exercice 4 — Complète avec le bon mot]
\begin{enumerate}
\item \eng{My English \textbf{level} is intermediate, but I want to reach an advanced level.} « Mon niveau d'anglais est intermédiaire, mais je veux atteindre un niveau avancé. »
\item \eng{This position has two language \textbf{requirements}: English and French.} « Ce poste exige deux langues : l'anglais et le français. »
\item \eng{Good communication \textbf{skills} are essential in international trade.} « De bonnes compétences en communication sont essentielles dans le commerce international. »
\item \eng{She enrolled in a three-month English \textbf{course} at the British Council.} « Elle s'est inscrite à un cours d'anglais de trois mois au British Council. »
\item \eng{Reading the press every day helps me \textbf{improve} my vocabulary.} « Lire la presse tous les jours m'aide à améliorer mon vocabulaire. »
\item \eng{He is \textbf{fluent} in three languages: English, French and Spanish.} « Il parle couramment trois langues : l'anglais, le français et l'espagnol. »
\end{enumerate}
\textbf{Point de grammaire :} après \eng{help}, on utilise la base verbale avec ou sans \eng{to} : \eng{helps me improve} ou \eng{helps me to improve}.
\end{corrige}

\begin{corrige}[Exercice 5 — Traduction]
\begin{enumerate}
\item \eng{I must improve my English for my work.} / \eng{I need to improve my English for my job.} — « Améliorer » se dit \eng{improve}, jamais \eng{ameliorate} (faux ami, très rare en anglais courant).
\item \eng{My company organises an English training course every year.} — « Formation » = \eng{training (course)}, pas \eng{formation}. « Tous les ans » = \eng{every year}, en fin de phrase.
\item \eng{She has a good command of English: she can negotiate with foreign customers.} — Collocation : \eng{a good command of}. « Clients » = \eng{customers} ou \eng{clients}.
\item \eng{I would like to attend evening classes to obtain a certificate.} — « Suivre des cours » = \eng{attend classes} ; \eng{I would like} traduit « j'aimerais » (conditionnel de politesse).
\item \eng{A good command of English is essential in international trade.} — « Le commerce international » = \eng{international trade} ; « indispensable » = \eng{essential} ou \eng{indispensable}.
\end{enumerate}
\textbf{Points de vigilance :} ordre des mots (le complément de temps en fin de phrase) ; pas de calque \eng{follow courses} : on dit \eng{attend a course} ou \eng{take a course}.
\end{corrige}

\begin{corrige}[Exercice 6 — Corrige les faux amis]
\begin{enumerate}
\item \eng{\underline{assisted a formation to ameliorate}} $\rightarrow$ \eng{I attended a training course to improve my English.} — \eng{Assist} = aider ; \eng{formation} et \eng{ameliorate} sont des calques du français.
\item \eng{\underline{actually following an English stage}} $\rightarrow$ \eng{She is currently doing an English internship in Buea.} — \eng{Actually} = en fait ; « actuellement » = \eng{currently}. « Stage » = \eng{internship} ou \eng{placement}.
\item \eng{\underline{society}} $\rightarrow$ \eng{My company sends me to conferences abroad every year.} — \eng{Society} = société (au sens de communauté humaine) ; « entreprise » = \eng{company} ou \eng{firm}.
\item \eng{\underline{knowledges}} $\rightarrow$ \eng{He has a good knowledge of English and computer science.} — \eng{Knowledge} est indénombrable : jamais de \eng{-s}.
\item \eng{\underline{linguistic diploma}} $\rightarrow$ \eng{I sat an examination to obtain my language certificate last June.} — \eng{Linguistic} = linguistique (la science du langage) ; « diplôme » au sens de certificat de langue = \eng{certificate}.
\item \eng{\underline{during three months}} $\rightarrow$ \eng{The training course will last (for) three months.} — \eng{During} + nom (\eng{during the holidays}) ; \eng{for} + durée.
\item \eng{\underline{perfect my English}} $\rightarrow$ \eng{I am very motivated to improve my English for my career.} — \eng{Perfect} comme verbe existe mais est rare ; « perfectionner/améliorer » = \eng{improve}.
\item \eng{\underline{made her studies}} $\rightarrow$ \eng{She studied at a bilingual secondary school in Yaoundé.} — « Faire ses études » = \eng{to study} ; « collège » = \eng{secondary school} (\eng{college} = université).
\end{enumerate}
\end{corrige}

\begin{corrige}[Exercice 7 — Texte à trous]
\eng{Amina works for a bank in Yaoundé. Her manager told her that she must \textbf{improve} her English if she wants a promotion. Her current \textbf{level} is A2, which is not enough to speak with foreign customers. She has identified her \textbf{needs}: she must work on her speaking and writing \textbf{skills}. She has decided to enrol in an evening \textbf{course} at a language centre in Bastos. After six months of \textbf{training}, she hopes to obtain an international \textbf{certificate} such as the TOEIC. Then she will be able to \textbf{apply} for the post of customer relations officer, which requires \textbf{fluent} English.}

\textbf{Mots en trop :} la boîte contenait douze mots pour dix trous ; les deux mots non utilisés sont \eng{brush up} et \eng{salary} (ni l'un ni l'autre ne convient au sens du texte : Amina suit une formation complète, elle ne « rafraîchit » pas simplement son anglais, et le texte ne parle pas de salaire).

\textbf{Justifications :} \eng{improve her English} (collocation) ; \eng{her current level} (A2 est un niveau du Cadre européen) ; \eng{identified her needs} (les besoins : speaking et writing) ; \eng{speaking and writing skills} (collocation) ; \eng{evening course} (cours du soir) ; \eng{six months of training} (nom abstrait après \eng{of}) ; \eng{obtain a certificate} (collocation : \eng{pass} s'emploie avec \eng{an examination}, pas avec \eng{a certificate}) ; \eng{apply for a post} (préposition \eng{for}) ; \eng{fluent English} (adjectif devant le nom).
\end{corrige}

\begin{corrige}[Exercice 8 — Appariement : collocations]
\begin{center}
\begin{tabular}{|l|l|}
\hline
\rowcolor{popblueL} Collocation correcte & Traduction \\
\hline
1 -- b : to brush up one's English & rafraîchir son anglais \\
\hline
2 -- a : to meet the requirements & satisfaire aux exigences \\
\hline
3 -- c : to attend a language course & suivre un cours de langue \\
\hline
4 -- f : to improve one's skills & améliorer ses compétences \\
\hline
5 -- d : to apply for a job & postuler à un emploi \\
\hline
6 -- e : to obtain a certificate & obtenir un certificat \\
\hline
7 -- h : to have a good command of English & bien maîtriser l'anglais \\
\hline
8 -- g : to get a promotion & obtenir une promotion \\
\hline
9 -- i : to sit an examination & passer un examen \\
\hline
10 -- j : to enrol in a training centre & s'inscrire dans un centre de formation \\
\hline
\end{tabular}
\end{center}
\textbf{Exemples de phrases personnelles :}
\begin{itemize}
\item \eng{I brush up my English every summer before the new school year.}
\item \eng{Last year, I attended a language course at the American Language Center in Douala.}
\item \eng{I want to improve my speaking skills because I am shy when I talk.}
\item \eng{After the Bac, I will apply for a job in a bilingual company.}
\item \eng{My brother sat the TOEFL examination and obtained a good score.}
\end{itemize}
\textbf{Point de vigilance :} \eng{sit an examination} (britannique) = \eng{take an examination} (américain) ; attention : \eng{pass an examination} signifie « réussir un examen », pas « se présenter à un examen ».
\end{corrige}

\begin{corrige}[Exercice 9 — Compréhension de texte — type Bac]
\textbf{A. Comprehension questions (réponses modèles)}
\begin{enumerate}
\item Companies invest in English training because business in Central and West Africa is increasingly conducted in English, and because major trading partners such as Nigeria, Ghana and Kenya use English for contracts, e-mails and meetings.
\item Three years ago, Njomba Logistics was losing contracts because its teams could not communicate with Nigerian clients in English.
\item He means that the company had all the material and human resources (vehicles and staff) but lacked the essential tool for international business: the English language.
\item First, the company assessed the English level of every employee. Second, it designed a training plan adapted to each level: general English for beginners and specialised business English for intermediate learners.
\item The company paid eighty per cent of the fees and the employees paid the remaining twenty per cent, so that they would feel involved and attend the classes regularly.
\item After the training, the company signed three new contracts with partners in Lagos and Accra, and the trained employees received a certificate and a salary increase.
\end{enumerate}
\textbf{B. Language work}
\begin{enumerate}
\item Noun of \eng{to employ} : \eng{employee} (« every employee »). Verb of the same family as \eng{a certificate} : \eng{to certify} (dans le texte, l'expression utilisée est \eng{received a certificate}).
\item \emph{Évaluer} : \eng{assessed} (« the company first assessed the English level »). \emph{Une augmentation de salaire} : \eng{a salary increase}. \emph{La volonté d'apprendre} : \eng{their willingness to learn}.
\item \eng{The director concludes that anticipating their language needs was the best investment they had ever made.} — Discours indirect : \eng{our} $\rightarrow$ \eng{their}, \eng{we ever made} $\rightarrow$ \eng{they had ever made} (recul du temps au pluperfect).
\end{enumerate}
\textbf{Barème indicatif (sur 20) :} A : 6 questions $\times$ 2 pts = 12 pts (1 pt pour le sens, 1 pt pour la langue) ; B : 8 pts (2 + 3 + 3). \textbf{Points de vigilance :} répondre par des phrases complètes ; ne pas recopier de longs passages ; au discours indirect, penser au recul des temps et au changement des pronoms.
\end{corrige}

\begin{corrige}[Exercice 10 — Traduction — type Bac]
\textbf{Proposition de traduction :}

\eng{Nowadays, English has become essential in the world of work. In Cameroon, many companies require their employees to speak English fluently. That is why I have decided to attend a training course at a language centre. I want to improve my level, obtain an international certificate and thus get a promotion. Anticipating one's language needs means preparing one's professional future.}

\textbf{Points de vigilance :}
\begin{itemize}
\item « De nos jours » = \eng{nowadays} (un seul mot, avec \eng{-s}).
\item « Exigent que leurs employés parlent » : construction \eng{require somebody to do something} $\rightarrow$ \eng{require their employees to speak}. Éviter le subjonctif calqué \eng{require that their employees speak}, lourd au Bac.
\item « Couramment » = \eng{fluently} ; ordre : \eng{to speak English fluently}.
\item « C'est pourquoi » = \eng{that is why} (pas \eng{it is why}).
\item « Décrocher une promotion » = \eng{get a promotion} (registre soutenu : \eng{obtain a promotion}).
\item « Anticiper ses besoins linguistiques, c'est préparer... » : tournure anglaise avec \eng{means + V-ing} : \eng{Anticipating one's language needs means preparing...}.
\end{itemize}
\textbf{Barème indicatif (sur 10) :} 5 segments $\times$ 2 pts (1 pt sens, 1 pt correction grammaticale).
\end{corrige}

\begin{corrige}[Exercice 11 — Expression écrite — type Bac]
\textbf{Topic 1 — Modèle de paragraphe (environ 90 mots) :}

\eng{Dear Sir, I am writing to ask for your support for an English training course. My current level is B1, but my position requires fluent spoken English to deal with our Nigerian partners. I would like to attend evening classes at the British Council for six months and then sit the TOEIC examination. This training will improve my communication skills and help me negotiate contracts more efficiently. Investing in my English means investing in the image of our company. I hope you will consider my request favourably.}

\textbf{Topic 2 — Modèle de composition (environ 160 mots) :}

\eng{Should every Cameroonian worker be trained in English? In my opinion, the answer is clearly yes, for three main reasons.

First, Cameroon is surrounded by English-speaking partners such as Nigeria and Ghana: trade with these countries is conducted in English. Secondly, a worker who has a good command of English has access to better jobs, promotions and international conferences. Finally, English is the language of technology and of most professional documentation; without it, workers remain dependent on translators.

A concrete example is the cocoa sector in the South-West region: exporters who trained their staff in business English now negotiate directly with buyers in London and Accra, without intermediaries, and earn more money.

To conclude, English training is not a luxury but a necessity. Anticipating one's language needs means preparing one's professional future, and companies as well as the State should invest in it.}

\textbf{Topic 3 — Grille d'évaluation du jeu de rôle :} le candidat doit utiliser au moins cinq mots du lexique de la leçon (\eng{level, fluent, certificate, training, command, skills}) ; l'examinateur pose cinq questions variées (\eng{What is your current English level? Do you have any language certificates? How do you practise your English? Why do you want to improve it? What is your training project?}).

\textbf{Barème indicatif (sur 20) :} respect de la consigne et de la longueur : 4 pts ; richesse et exactitude du lexique thématique : 6 pts ; grammaire et orthographe : 6 pts ; organisation et articulation du texte : 4 pts.

\textbf{Points de vigilance :} respecter le nombre de mots ; éviter les calques (\eng{formation}, \eng{ameliorate}, \eng{actually}) ; articuler avec \eng{first, secondly, finally, to conclude} ; relire les accords sujet--verbe et les prépositions (\eng{fluent in}, \eng{apply for}, \eng{enrol in}).
\end{corrige}

\newpage

\coursec{Fiche bilan}

\begin{popfigure}
\begin{tikzpicture}[mindmap, grow cyclic, every node/.style={concept, text=white, align=center, font=\small},
  root concept/.append style={concept color=popblue, font=\bfseries\small},
  level 1/.append style={level distance=3.2cm, sibling angle=60},
  level 2/.append style={level distance=2.1cm, sibling angle=45, font=\scriptsize}]
\node[root concept]{Second language needs at work}
  child[concept color=poporange]{ node {Assessing my level}
    child { node {self-assessment} }
    child { node {language needs} }
  }
  child[concept color=popgreen]{ node {Training and skills}
    child { node {attend a course} }
    child { node {improve fluency} }
  }
  child[concept color=popgold]{ node {English at work}
    child { node {job interview} }
    child { node {CV / résumé} }
  }
  child[concept color=poppink]{ node {False friends}
    child { node {actually $\neq$ actuellement} }
    child { node {formation $\neq$ training} }
  }
  child[concept color=poppurple]{ node {Word families}
    child { node {employ / employer} }
    child { node {apply / applicant} }
  }
  child[concept color=popblue!70]{ node {Using words in context}
    child { node {dialogues} }
    child { node {short texts} }
  };
\end{tikzpicture}
\legende{Carte mentale de la leçon : anticiper ses besoins en anglais au travail.}
\end{popfigure}

\begin{bilan}
\textbf{1. Lexique clé à connaître par cœur}
\begin{itemize}
  \item \eng{language needs} : besoins linguistiques ; \eng{to assess one's level} : évaluer son niveau ; \eng{self-assessment} : auto-évaluation.
  \item \eng{to attend a training course} : suivre une formation ; \eng{to improve one's fluency} : améliorer son aisance ; \eng{to brush up one's English} : se remettre à l'anglais.
  \item \eng{a job interview} : un entretien d'embauche ; \eng{a CV / a résumé} : un CV ; \eng{to apply for a job} : postuler à un emploi ; \eng{an applicant} : un candidat.
  \item \eng{proficiency} : maîtrise ; \eng{a native speaker} : un locuteur natif ; \eng{bilingual} : bilingue ; \eng{workplace} : lieu de travail.
  \item Collocations : \eng{to meet a requirement} (satisfaire à une exigence), \eng{to acquire skills} (acquérir des compétences), \eng{to communicate effectively} (communiquer efficacement).
\end{itemize}

\textbf{2. Faux amis et pièges des francophones}
\begin{center}
\begin{tabularx}{\linewidth}{|X|X|X|}\hline
\rowcolor{popblueL} Français & English correct & Piège à éviter \\ \hline
actuellement & currently, at present & \eng{actually} = en fait \\ \hline
une formation & training, a course & \eng{a formation} = une formation géologique \\ \hline
postuler & to apply for & jamais \eng{to postulate} \\ \hline
assister à (un cours) & to attend & \eng{to assist} = aider \\ \hline
un stage & an internship, a placement & \eng{a stage} = une étape / une scène \\ \hline
\end{tabularx}
\end{center}

\textbf{3. Familles de mots et prononciation}
\begin{itemize}
  \item \eng{to employ} $\rightarrow$ \eng{employer} (employeur), \eng{employee} (employé), \eng{employment} (emploi), \eng{unemployed} (au chômage).
  \item \eng{to apply} $\rightarrow$ \eng{application}, \eng{applicant} ; \eng{to train} $\rightarrow$ \eng{training}, \eng{trainer}, \eng{trainee} (stagiaire).
  \item Prononciation : \eng{employee} « an-plo-i-i » (accent sur la dernière syllabe) ; \eng{interview} « in-teur-viou » ; \eng{fluent} « flou-annt ».
\end{itemize}

\textbf{4. Méthodes express}
\begin{itemize}
  \item Pour parler de ses besoins : \eng{I need to improve my English because...} / \eng{My level is not good enough to...}.
  \item Pour un entretien : saluer, se présenter (\eng{I hold a degree in...}), parler de ses compétences (\eng{I am fluent in English and French}), remercier.
  \item Toujours réemployer le nouveau mot dans une phrase personnelle pour le mémoriser.
\end{itemize}
\end{bilan}

\begin{aretenir}[Checklist avant le Bac]
\begin{itemize}
  \item $\square$ Je sais définir \eng{language needs} et \eng{self-assessment} en anglais.
  \item $\square$ Je connais au moins dix mots du champ lexical de la formation et du travail.
  \item $\square$ Je distingue \eng{actually} et « actuellement », \eng{to attend} et « assister ».
  \item $\square$ Je sais conjuguer et employer \eng{to apply for a job} sans oublier la préposition \eng{for}.
  \item $\square$ Je connais la famille de \eng{employ} : employer, employee, employment, unemployed.
  \item $\square$ Je sais dire \eng{to attend a training course} et non \eng{a formation}.
  \item $\square$ Je sais présenter mon niveau d'anglais : \eng{I am fluent in...}, \eng{I have a good command of...}.
  \item $\square$ Je sais prononcer \eng{employee}, \eng{interview} et \eng{fluency} correctement.
  \item $\square$ Je peux rédiger cinq phrases personnelles réutilisant le lexique de la leçon.
  \item $\square$ Je peux jouer un court dialogue d'entretien d'embauche avec un camarade.
\end{itemize}
\end{aretenir}

\findecours

\end{document}
